\documentclass{article}

\usepackage{amsmath, amssymb}
\usepackage[margin=1in]{geometry}
\usepackage{csquotes}
\usepackage{hyperref}
\usepackage{algorithm}
\usepackage{algpseudocode}

\title{Self similar diffusion}
\author{}
\date{\today}

\begin{document}
\maketitle

\section{The self-similar framework}\label{sec:framework}


\subsection{Rescaling into a self-similar frame}\label{ssec:fw-rescaling}

Let $u_i^*(\mathbf{x}^*, t^*)$, $i = 1, \dots, d$, satisfy an evolution
equation of the form
\begin{equation}\label{eq:fw-general-pde}
  \frac{\partial u_i^*}{\partial t^*} = R_i^*(u^*),
\end{equation}
on $\mathbb{R}^d$. We rescale amplitude, length, and time by
\begin{equation}\label{eq:fw-nondim}
  u_i = \frac{u_i^*}{U}, \qquad
  x_j = \frac{x_j^*}{L}, \qquad
  t = \frac{t^*}{s(t)},
\end{equation}
where the amplitude scale $U$ and the length scale $L$ depend on the non dimensional time $t$. The timescale
is determined by specifiying the clock
\begin{equation}\label{eq:fw-clock}
  \tau = \frac{\mathrm{d}t^*}{\mathrm{d}t}.
\end{equation}
The clock $\tau$ is problem specific and can contain the time dependent length and amplitude scales.
It is defined in \S\ref{ssec:1d-clock} and \S\ref{ssec:ns-clock} for the two cases considered here.
\par
Throughout, a prime denotes differentiation with respect to $t$ and we use the following abbreviations
\begin{equation}\label{eq:fw-rates}
  \sigma_L = \frac{L^\prime}{L}, \qquad \sigma_U = \frac{U^\prime}{U}.
\end{equation}
Differentiating $x_j = x_j^*/L$ with respect to $t$ at fixed $x_j^*$ gives,
for each component $j$,
\begin{equation}\label{eq:fw-dxdt}
  \frac{\mathrm{d} x_j}{\mathrm{d} t} = -x_j^* \frac{L^\prime}{L^2} = -x_j\, \sigma_L .
\end{equation}
Applying the chain rule to the left-hand side of \eqref{eq:fw-general-pde},
and using \eqref{eq:fw-clock} and \eqref{eq:fw-dxdt},
\begin{equation}\label{eq:fw-lhs-chain}
  \begin{aligned}
  \frac{\partial u_i^*}{\partial t^*}
  &= \frac{\mathrm{d} t}{\mathrm{d} t^*}\left(
      \frac{\partial (U u_i)}{\partial t}
      + U \frac{\partial u_i}{\partial x_j} \frac{\mathrm{d} x_j}{\mathrm{d} t}
    \right) \\
  &= \frac{U}{\tau}\left(
      \frac{\partial u_i}{\partial t}
      + \sigma_U u_i
      - \sigma_L\, x_j \frac{\partial u_i}{\partial x_j}
    \right).
  \end{aligned}
\end{equation}
Equating \eqref{eq:fw-lhs-chain} to $R_i^*$ and solving for
$\partial u_i / \partial t$ gives the result used throughout this document.
Under the rescaling \eqref{eq:fw-nondim}, the equation \eqref{eq:fw-general-pde} becomes
\begin{equation}\label{eq:fw-lemma}
  \frac{\partial u_i}{\partial t}
  = \underbrace{\sigma_L\, x_j \frac{\partial u_i}{\partial x_j}
    - \sigma_U\, u_i}_{\text{scaling terms}}
  + \underbrace{\frac{\tau}{U}\, R_i^*(U u,\, L x)}_{\text{physical terms}} .
\end{equation}
\par

Accordingly we split \eqref{eq:fw-lemma} into two flows. The \emph{scaling
operator} $\Phi_{\Delta t}^{(\mathrm{S})}$ advances the scaling terms,
\begin{equation}\label{eq:fw-scaling-flow}
  \Phi_{\Delta t}^{(\mathrm{S})}(u_i)
  = u_i + \int_t^{t + \Delta t}
    \left( \sigma_L\, x_j \frac{\partial u_i}{\partial x_j} - \sigma_U\, u_i \right)
    \mathrm{d} t,
\end{equation}
and the \emph{physics operator} $\Phi_{\Delta t}^{(\mathrm{P})}$ advances the
rest,
\begin{equation}\label{eq:fw-physics-flow}
  \Phi_{\Delta t}^{(\mathrm{P})}(u_i)
  = u_i + \int_t^{t + \Delta t} \frac{\tau}{U}\, R_i^*(U u,\, L x)\, \mathrm{d} t .
\end{equation}
Both are written as the field plus the integral of the terms they advance.
$\Phi^{(\mathrm{S})}$ is obtained in closed form below and is shared;
$\Phi^{(\mathrm{P})}$ is equation-specific and is supplied in
\S\ref{ssec:1d-physics-step} and \S\ref{ssec:ns-physics-step}.

\subsection{The scaling operator}\label{ssec:fw-scaling-operator}

Evaluating $\Phi_{\Delta t}^{(\mathrm{S})}$ amounts to solving
\begin{equation}\label{eq:fw-scaling-pde}
  \frac{\partial u_i}{\partial t}
  = \sigma_L\, x_j \frac{\partial u_i}{\partial x_j} - \sigma_U\, u_i .
\end{equation}
This is a linear, first-order PDE and can be solved by the method of
characteristics. Note that the scaling terms do not couple the components of
$u_i$, so the argument below is componentwise.

Let $x_j(t)$ denote a characteristic curve satisfying
\begin{equation}\label{eq:fw-characteristic}
  \frac{\mathrm{d} x_j}{\mathrm{d} t} = -x_j\, \sigma_L.
\end{equation}
 The total derivative of $u_i$ along such a
curve is
\begin{equation}\label{eq:fw-total-derivative}
  \frac{\mathrm{d} u_i}{\mathrm{d} t}
  = \frac{\partial u_i}{\partial t} + \frac{\partial u_i}{\partial x_j} \frac{\mathrm{d} x_j}{\mathrm{d} t}
  = \frac{\partial u_i}{\partial t} - \sigma_L\, x_j \frac{\partial u_i}{\partial x_j},
\end{equation}
so along the characteristic \eqref{eq:fw-scaling-pde} reduces to the ordinary
differential equation
\begin{equation}\label{eq:fw-scaling-ode}
  \frac{\mathrm{d} u_i}{\mathrm{d} t} = -\sigma_U\, u_i .
\end{equation}

Integrating \eqref{eq:fw-characteristic} from $t_k$ to $t_{k+1}$ gives
\begin{equation}\label{eq:fw-char-integrated}
  \ln x_j(t_{k+1}) - \ln x_j(t_k) = -\bigl[\ln L_{k+1} - \ln L_k\bigr]
  \quad\Longrightarrow\quad
  x_j(t_{k+1}) = x_j(t_k)\, \frac{L_k}{L_{k+1}},
\end{equation}
so that the characteristic arriving at the evaluation point
$x_j \equiv x_j(t_{k+1})$ originated at
\begin{equation}\label{eq:fw-char-origin}
  x_j(t_k) = x_j\, \frac{L_{k+1}}{L_k} .
\end{equation}
Integrating \eqref{eq:fw-scaling-ode} gives
\begin{equation}
  u_i(t_{k+1}) = u_i(t_k)\, U_k / U_{k+1}.
\end{equation}
Combining this with the characteristic origin \eqref{eq:fw-char-origin} and the initial condition
$u_i(t_k, \cdot)$ gives the scaling operator in closed form,
\begin{equation}\label{eq:fw-scaling-closed}
  \Phi_{\Delta t}^{(\mathrm{S})}(u_i)(\mathbf{x})
  = \frac{U_k}{U_{k+1}}\, u_i\!\left( \mathbf{x}\, \frac{L_{k+1}}{L_k} \right).
\end{equation}
Writing the step ratios
\begin{equation}\label{eq:fw-step-ratios}
  \lambda_{U,k+1} = \frac{U_{k+1}}{U_k}, \qquad
  \lambda_{L,k+1} = \frac{L_{k+1}}{L_k},
\end{equation}
this is
\begin{equation}\label{eq:fw-scale-step}
  \Phi_{\Delta t}^{(\mathrm{S})}(u_i)(\mathbf{x})
  = \frac{1}{\lambda_{U,k+1}}\, u_i\!\left( \mathbf{x}\, \lambda_{L,k+1} \right).
\end{equation}

\subsubsection{The scaling operator in \texorpdfstring{$\boldsymbol{\kappa}$}{kappa}-space}\label{sssec:fw-scaling-fourier}

For equations formulated spectrally it is useful to have the transform of the
scaling terms. Using $\widehat{\cdot}$ for the spatial Fourier transform,
\begin{equation}\label{eq:fw-fourier-convention}
  \hat{u}_i(\boldsymbol{\kappa}) = \int u_i(\mathbf{x})\, e^{-\mathrm{i} \kappa_j x_j}\, \mathrm{d}\mathbf{x},
\end{equation}
and writing the spread term in conservative form,
\begin{equation}\label{eq:fw-scaling-divform}
  x_j \frac{\partial u_i}{\partial x_j}
  = \frac{\partial}{\partial x_j}\left( x_j u_i \right) - d\, u_i,
\end{equation}
gives
\begin{equation}\label{eq:fw-fourier-scaling}
  \widehat{x_j \frac{\partial u_i}{\partial x_j}}
  = \mathrm{i} \kappa_j \left( \mathrm{i} \frac{\partial \hat{u}_i}{\partial \kappa_j} \right) - d\, \hat{u}_i
  = -\left( \kappa_j \frac{\partial \hat{u}_i}{\partial \kappa_j} + d\, \hat{u}_i \right),
\end{equation}
so that the transform of \eqref{eq:fw-scaling-pde} is
\begin{equation}\label{eq:fw-scaling-pde-fourier}
  \frac{\partial \hat{u}_i}{\partial t}
  = -\sigma_L \left( \kappa_j \frac{\partial \hat{u}_i}{\partial \kappa_j} + d\, \hat{u}_i \right)
  - \sigma_U\, \hat{u}_i .
\end{equation}
This is again a linear, first-order PDE, now with characteristics in
$\boldsymbol{\kappa}$ satisfying
\begin{equation}\label{eq:fw-characteristic-fourier}
  \frac{\mathrm{d} \kappa_j}{\mathrm{d} t} = \kappa_j\, \sigma_L ,
\end{equation}
along which \eqref{eq:fw-scaling-pde-fourier} reduces to
\begin{equation}\label{eq:fw-scaling-ode-fourier}
  \frac{\mathrm{d} \hat{u}_i}{\mathrm{d} t} = -\left( d\, \sigma_L + \sigma_U \right) \hat{u}_i .
\end{equation}
Integrating \eqref{eq:fw-characteristic-fourier} from $t_k$ to $t_{k+1}$ gives
$\kappa_j(t_{k+1}) = \kappa_j(t_k)\, \lambda_{L,k+1}$, so the characteristic arriving at
$\boldsymbol{\kappa}$ originated at $\boldsymbol{\kappa} / \lambda_{L,k+1}$, while integrating
\eqref{eq:fw-scaling-ode-fourier} gives
$\hat{u}_i(t_{k+1}) = \hat{u}_i(t_k) / (\lambda_{U,k+1} \lambda_{L,k+1}^{d})$.
The scaling operator in $\boldsymbol{\kappa}$-space is therefore
\begin{equation}\label{eq:fw-scale-step-fourier}
  \widehat{\Phi_{\Delta t}^{(\mathrm{S})}(u_i)}(\boldsymbol{\kappa})
  = \frac{1}{\lambda_{U,k+1}\, \lambda_{L,k+1}^{d}}\,
    \hat{u}_i\!\left( \frac{\boldsymbol{\kappa}}{\lambda_{L,k+1}} \right).
\end{equation}
This agrees with transforming \eqref{eq:fw-scale-step} directly: the
substitution $\mathbf{x}' = \mathbf{x}\, \lambda_{L,k+1}$ in
\eqref{eq:fw-fourier-convention} contributes the Jacobian
$\lambda_{L,k+1}^{-d}$ and maps $\boldsymbol{\kappa}$ to $\boldsymbol{\kappa} / \lambda_{L,k+1}$.
A growing length scale ($\lambda_{L,k+1} > 1$) thus compresses $u_i$ toward
$\mathbf{x} = 0$ and correspondingly stretches $\hat{u}_i$ toward larger
$|\boldsymbol{\kappa}|$.

\subsection{Spectral interpolation}\label{ssec:fw-interpolation}

The closed form \eqref{eq:fw-scale-step} requires $u_i$ at the rescaled points
$\mathbf{x}\,\lambda_{L,k+1}$, which are not grid points. On a uniform
periodic grid the represented field is band-limited, so it can be evaluated at
arbitrary points exactly, with no interpolation error, by summing its
truncated Fourier series there:
\begin{equation}\label{eq:fw-spectral-interp}
  u_i\!\left(\mathbf{x}\, \lambda_{L,k+1}\right)
  = \sum_{\boldsymbol{\kappa} \in \mathcal{K}} \hat{u}_i(\boldsymbol{\kappa})\,
    \exp\!\left( \mathrm{i}\, \kappa_j x_j\, \lambda_{L,k+1} \right),
\end{equation}
where $\mathcal{K}$ is the set of represented wavenumbers and
$\hat{u}_i(\boldsymbol{\kappa})$ are the normalized transform coefficients of $u_i$.
This is the defining property of the scheme: the scaling step introduces no
spatial discretization error of its own.

\paragraph{Queries outside the box.}
The true problem we are trying to approximate is posed on an infinite box, with $u_i$ going to zero as it approaches the boundaries.
For $\lambda_{L,k+1} > 1$ it is possible that $|x_j\,\lambda_{L,k+1}| > H/2$, in which case we would like the rescaling in \eqref{eq:fw-spectral-interp} to return $0$.
However, this is not the current behavior: instead, it wraps around the periodic box and returns a potentially nonzero value from the other end of the domain.
Therefore, to obtain the desired behavior we modify the scaling step to
\begin{equation}\label{eq:fw-dilate-cutoff}
  u_i\!\left(\mathbf{x}\, \lambda_{L,k+1}\right)
  = \chi\!\left(\mathbf{x}\, \lambda_{L,k+1}\right)
    \sum_{\boldsymbol{\kappa} \in \mathcal{K}} \hat{u}_i(\boldsymbol{\kappa})\,
    \exp\!\left( \mathrm{i}\, \kappa_j x_j\, \lambda_{L,k+1} \right),
  \qquad
  \chi(\mathbf{y}) = \prod_{d} \mathbb{1}\!\left[\, |y_d| \le H/2 \,\right].
\end{equation}

\subsection{Constraint functionals}\label{ssec:fw-constraints}
We will set the ratios $\lambda_{U,k+1}, \lambda_{L,k+1}$ that the
scaling step \eqref{eq:fw-scale-step} depends on by requiring
the rescaled field to carry prescribed values of two functionals $F_1, F_2$ of the form
\begin{equation}\label{eq:fw-homogeneity}
  F_m\!\left[ \frac{1}{\lambda_U}\, u\!\left( x \lambda_L \right) \right]
  = \lambda_U^{-p_m}\, \lambda_L^{-q_m}\, F_m[u],
  \qquad m = 1, 2 .
\end{equation}
The targets $F_1, F_2$ are fixed once, from the initial condition
(\S\ref{ssec:alg-init}). Given a field $u$ with current values
$F_m^{\mathrm{in}} = F_m[u]$, imposing
$\lambda_U^{-p_m} \lambda_L^{-q_m} F_m^{\mathrm{in}} = F_m$ and taking
logarithms gives the linear system
\begin{equation}\label{eq:fw-constraint-system}
  \begin{bmatrix} p_1 & q_1 \\ p_2 & q_2 \end{bmatrix}
  \begin{bmatrix} \ln \lambda_U \\ \ln \lambda_L \end{bmatrix}
  = \begin{bmatrix} \ln F_1^{\mathrm{in}} - \ln F_1 \\ \ln F_2^{\mathrm{in}} - \ln F_2 \end{bmatrix},
\end{equation}
which is uniquely solvable when $p_1 q_2 - p_2 q_1 \neq 0$, i.e.\ when each
functional responds to amplitude and to length rescaling differently from the
other, and is solved in closed form. Once $F_m^{\mathrm{in}}$ is evaluated, no
further evaluation of the scaled field is needed to find the ratios.
\par
In both cases the second functional is a second moment, and we normalize it
by the target value of the zeroth moment $F_1$. As the normalization is a
constant, it changes neither the exponents nor the ratios the constraints
produce; it gives $F_2$ the units of a squared length, a mean-square width.
The functionals themselves are equation-specific
(\S\ref{ssec:1d-moments}, \S\ref{ssec:ns-constraints}); the substep that
applies \eqref{eq:fw-constraint-system} is \textsc{ConstrainedScale}
(\S\ref{ssec:alg-substeps}).

\subsection{Power-law ansatz and the step in physical time}\label{ssec:fw-powerlaw}

The algorithm also needs to \emph{predict} scale ratios over a fraction of a
step, at points where no constraint is available to fix them
(\S\ref{ssec:alg-strang}). For this we assume that both scales follow a power
law in physical time,
\begin{equation}\label{eq:fw-powerlaw}
  X = A_X\, (t^* + \alpha_X)^{\beta_X}, \qquad X \in \{U, L\},
\end{equation}
with a prefactor $A_X > 0$, an offset $\alpha_X$, and an exponent $\beta_X$.
The prefactor cannot be dropped: the scales are normalized to $U_0 = L_0 = 1$
at the initial condition (\S\ref{ssec:alg-init}), before the flow has reached
the self-similar state, so nothing ties that normalization to the one of the
power law. The prefactor cancels from the ratio over a fraction $i$ of step
$k$,
\begin{equation}\label{eq:fw-powerlaw-ratio}
  \frac{X_{k+i}}{X_k} = \left( 1 + \frac{i\, \Delta t_k^*}{t_k^* + \alpha_X} \right)^{\beta_X}.
\end{equation}
Since the time step $\Delta t_k = t_{k+1} - t_k$ is taken in the
self-similar time $t$, the corresponding physical increment must be obtained
from the clock,
\begin{equation}
  \Delta t_k^* = \int_{t_k}^{t_{k+1}} \tau\, \mathrm{d} t,
\end{equation}
which can be approximated with a three-point Adams--Bashforth-style quadrature:
\begin{equation}\label{eq:fw-dtstar}
  \Delta t_k^* = \int_{t_k}^{t_{k+1}} \tau\, \mathrm{d} t
  \approx w_0\, \tau_k + w_1\, \tau_{k-1} + w_2\, \tau_{k-2}.
\end{equation}
The coefficients can be computed as follows.
\begin{equation}\label{eq:fw-dtstar-weights}
  \begin{aligned}
    w_0 &= \frac{h^3/3 + (2h_1 + h_2)\, h^2/2 + h_1 (h_1 + h_2)\, h}{h_1 (h_1 + h_2)}, \\
    w_1 &= -\frac{h^3/3 + (h_1 + h_2)\, h^2/2}{h_1 h_2}, \\
    w_2 &= \frac{h^3/3 + h_1\, h^2/2}{h_2 (h_1 + h_2)},
  \end{aligned}
\end{equation}
where $h = \Delta t_k$, $h_1 = \Delta t_{k-1}$, $h_2 = \Delta t_{k-2}$.
In this work the step is constant, $\Delta t_k = \Delta t$ for every $k$
(\S\ref{ssec:alg-timestep}). With $h = h_1 = h_2 = \Delta t$ the weights
reduce to those of the third-order Adams--Bashforth method,
\begin{equation}\label{eq:fw-dtstar-weights-const}
  w_0 = \frac{23}{12}\, \Delta t, \qquad
  w_1 = -\frac{16}{12}\, \Delta t, \qquad
  w_2 = \frac{5}{12}\, \Delta t .
\end{equation}
\par
Eliminating $t_k^* + \alpha_X = (X_k/A_X)^{1/\beta_X}$ from
\eqref{eq:fw-powerlaw-ratio} and taking logarithms gives
\begin{equation}\label{eq:fw-log-ratio-A}
  \ln \frac{X_{k+i}}{X_k}
  = \beta_X \ln\!\left( 1 + \exp\!\left( \ln i + \ln \Delta t_k^* - \frac{\ln X_k - \ln A_X}{\beta_X} \right) \right).
\end{equation}
The algorithm uses this with $\ln A_X = 0$, as the function
\begin{equation}\label{eq:fw-log-ratio}
  \ln \frac{X_{k+i}}{X_k} = G_X(i, \beta_X; k),
  \qquad
  G_X(i, \beta; k) = \beta \ln\!\left( 1 + \exp\!\left( \ln i + \ln \Delta t_k^* - \frac{\ln X_k}{\beta} \right) \right),
\end{equation}
in both directions: it is evaluated forward, with $i = \tfrac12$, to predict
a half-step ratio (\S\ref{ssec:alg-substeps}), and inverted for $\beta$, with
$i = 1$, to estimate the exponents from the realized history
(\S\ref{ssec:alg-beta-estimation}).

\paragraph{The prefactor within a step.}
Setting $\ln A_X = 0$ is not an assumption that $A_X = 1$; the step does
not need $A_X$. By \eqref{eq:fw-powerlaw-ratio} the ratio depends on the
pair $(\beta_X,\, t_k^* + \alpha_X)$, and expanding in the step,
\begin{equation}\label{eq:fw-ratio-expansion}
  \ln \frac{X_{k+i}}{X_k}
  = \frac{\beta_X}{t_k^* + \alpha_X}\, i\, \Delta t_k^*
    \left( 1 - \frac{i\, \Delta t_k^*}{2\,(t_k^* + \alpha_X)} + \dots \right).
\end{equation}
The leading term depends on the pair only through
$\beta_X / (t_k^* + \alpha_X)$, which is the local rate $\sigma_X/\tau$
(\S\ref{ssec:pp-rates}). The inversion of \S\ref{ssec:alg-beta-estimation}
fixes exactly this combination, by matching the realized ratio of the step
just taken, whatever $\ln A_X$ is assumed. With $\ln A_X = 0$ the $\beta$
that does so is an \emph{effective} exponent $\tilde\beta_X$ that absorbs
the prefactor: matching $(\ln X_k - \ln A_X)/\beta_X$ with
$\ln X_k / \tilde\beta_X$ gives
$\tilde\beta_X \approx \beta_X\, (1 + \ln A_X / \ln X_k)$, which tends to
$\beta_X$ as the scales grow. The prefactor only decides how the fixed
$\beta_X/(t_k^*+\alpha_X)$ is split between $\beta_X$ and $t_k^* + \alpha_X$,
which enters the correction in \eqref{eq:fw-ratio-expansion}; that is
smaller than the leading term by
$i\,\Delta t_k^*/(t_k^*+\alpha_X) = i\, \sigma_X \Delta t_k / \beta_X$,
the order of the splitting error itself. The algorithm therefore carries only
the effective exponents, and $A_X$, $\alpha_X$, $\beta_X$ are recovered
separately, from the whole history, in \S\ref{sec:postproc}.

\paragraph{Logarithmic arithmetic.}
$U_k$ and $L_k$ may become extremely large or small, so we store
$\ln U_k$ and $\ln L_k$ rather than $U_k$ and $L_k$, and perform all
arithmetic on them in the logarithmic number system, where multiplication and
division become addition and subtraction, and where sums and differences are
evaluated by
\begin{equation}\label{eq:fw-log-sum-exp}
  \log_b(X+Y) = \log_b(X) + \log_b\!\bigl(1+b^{\,y-x}\bigr),
  \qquad
  \log_b(X-Y) = \log_b(X) + \log_b\!\bigl(1-b^{\,y-x}\bigr) \quad (X>Y>0),
\end{equation}
with $x = \log_b(X)$, $y = \log_b(Y)$. In particular
\eqref{eq:fw-dtstar} is evaluated as
\begin{equation}\label{eq:fw-ln-dtstar}
  \ln \Delta t_k^* =
  \ln\!\left( w_0\, e^{\ln \tau_k} + w_1\, e^{\ln \tau_{k-1}} + w_2\, e^{\ln \tau_{k-2}} \right).
\end{equation}

\section{Algorithm}\label{sec:algorithm}

The integrator below is stated once, for any equation that supplies the three
items of \S\ref{sec:framework}: a clock $\tau$, a physics operator
$\Phi^{(\mathrm{P})}$, and a pair of constraint functionals. Everything else
--- the splitting, the scale bookkeeping, the exponent estimation --- is
common.

\subsection{Domain, grid, and notation}\label{ssec:alg-domain}

The spatial domain is periodic with a fixed physical side length $H$ and $N$
grid points per direction, centred at the origin, so that
$x_j \in [-H/2, H/2)$ with uniform spacing $\Delta x = H/N$ and grid points
\begin{equation}\label{eq:alg-grid}
  x_n = -\frac{H}{2} + n\, \Delta x, \qquad n = 0, 1, \dots, N-1
\end{equation}
in each of the $d$ directions. Every integral over the domain below is
approximated by the Riemann sum
\begin{equation}\label{eq:alg-riemann}
  \int f(\mathbf{x})\,\mathrm{d}\mathbf{x} \approx (\Delta x)^d \sum_{\mathbf{n}} f(\mathbf{x_n})
\end{equation}
on this grid.

An index range abbreviates the corresponding tuple of stored values, e.g.\
$\ln \tau_{k-2:k} = (\ln \tau_{k-2}, \ln \tau_{k-1}, \ln \tau_k)$. The run
constants $H, N$, the time step $\Delta t$ of \S\ref{ssec:alg-timestep}, the
constraint targets $F_1, F_2$, and whatever physical parameters the equation
carries are visible to every subroutine and are omitted from their signatures.
The step is constant, $\Delta t_k = \Delta t$; the subscript is kept where a
formula holds for a general step. As set out in \S\ref{ssec:fw-powerlaw}, all
arithmetic on the scales is done on $\ln U_k, \ln L_k$.

\subsection{Initialization}\label{ssec:alg-init}

A run is specified by the initial field $u_0$, the domain parameters $H, N$,
the time step $\Delta t$ of
\S\ref{ssec:alg-timestep}, a time horizon $T_{\mathrm{end}}$, the
splitting scheme
$\mathit{scheme} \in \{\textsc{LieTrotter}, \textsc{Strang}\}$
(\S\ref{ssec:alg-lie-trotter}, \S\ref{ssec:alg-strang}), for equations that
define one, the switch $\mathit{project}$ of the optional projection inside
\textsc{Scale} (\S\ref{ssec:alg-substeps}), for equations that use it (here
only Navier--Stokes), the gain $\gamma > 0$ and initial clock rate
$c_{\mathrm{init}}$ of the clock control (\S\ref{ssec:ns-clock-control};
$\gamma = \varnothing$ for none), and initial guesses
$\beta_U^{(0)}, \beta_L^{(0)}$, which only seed the bracket of the first
root-find for $\beta_U, \beta_L$ (\S\ref{ssec:alg-beta-estimation}). From these:
\begin{itemize}
  \item The integration loop starts at $t_0 = 0$ and runs while $t_k < T_{\mathrm{end}}$.
  \item The self-similar frame is initialized at the initial condition itself,
  $L_0 = U_0 = 1$, i.e.\ $\ln L_0 = \ln U_0 = 0$.
  \item The \emph{target} values $F_1, F_2$ of the constraint functionals
  (\S\ref{ssec:fw-constraints}) are fixed once, from $u_0$ itself,
  \begin{equation}\label{eq:alg-targets}
    F_1 = F_1[u_0], \qquad F_2 = F_2[u_0].
  \end{equation}
  \item The flag $\beta_{\mathrm{ok}}$, which records whether the current
  $\beta_U, \beta_L$ are usable by \textsc{Strang}, starts as
  $\beta_{\mathrm{ok}} = \text{false}$.
  \item Under clock control the clock rate starts at $c_{-1} = c_{\mathrm{init}}$,
  the rate the first step relaxes from (\S\ref{ssec:ns-clock-control}), and the miss
  \eqref{eq:ns-cc-miss-def} of $u_0$ is $\mu_0 = 0$ by
  \eqref{eq:alg-targets}.
\end{itemize}

\subsection{Time step}\label{ssec:alg-timestep}

Every step is taken at the same step $\Delta t_k = \Delta t$, a run
parameter, so that $t_k = k\, \Delta t$ and the quadrature
\eqref{eq:fw-dtstar} uses the constant weights
\eqref{eq:fw-dtstar-weights-const}. A constant step also keeps the samples of
every recorded quantity equally spaced in $t$, which simplifies the time
averages of Appendix~\ref{app:stat-error}.
\par
Nothing adapts the step to the field. When the physics operator has a
state-dependent stability limit, $\Delta t$ must be chosen small enough for
the field the run settles into, and whether it is is monitored rather than
enforced. Let $\Delta t_{\mathrm{CFL}}(u)$ be the largest step at which
\textsc{Evolve} from $u$ stays within the stability limit of its time
integrator; for Navier--Stokes it is the advective bound
\eqref{eq:ns-cfl-dt}. The \emph{realized} CFL number of step $k$,
\begin{equation}\label{eq:alg-realized-cfl}
  \mathrm{CFL}_k = \frac{\Delta t}{\Delta t_{\mathrm{CFL}}(u_{k+1})},
\end{equation}
measures the step against the field it produced. It is recorded every step,
and the run warns the first time it exceeds $1$; the remedy is a smaller
$\Delta t$. The 1D diffusion equation has no such limit, since its
propagator \eqref{eq:1d-propagator} is exact at any step.

\paragraph{The end of the run.} The loop runs while $t_k < T_{\mathrm{end}}$,
and the last step is \emph{not} shortened to land on $T_{\mathrm{end}}$, which
would break the constant step. The run overshoots $T_{\mathrm{end}}$ by less
than one step, which is immaterial on the logarithmic clock $t$. For the same
reason, output taken at prescribed times $t$ is taken at the first step
reaching each time, never by shortening a step.

\subsection{Top-level integrator}\label{ssec:alg-top-level}
\begin{algorithm}[H]
\caption{Top-level integrator}
\begin{algorithmic}[1]
\Require $u_0$, $H$, $N$, $\Delta t$, $T_{\mathrm{end}}$, $\mathit{scheme}$, $\gamma$, $c_{\mathrm{init}}$, $\beta_U^{(0)}$, $\beta_L^{(0)}$
\State $t_0 \gets 0$, $k \gets 0$
\State $\ln U_0 \gets 0$, $\ln L_0 \gets 0$
\State Compute constraint targets $F_1, F_2$ from $u_0$ via \eqref{eq:alg-targets}
\State $\beta_U \gets \beta_U^{(0)}$, $\beta_L \gets \beta_L^{(0)}$, $\beta_{\mathrm{ok}} \gets \text{false}$
\State $c_{-1} \gets c_{\mathrm{init}}$, $\mu_0 \gets 0$ \Comment{under clock control, \S\ref{ssec:ns-clock-control}}
\State Store $(u_0, \ln U_0, \ln L_0, \ln \tau_0, \beta_U, \beta_L, \beta_{\mathrm{ok}}, \mu_0)$
\While{$t_k < T_{\mathrm{end}}$}
  \If{$\gamma \neq \varnothing$}
    \State $C \gets (c_{k-1}, \mu_k)$ \Comment{\S\ref{ssec:ns-clock-control}}
  \Else
    \State $C \gets \varnothing$
  \EndIf
  \If{$\mathit{scheme} = \textsc{Strang}$ \textbf{and} $\beta_{\mathrm{ok}}$}
    \State $(u_{k+1}, \ln U_{k+1}, \ln L_{k+1}, c_k) \gets \Call{Strang}{u_k, \ln U_k, \ln L_k, \ln \tau_{k-2:k}, \Delta t, \beta_U, \beta_L, C}$ \Comment{\S\ref{ssec:alg-strang}}
  \Else
    \State $(u_{k+1}, \ln U_{k+1}, \ln L_{k+1}, c_k) \gets \Call{LieTrotter}{u_k, \ln U_k, \ln L_k, \Delta t, C}$ \Comment{\S\ref{ssec:alg-lie-trotter}}
  \EndIf
  \State $u_{k+1} \gets \Call{ApplyBC}{u_{k+1}}$ \Comment{equation-specific, \S\ref{ssec:alg-substeps}}
  \State $\ln \tau_{k+1} \gets \ln \tau(\ln U_{k+1}, \ln L_{k+1})$ \Comment{\eqref{eq:1d-ln-clock}, \eqref{eq:ns-ln-clock}}
  \State $t_{k+1} \gets t_k + \Delta t$
  \If{$k+1 \geq 3$}
    \State $(\beta_U, \beta_L, \beta_{\mathrm{ok}}) \gets \Call{EstimateBeta}{\ln U_{k:k+1}, \ln L_{k:k+1}, \ln \tau_{k-2:k}, \Delta t, \beta_U, \beta_L}$ \Comment{\S\ref{ssec:alg-beta-estimation}}
  \EndIf
  \State $\mu_{k+1} \gets \mu[u_{k+1}]$ \Comment{\eqref{eq:ns-cc-miss-def}}
  \State Store $(u_{k+1}, \ln U_{k+1}, \ln L_{k+1}, \ln \tau_{k+1}, t_{k+1}, \mathrm{CFL}_k, \beta_U, \beta_L, \beta_{\mathrm{ok}}, c_k, \mu_{k+1})$ \Comment{\eqref{eq:alg-realized-cfl}}
  \State $k \gets k + 1$
\EndWhile
\end{algorithmic}
\end{algorithm}

The run parameter $\mathit{scheme}$ selects the step function. \textsc{Strang}
needs usable estimates of $\beta_U, \beta_L$, so the loop falls back to
\textsc{LieTrotter} whenever $\beta_{\mathrm{ok}}$ is false. This is always
the case for the first three steps $k = 0, 1, 2$: the first $\beta$ update
needs four stored states, so it first runs at the end of step $k = 2$. It is
also the case for every step that follows a failed $\beta$ update
(\S\ref{ssec:alg-beta-estimation}), until a later update succeeds. The $\beta$
update itself runs regardless of $\mathit{scheme}$; under
$\mathit{scheme} = \textsc{LieTrotter}$ the estimates are diagnostic only.
Clock control needs no $\beta$ and so acts from the first step $k = 0$.
Without it ($C = \varnothing$) the returned $c_k$ is unused.

\subsection{Substeps}\label{ssec:alg-substeps}
Both splitting schemes are assembled from the substeps below. None of them
applies the boundary treatment, which is handled only by the top-level loop.

\paragraph{\textsc{Evolve}$(u, \Delta t_k)$.}
Applies the physics operator $\Phi_{\Delta t_k}^{(\mathrm{P})}$ of
\eqref{eq:fw-physics-flow}: it advances the non-scaling terms of
\eqref{eq:fw-lemma} over $\Delta t_k$, leaving the scales $\ln U, \ln L$
untouched.

\paragraph{\textsc{Scale}$(u, \ln \lambda_U, \ln \lambda_L)$.}
Applies the scaling operator \eqref{eq:fw-scale-step} with prescribed ratios
$\lambda_U, \lambda_L$, returning $\lambda_U^{-1} u(\mathbf{x_n} \lambda_L)$ at each grid
point, evaluated with the spectral interpolation of
\S\ref{ssec:fw-interpolation} and set to zero wherever $\mathbf{x_n}\lambda_L$
leaves the box \eqref{eq:fw-dilate-cutoff}. Here, and wherever the substeps are used,
$\lambda_U$ and $\lambda_L$ denote the ratio of $U$, respectively $L$, after the
substep to its value before it. Inside a Lie--Trotter step this is
$\lambda_{U,k+1} = U_{k+1}/U_k$ of \eqref{eq:fw-step-ratios}; inside a Strang step it
is a ratio across a half step, $U_{k+1/2}/U_k$ or $U_{k+1}/U_{k+1/2}$
(\S\ref{ssec:alg-strang}).

\textbf{Optional projection.} With the run parameter $\mathit{project}$
enabled, \textsc{Scale} ends by applying an equation-specific projection
$\mathcal{P}$ onto the admissible fields, returning
$\mathcal{P}\bigl(\lambda_U^{-1} u(\mathbf{x_n} \lambda_L)\bigr)$; with it disabled, or
for an equation that defines no $\mathcal{P}$, the projection is the identity. The
continuous operator \eqref{eq:fw-scale-step} maps a divergence-free field to a
divergence-free field, but its evaluation on the grid by interpolation and
the out-of-box cutoff does not preserve this exactly. For Navier--Stokes,
$\mathcal{P}$ is the Leray projection \eqref{eq:ns-leray-scale}, which restores
$\nabla \cdot \tilde{\mathbf{u}} = 0$ (\S\ref{ssec:ns-interpolation}); the 1D
diffusion equation defines none. Because every scaling step runs through
\textsc{Scale}, the projection follows the scaling step of \textsc{LieTrotter}
and \emph{both} scaling half steps of \textsc{Strang}. Inside
\textsc{ConstrainedScale} it is applied after the ratios have been solved, so
the returned field meets the targets $F_1, F_2$ only up to the part the
projection removes. The miss does not accumulate: the next
\textsc{ConstrainedScale} measures the functionals of the projected field
afresh and absorbs it into the next ratios.

\paragraph{\textsc{ConstrainedScale}$(u, \ln U, \ln L, C, \Delta^{\mathrm{pred}}_{\ln\tau,k} = 0)$.}
Rescales $u$ so that the result carries the fixed targets $F_1, F_2$ of
\eqref{eq:alg-targets} --- under the clock control of Navier--Stokes, $F_1$
and a prescribed clock increment instead --- and advances the log-scale state
accordingly:
\begin{enumerate}
  \item \textbf{Functionals of the input,}
  $F_m^{\mathrm{in}} = F_m[u]$, evaluated on the grid via
  \eqref{eq:alg-riemann}.
  \item \textbf{Ratios} $\ln \lambda_U, \ln \lambda_L$ from the linear system
  \eqref{eq:fw-constraint-system}, in closed form. This works on the scalars
  $F_m^{\mathrm{in}}$ alone and does not touch the field.
  \item \textbf{Clock (Navier--Stokes only).} If a clock state
  $C = (c_{k-1}, \mu_k)$ is given, advance the clock rate to $c_k$ by
  \eqref{eq:ns-cc-update} and replace the ratios by the solution of
  \eqref{eq:ns-cc-system} (\S\ref{ssec:ns-clock-control}), using the
  increment $\Delta^{\mathrm{pred}}_{\ln\tau,k}$ already applied by a predicted
  half step; with $C = \varnothing$ they are kept.
  \item \textbf{Rescale,} $u' = \textsc{Scale}(u, \ln \lambda_U, \ln \lambda_L)$.
  \item \textbf{Return} $u'$ together with the advanced state
  $\ln U + \ln \lambda_U$, $\ln L + \ln \lambda_L$, and $c_k$.
\end{enumerate}
Without clock control no time increment enters \textsc{ConstrainedScale}: the
ratios are fixed by the functionals of the field it is given, not by the
elapsed time. The same substep therefore serves as the full scaling step of
Lie--Trotter and as the closing scaling half step of Strang; under clock
control the two differ only in the part $\Delta^{\mathrm{pred}}_{\ln\tau,k}$ of
the step's clock increment already applied.

\paragraph{\textsc{PredictScale}$(i, \ln U_k, \ln L_k, \ln \tau_{k-2:k}, \Delta t, \beta_U, \beta_L)$.}
Predicts the ratios over the fraction $i$ of step $k$ from the power-law form
\eqref{eq:fw-powerlaw} rather than from the constraints. It evaluates
$\ln \Delta t_k^*$ from \eqref{eq:fw-ln-dtstar}, with the constant-step
weights \eqref{eq:fw-dtstar-weights-const}, and returns
$\bigl(G_U(i, \beta_U; k),\, G_L(i, \beta_L; k)\bigr)$ using \eqref{eq:fw-log-ratio}.
The same function $G_X$ is inverted for $\beta$ in
\S\ref{ssec:alg-beta-estimation}. Both use $\ln A_X = 0$, which the step
does not need (\S\ref{ssec:fw-powerlaw}).

\paragraph{\textsc{ApplyBC}$(u)$.}
Imposes whatever boundary treatment the equation requires, after the full step
has been assembled. The 1D diffusion equation clamps its two edge grid
points to zero (\S\ref{ssec:1d-boundary}). For the Navier--Stokes equations
\textsc{ApplyBC} is the identity: the far-field condition is left to the
sponge forcing of \S\ref{ssec:ns-sponge}, which acts inside \textsc{Evolve}.
\textsc{ApplyBC} performs no projection; that is done only inside
\textsc{Scale}.

\subsection{\textsc{LieTrotter}: Lie--Trotter splitting}\label{ssec:alg-lie-trotter}
A Lie--Trotter step composes a full physics step with a full scaling step,
$u_{k+1} = \Phi_{\Delta t_k}^{(\mathrm{S})}\bigl(\Phi_{\Delta t_k}^{(\mathrm{P})}(u_k)\bigr)$,
with the scaling ratios fixed by the constraints, or under clock control
by $F_1$ and the clock (\S\ref{ssec:ns-clock-control}). It uses no estimate of
$\beta_U, \beta_L$, so it is valid from the first step. It is written for a
general input state $(u, \ln U, \ln L)$ and step $\Delta t_k$ rather than only a stored state
$(u_k, \ln U_k, \ln L_k)$, so that \textsc{Strang} can call it on a
half-stepped state.
\begin{algorithm}[H]
\caption{Lie--Trotter step}
\begin{algorithmic}[1]
\Function{LieTrotter}{$u, \ln U, \ln L, \Delta t_k, C, \Delta^{\mathrm{pred}}_{\ln\tau,k} = 0$}
  \State $u^{(\mathrm{P})} \gets \Call{Evolve}{u, \Delta t_k}$ \Comment{\S\ref{ssec:alg-substeps}}
  \State \Return $\Call{ConstrainedScale}{u^{(\mathrm{P})}, \ln U, \ln L, C, \Delta^{\mathrm{pred}}_{\ln\tau,k}}$ \Comment{\S\ref{ssec:alg-substeps}}
\EndFunction
\end{algorithmic}
\end{algorithm}

\subsection{\textsc{Strang}: Strang splitting}\label{ssec:alg-strang}
A Strang step is the symmetric composition
\begin{equation}\label{eq:alg-strang-composition}
  u_{k+1} = \Phi_{\Delta t_k/2}^{(\mathrm{S})}\Bigl(\Phi_{\Delta t_k}^{(\mathrm{P})}\bigl(\Phi_{\Delta t_k/2}^{(\mathrm{S})}(u_k)\bigr)\Bigr)
\end{equation}
of a scaling half step, a full physics step, and a second scaling half step.
The last two of these are exactly a Lie--Trotter step --- a full physics step
followed by a scaling step whose ratios come from the constraints --- so
\textsc{Strang} is a predicted scaling half step followed by a call to
\textsc{LieTrotter}.
\begin{algorithm}[H]
\caption{Strang step}
\begin{algorithmic}[1]
\Function{Strang}{$u_k, \ln U_k, \ln L_k, \ln \tau_{k-2:k}, \Delta t, \beta_U, \beta_L, C$}
  \State $\bigl(\ln (U_{k+1/2}/U_k),\, \ln (L_{k+1/2}/L_k)\bigr) \gets \Call{PredictScale}{\tfrac12, \ln U_k, \ln L_k, \ln \tau_{k-2:k}, \Delta t, \beta_U, \beta_L}$ \Comment{\eqref{eq:alg-strang-predict}}
  \State $u_{k+1/2} \gets \Call{Scale}{u_k, \ln (U_{k+1/2}/U_k), \ln (L_{k+1/2}/L_k)}$
  \State $\ln U_{k+1/2} \gets \ln U_k + \ln (U_{k+1/2}/U_k)$;\ \ $\ln L_{k+1/2} \gets \ln L_k + \ln (L_{k+1/2}/L_k)$
  \State $\Delta^{\mathrm{pred}}_{\ln\tau,k} \gets \ln \tau(\ln U_{k+1/2}, \ln L_{k+1/2}) - \ln \tau_k$ \Comment{used only if $C \neq \varnothing$, \S\ref{ssec:ns-clock-control}}
  \State \Return $\Call{LieTrotter}{u_{k+1/2}, \ln U_{k+1/2}, \ln L_{k+1/2}, \Delta t_k, C, \Delta^{\mathrm{pred}}_{\ln\tau,k}}$ \Comment{\S\ref{ssec:alg-lie-trotter}}
\EndFunction
\end{algorithmic}
\end{algorithm}
No constraint fixes the state at $t_k + \Delta t_k/2$, so the first scaling half
step takes its ratios from the current power-law estimates, via
\textsc{PredictScale} with $i = 1/2$:
\begin{equation}\label{eq:alg-strang-predict}
  \ln\frac{U_{k+1/2}}{U_k} = G_U\!\left(\tfrac12, \beta_U; k\right), \qquad
  \ln\frac{L_{k+1/2}}{L_k} = G_L\!\left(\tfrac12, \beta_L; k\right).
\end{equation}
The call to \textsc{LieTrotter} then advances $u_{k+1/2}$ over the full
$\Delta t_k$, and its \textsc{ConstrainedScale} substep --- the second scaling
half step --- sets $U_{k+1}/U_{k+1/2}$ and $L_{k+1}/L_{k+1/2}$ so that
$u_{k+1}$ satisfies the constraints. Under clock control the predicted half
step is unchanged and the clock constrains the \emph{full} step: the closing
half step advances $\ln \tau$ by $c_k \Delta t$ less the increment
$\Delta^{\mathrm{pred}}_{\ln\tau,k}$ the predicted half step already applied
(\S\ref{ssec:ns-clock-control}).
The half-step state
$(u_{k+1/2}, \ln U_{k+1/2}, \ln L_{k+1/2})$ is internal to the step and is not
stored. The top-level loop calls \textsc{Strang} only when
$\beta_{\mathrm{ok}}$ is true (\S\ref{ssec:alg-top-level}), which guarantees
that $\beta_U, \beta_L$ were fitted to the stored history at the end of the
previous step and that $\ln \tau_{k-2}$ exists.

\subsection{\textsc{EstimateBeta}: estimating \texorpdfstring{$\beta_U$}{beta\_U} and \texorpdfstring{$\beta_L$}{beta\_L}}\label{ssec:alg-beta-estimation}
After every step with $k+1 \geq 3$, regardless of $\mathit{scheme}$, $\beta_U$
and $\beta_L$ are re-estimated from the log-scale history. Four stored states
are needed: $\ln X_k$ and $\ln X_{k+1}$ for the ratio, and $\ln \tau_{k-2:k}$
for $\Delta t_k^*$. The update inverts the
function $G_X$ of \eqref{eq:fw-log-ratio} with $i = 1$, the same function that
\textsc{PredictScale} evaluates forward with $i = 1/2$:
\begin{equation}\label{eq:alg-beta-rootfind}
  G_U(1, \beta_U; k) = \ln U_{k+1} - \ln U_k,
\end{equation}
and identically with $U \to L$ for $\beta_L$. The $\Delta t_k^*$ inverted
here is the one \textsc{PredictScale} used for that step's half step, so
prediction and estimation are exact inverses on the same step. Since $G_X$
carries $\ln A_X = 0$, the roots are the effective exponents
$\tilde\beta_X$ of \S\ref{ssec:fw-powerlaw}, which absorb the prefactor and
tend to $\beta_X$ as the scales grow; we write them $\beta_X$ in the
algorithm.

The sensitivity of $\beta_X$ to an error in $\Delta t_k^*$ is weak. Taking
logarithms of \eqref{eq:fw-log-ratio} for a small step,
$\ln(\ln X_{k+1} - \ln X_k) \approx \ln \beta_X + \ln \Delta t_k^* - \ln X_k/\beta_X$,
so a relative error $\eta$ in $\Delta t_k^*$ moves $\beta_X$ by
$\delta\beta_X \approx \beta_X^2 \eta / \ln X_k$, damped as the scales grow.

Each of
\eqref{eq:alg-beta-rootfind} and its $L$-counterpart is a single scalar
equation in a single unknown ($\beta_U$, resp.\ $\beta_L$), solved by
\texttt{brentq}; the implementer will need to choose a bracketing interval
(e.g.\ centred on the previous $\beta_U, \beta_L$ estimate --- on
$\beta_U^{(0)}, \beta_L^{(0)}$ for the first update --- widened adaptively if
it fails to bracket a root).

If both root-finds succeed, \textsc{EstimateBeta} returns the new
$\beta_U, \beta_L$ with $\beta_{\mathrm{ok}} = \text{true}$. If either fails
to bracket a root even after widening, the failed estimate keeps its previous
value (which seeds the next bracket), and \textsc{EstimateBeta} returns
$\beta_{\mathrm{ok}} = \text{false}$, so the next step falls back to
\textsc{LieTrotter}. Before $k+1 \geq 3$ no update is made: $\beta_U, \beta_L$
stay at $\beta_U^{(0)}, \beta_L^{(0)}$ and $\beta_{\mathrm{ok}}$ stays false.

\section{Post-processing}\label{sec:postproc}

A completed run leaves behind the log-scale history $\ln U_k, \ln L_k$ and the
clock history $\ln \tau_k$. From these we recover the parameters
$A_X, \alpha_X, \beta_X$ of the ansatz \eqref{eq:fw-powerlaw}. This is the
offline, global counterpart of \textsc{EstimateBeta}
(\S\ref{ssec:alg-beta-estimation}), which fits an effective exponent one
step at a time with $\ln A_X = 0$; here the whole history is used at once,
which separates the prefactor from the exponent, and the fit also serves as a
check on the run.

Write
\begin{equation}\label{eq:pp-logs}
  \mathcal{L}_k = \ln L_k, \qquad \mathcal{U}_k = \ln U_k,
\end{equation}
together with the per-step log-ratios $\ln \lambda_{X,k} = \ln X_k - \ln X_{k-1}$
of \eqref{eq:fw-step-ratios}, which for a Strang step combine both half-step
ratios. The stored scales are their cumulative sums,
\begin{equation}\label{eq:pp-cumsum}
  \mathcal{L}_k = \mathcal{L}_0 + \sum_{i=1}^{k} \ln \lambda_{L,i},
  \qquad
  \mathcal{U}_k = \mathcal{U}_0 + \sum_{i=1}^{k} \ln \lambda_{U,i},
\end{equation}
which grow only linearly in $k$; unlike $L_k$ and $U_k$ themselves,
$\mathcal{L}_k$ and $\mathcal{U}_k$ never overflow.

\paragraph{The fitted range.} The power law describes the run only once the
initial transient has passed. We discard the record before a cutoff time
$t_{\mathrm{cut}}$, the same discard as in step~1 of
\S\ref{ssec:stat-batch-means}, and fit over the steps
$k_{\mathrm{c}} \le k \le K$, where $t_{k_{\mathrm{c}}} = t_{\mathrm{cut}}$
and $K$ is the last step. The origin of physical time is placed at the
cutoff,
\begin{equation}\label{eq:pp-tstar-origin}
  t_{k_{\mathrm{c}}}^* = 0,
\end{equation}
so that $\alpha_X$ is the offset of the power law at the start of the fitted
record, $X_{k_{\mathrm{c}}} = A_X\, \alpha_X^{\beta_X}$.

\subsection{The scaling rates under the ansatz}\label{ssec:pp-rates}

Differentiating the ansatz \eqref{eq:fw-powerlaw} with respect to $t$, and
using the clock \eqref{eq:fw-clock} to convert from $t^*$, the logarithmic
rates \eqref{eq:fw-rates} are
\begin{equation}\label{eq:pp-rate-chain}
  \sigma_X = \frac{1}{X}\frac{\mathrm{d} X}{\mathrm{d} t^*}\frac{\mathrm{d} t^*}{\mathrm{d} t}
  = \frac{\beta_X\, \tau}{t^* + \alpha_X},
  \qquad X \in \{U, L\},
\end{equation}
independent of the prefactor. Inverting the ansatz to eliminate the unknown
offset,
\begin{equation}\label{eq:pp-tstar-offset}
  t^* + \alpha_X = \left( \frac{X}{A_X} \right)^{1/\beta_X},
\end{equation}
and substituting into \eqref{eq:pp-rate-chain} gives the rates in terms of the
scales alone,
\begin{equation}\label{eq:pp-rate}
  \sigma_X = \beta_X\, \tau \left( \frac{X}{A_X} \right)^{-1/\beta_X}.
\end{equation}
The rate has the sign of $\beta_X$. Taking logarithms of its magnitude,
\begin{equation}\label{eq:pp-rate-log}
  \ln |\sigma_{X,k}| = \ln |\beta_X| + \frac{\ln A_X}{\beta_X} + \ln \tau_k - \frac{\mathcal{X}_k}{\beta_X},
  \qquad \mathcal{X}_k = \ln X_k .
\end{equation}
Finally, the rates themselves are estimated directly from the log-ratios
already in hand, without any finite difference of $L_k$ or $U_k$ themselves,
\begin{equation}\label{eq:pp-rate-fd}
  \sigma_{X,k} \approx \frac{\ln \lambda_{X,k+1}}{\Delta t_k},
  \qquad X \in \{U, L\},
\end{equation}
The rates, not the log-ratios $\ln \lambda_{X,k+1}$ themselves, are also
what the run's diagnostics plot, since a ratio scales with the step and the
rates do not, so runs at different $\Delta t$ compare directly.

\subsection{Recovering \texorpdfstring{$A_X$, $\alpha_X$ and $\beta_X$}{A\_X, alpha\_X and beta\_X}}\label{ssec:pp-fit}

Equation \eqref{eq:pp-rate-log} is linear in $\mathcal{X}_k$ once the clock is
divided out. Regressing
\begin{equation}\label{eq:pp-regression}
  \ln |\sigma_{X,k}| - \ln \tau_k = b_X - \frac{1}{\beta_X}\, \mathcal{X}_k,
  \qquad
  b_X = \ln |\beta_X| + \frac{\ln A_X}{\beta_X},
\end{equation}
against $\mathcal{X}_k$ by least squares over the fitted range gives the
slope $-1/\beta_X$ and the intercept $b_X$, and hence
\begin{equation}\label{eq:pp-beta-A}
  \beta_X = -\frac{1}{\text{slope}},
  \qquad
  \ln A_X = \beta_X \left( b_X - \ln |\beta_X| \right).
\end{equation}
The exponent is read off the slope alone and does not depend on the
prefactor; the prefactor is whatever the intercept leaves over once
$\beta_X$ is known. The scatter of \eqref{eq:pp-regression} about the fitted
line is the first check on the ansatz: on a power law the points fall on a
straight line, and a range that includes the initial transient, before the
frame has locked on, bends it.

With $A_X$ and $\beta_X$ in hand, \eqref{eq:pp-tstar-offset} gives the
power-law time at each step,
\begin{equation}\label{eq:pp-theta}
  \theta_{X,k} = \left( \frac{X_k}{A_X} \right)^{1/\beta_X}
  = \exp\!\left( \frac{\mathcal{X}_k - \ln A_X}{\beta_X} \right),
\end{equation}
which should equal $t_k^* + \alpha_X$. The physical time is accumulated from
the per-step increments of \eqref{eq:fw-dtstar}, with the constant-step
weights \eqref{eq:fw-dtstar-weights-const} (the first two steps of the run,
which lack a three-point history, take $\Delta t_i^* \approx \Delta t_i \tau_i$),
from the origin \eqref{eq:pp-tstar-origin},
\begin{equation}\label{eq:pp-tstar-cumsum}
  t_k^* = \sum_{i=k_{\mathrm{c}}}^{k-1} \Delta t_i^* , \qquad k \ge k_{\mathrm{c}} .
\end{equation}
In exact arithmetic $\theta_{X,k} - t_k^*$ is independent of $k$ and equals
$\alpha_X$. In practice $\theta_{X,k}$ inherits a \emph{relative} error from
the fit --- an error $\delta$ in $\ln A_X$ multiplies it by
$e^{-\delta/\beta_X}$ --- while $t_k^*$ grows roughly exponentially in $t$ on
the logarithmic clock, so the absolute scatter of $\theta_{X,k} - t_k^*$
grows with $\theta_{X,k}$. We therefore take $\alpha_X$ as the least-squares
fit of the relative misfit $(\theta_{X,k} - t_k^* - \alpha_X)/\theta_{X,k}$,
\begin{equation}\label{eq:pp-alpha}
  \alpha_X = \frac{\sum_k \theta_{X,k}^{-2} \left( \theta_{X,k} - t_k^* \right)}
                  {\sum_k \theta_{X,k}^{-2}},
\end{equation}
summed over the fitted range. The weights concentrate it at the start of the
range, where $t_k^*$ is comparable to $\alpha_X$ and the offset is resolved at
all. The root-mean-square of the relative misfit
$(t_k^* + \alpha_X)/\theta_{X,k} - 1$ over the range is the second check on
the ansatz: it measures how well the power law, with all three parameters
fixed, reproduces the physical time the clock accumulates.

For the 1D diffusion equation the clock is $\tau = L^2/D$, and
\eqref{eq:pp-rate-log} specializes to
\begin{equation}\label{eq:pp-1d-length-rate}
  \ln \sigma_{L,k} = \ln \frac{\beta_L}{D} + \frac{\ln A_L}{\beta_L} + \left( 2 - \frac{1}{\beta_L} \right) \mathcal{L}_k,
\end{equation}
\begin{equation}\label{eq:pp-1d-velocity-rate}
  \ln |\sigma_{U,k}| = \ln \frac{|\beta_U|}{D} + \frac{\ln A_U}{\beta_U} + 2\, \mathcal{L}_k - \frac{\mathcal{U}_k}{\beta_U},
\end{equation}
so that for $X = L$ the regression \eqref{eq:pp-regression} collapses onto a
single variable, while for $X = U$ it is a two-variable fit in
$\mathcal{L}_k$ and $\mathcal{U}_k$.

\section{The 1D diffusion equation}\label{sec:1d}

The first instantiation of the framework is the one-dimensional diffusion
equation, $d = 1$, with a single scalar component; we drop the component index
throughout this section. It serves as a testbed: every operator the framework
needs is available in closed form, so the only approximation in a run is the
operator splitting itself.

\subsection{Governing equation and clock}\label{ssec:1d-clock}

The dimensional equation is
\begin{equation}\label{eq:1d-dimensional}
  \frac{\partial u^*}{\partial t^*} = D \frac{\partial^2 u^*}{\partial (x^*)^2},
\end{equation}
so that, in the notation of \eqref{eq:fw-general-pde},
\begin{equation}\label{eq:1d-rhs}
  R^*(u^*) = D \frac{\partial^2 u^*}{\partial (x^*)^2}.
\end{equation}
Under the rescaling \eqref{eq:fw-nondim} this term becomes
\begin{equation}\label{eq:1d-rhs-rescaled}
  D \frac{\partial^2 (u U)}{\partial (x^*)^2}
  = D U \frac{\partial}{\partial x}\!\left( \frac{\partial u}{\partial x} \frac{1}{L} \right) \frac{1}{L}
  = \frac{D U}{L^2} \frac{\partial^2 u}{\partial x^2},
\end{equation}
having used $\mathrm{d}x/\mathrm{d}x^* = 1/L$. The clock is then chosen to make
the physical term of the lemma \eqref{eq:fw-lemma} carry unit coefficient,
\begin{equation}\label{eq:1d-clock}
  \frac{\tau}{U} \cdot \frac{D U}{L^2} = 1
  \quad\Longrightarrow\quad
  \tau = \frac{L^2}{D},
\end{equation}
which is the diffusive time scale of the current length scale, as expected.
In log form, as required by \eqref{eq:fw-ln-dtstar},
\begin{equation}\label{eq:1d-ln-clock}
  \ln \tau_k = 2 \ln L_k - \ln D .
\end{equation}
Substituting \eqref{eq:1d-clock} into the rescaling lemma
\eqref{eq:fw-lemma} gives the equation the algorithm integrates,
\begin{equation}\label{eq:1d-transformed}
  \frac{\partial u}{\partial t}
  = \sigma_L\, x \frac{\partial u}{\partial x}
  - \sigma_U\, u
  + \frac{\partial^2 u}{\partial x^2}.
\end{equation}

\subsection{The physics step}\label{ssec:1d-physics-step}

Because the diffusion term in \eqref{eq:1d-transformed} has unit coefficient,
the physics operator \eqref{eq:fw-physics-flow},
\begin{equation}\label{eq:1d-physics-flow}
  \Phi_{\Delta t}^{(\mathrm{P})}(u_k)
  = u_k + \int_{t_k}^{t_k + \Delta t} \frac{\partial^2 u}{\partial x^2}\, \mathrm{d}t,
\end{equation}
has the closed-form spectral solution
\begin{equation}\label{eq:1d-propagator}
  \hat u_j^{(\mathrm{P})} = \hat u_j\, e^{-\kappa_j^2 \Delta t},
  \qquad \kappa_j = \frac{2\pi j}{H}, \quad j = 0, \dots, \left\lfloor N/2 \right\rfloor,
\end{equation}
where $\hat u_j = \mathrm{RFFT}(u)_j$. This is an exact one-shot spectral
propagator, not an approximate time-marching scheme, so no sub-stepping or
implicit solve is needed, and the accuracy requirement on \textsc{Evolve}
(\S\ref{ssec:alg-substeps}) is met trivially at any order. Inverse-transforming
gives the advanced field $u^{(\mathrm{P})}(x_n)$.

\subsection{Moment constraints}\label{ssec:1d-moments}

The constraint functionals of \S\ref{ssec:fw-constraints} are taken to be the
zeroth and the normalized second moment of the field,
\begin{equation}\label{eq:1d-functionals}
  F_1[u] = \int u\, \mathrm{d}x \equiv M_0,
  \qquad
  F_2[u] = \frac{1}{M_0} \int x^2 u\, \mathrm{d}x \equiv M_2,
\end{equation}
where the $M_0$ in $F_2$ is the target value \eqref{eq:alg-targets}, a
constant, not the zeroth moment of the field the functional is applied to. Thus
$M_0$ measures the amount of the conserved quantity and $M_2$, a mean-square
width, its spread. Writing $M_0^{\mathrm{in}}, M_2^{\mathrm{in}}$ for the moments of the field being scaled, and
substituting $y = x\, \lambda_L$ in each integral, the scaled field carries
\begin{equation}\label{eq:1d-m0-homogeneity}
  \int \frac{1}{\lambda_U} u\!\left( x\, \lambda_L \right) \mathrm{d}x
  = \frac{1}{\lambda_U \lambda_L}\, M_0^{\mathrm{in}},
\end{equation}
\begin{equation}\label{eq:1d-m2-homogeneity}
  \int x^2\, \frac{1}{\lambda_U} u\!\left( x\, \lambda_L \right) \mathrm{d}x
  = \frac{1}{\lambda_U \lambda_L^{3}}\, M_2^{\mathrm{in}},
\end{equation}
so the homogeneity exponents of \eqref{eq:fw-homogeneity} are
\begin{equation}\label{eq:1d-exponents}
  (p_1, q_1) = (1, 1), \qquad (p_2, q_2) = (1, 3),
\end{equation}
with $p_1 q_2 - p_2 q_1 = 2 \neq 0$, so the constraint system
\eqref{eq:fw-constraint-system} is uniquely solvable: the two moments do
respond independently to amplitude and to length rescaling. Solving it,
\begin{equation}\label{eq:1d-length-ratio}
  \lambda_L = \sqrt{\frac{M_0\, M_2^{\mathrm{in}}}{M_0^{\mathrm{in}}\, M_2}},
\end{equation}
\begin{equation}\label{eq:1d-velocity-ratio}
  \lambda_U = \frac{M_0^{\mathrm{in}}}{M_0\, \lambda_L}
  = \left(\frac{M_0^{\mathrm{in}}}{M_0}\right)^{3/2} \sqrt{\frac{M_2}{M_2^{\mathrm{in}}}} .
\end{equation}

Because the normalization of $F_2$ is a constant, these ratios are the same
as for the raw second moment $\int x^2 u\, \mathrm{d}x$; written in terms of the
raw moments, $M_0$ cancels from $M_2^{\mathrm{in}}/M_2$.

\subsection{Spectral interpolation on the real transform}\label{ssec:1d-interpolation}

For $d = 1$ and a real-valued field, the generic interpolation formula
\eqref{eq:fw-spectral-interp} is realized on the non-negative frequency
components of a real Fast Fourier Transform ($\mathrm{RFFT}$):
\begin{equation}\label{eq:1d-interp}
  \begin{split}
    \Phi_{\Delta t}^{(\mathrm{S})}(u_k)(x) = \frac{1}{\lambda_{U,k+1}}\biggl\{ &\hat{u}_0^{(k)}
      + 2 \sum_{j=1}^{N/2-1} (-1)^j \biggl[ \operatorname{Re}\bigl(\hat{u}_j^{(k)}\bigr) \cos\left(\frac{2\pi j \lambda_{L,k+1} x}{H}\right) \\
      &\qquad - \operatorname{Im}\bigl(\hat{u}_j^{(k)}\bigr) \sin\left(\frac{2\pi j \lambda_{L,k+1} x}{H}\right) \biggr] \\
      &+ (-1)^{N/2}\, \hat{u}_{N/2}^{(k)} \cos\left(\frac{\pi N \lambda_{L,k+1} x}{H}\right) \biggr\},
  \end{split}
\end{equation}
where $\hat{u}_j^{(k)} = \frac{1}{N} \mathrm{RFFT}(u_k)_j$ for
$j = 0, 1, \dots, N/2$, and the Nyquist term at $j = N/2$ is included only if
$N$ is even. The factor of $2$ accounts for the Hermitian conjugate modes
omitted by the real transform. \eqref{eq:1d-interp} holds for
$|\lambda_{L,k+1} x| \le H/2$; outside that the cutoff
\eqref{eq:fw-dilate-cutoff} sets $\Phi_{\Delta t}^{(\mathrm{S})}(u_k)(x) = 0$.

The alternating sign $(-1)^j$ is the phase-origin correction, and it is not
optional. The transform phases its coefficients to the sample index $n$, whose
origin $n = 0$ sits at $x = -H/2$ on the centred grid \eqref{eq:alg-grid}
rather than at $x = 0$. Since the scaling operator \eqref{eq:fw-scale-step}
dilates about the centre of the domain, each mode must be evaluated at
$\lambda_{L,k+1} x + H/2$ --- dilate about $x = 0$, then shift to the
transform's phase origin --- and not at $\lambda_{L,k+1}(x + H/2)$, which
dilates the shift along with the coordinate and agrees with it only at
$\lambda_{L,k+1} = 1$. As $H/2$ is half a period, that shift contributes the
constant phase $\exp(\mathrm{i} \kappa_j H / 2) = (-1)^j$. Dropping it altogether
returns $u_k\!\left(\lambda_{L,k+1} x - H/2\right)$, the wanted profile
displaced by half the domain, rather than the
$u_k\!\left(\lambda_{L,k+1} x\right)$ that \eqref{eq:fw-scale-step} calls for.
The same factor carries to each axis in three dimensions
(\S\ref{ssec:ns-interpolation}).

\subsection{Boundary treatment}\label{ssec:1d-boundary}

The substep \textsc{ApplyBC} of \S\ref{ssec:alg-substeps} imposes a homogeneous
Dirichlet condition by clamping the two edge grid points of the centred domain
\eqref{eq:alg-grid} after each completed step,
\begin{equation}\label{eq:1d-clamp}
  u_{k+1}(x_0) \gets 0, \qquad u_{k+1}(x_{N-1}) \gets 0 .
\end{equation}
On the centred grid $x \in [-H/2, H/2)$ these are two distinct grid points, not
periodic images of one another.

\emph{Open.} This sits in tension with the rest of the discretization: the
spectral propagator \eqref{eq:1d-propagator} presumes a periodic field, whereas
\eqref{eq:1d-clamp} imposes a Dirichlet condition. The spectral interpolation
\eqref{eq:1d-interp} no longer does at the edges: with the cutoff
\eqref{eq:fw-dilate-cutoff} it treats everything beyond the box as zero, which
is consistent with the clamp, but it is still a periodic interpolant inside the
box. The two are consistent only
so long as $u$ is negligible near the domain edges, which requires the initial
condition to be localized well inside the domain and the run to be stopped
before the profile spreads to the boundary. The scaling step works against
this --- it is what keeps the profile from spreading in the self-similar frame
--- but the condition is an assumption of the scheme rather than a consequence
of it, and is not currently monitored.

\section{The Navier--Stokes equations}\label{sec:ns}

The second instantiation is the incompressible Navier--Stokes equations, with
$d = 3$ and $u_i$ the velocity field. The initial condition and the
three-dimensional spectral interpolation used by \textsc{Scale} are given in
Appendix~\ref{app:ns-numerics}.
\subsection{Governing equations and clock}\label{ssec:ns-clock}

The equations are solved as a large-eddy simulation (LES). Let
$\mathcal{G}_{\Delta^*}$ denote a spatial filter of width $\Delta^*$. A tilde
marks a quantity of the filtered (resolved) field, and the same symbol without
a tilde the corresponding quantity of the full field: $\tilde u_i^* =
\mathcal{G}_{\Delta^*}(u_i^*)$ and $\tilde p^*$ are the filtered velocity and
pressure, while $u_i^*$ and $p^*$ are the full ones. For a quantity derived
from the velocity, the tilde means it is evaluated from $\tilde u_i^*$, not
that it is filtered; thus $\tfrac12 \tilde u_i^* \tilde u_i^*$ is not
$\mathcal{G}_{\Delta^*}(\tfrac12 u_i^* u_i^*)$. The effect of the unresolved
scales enters through the kinematic subgrid-scale stress
\begin{equation}\label{eq:ns-sgs-def}
  \tau_{ij}^* = \mathcal{G}_{\Delta^*}(u_i^* u_j^*) - \tilde u_i^* \tilde u_j^*,
\end{equation}
which, like the eddy viscosity below, exists only in the LES and carries no
tilde. The dimensional filtered equations are
\begin{equation}\label{eq:ns-momentum}
  \frac{\partial \tilde u_i^*}{\partial t^*} + \tilde u_j^* \frac{\partial \tilde u_i^*}{\partial x_j^*}= -\frac{1}{\rho}\frac{\partial \tilde p^*}{\partial x_i^*} + \nu^* \frac{\partial^2 \tilde u_i^*}{\partial x_j^* \partial x_j^*} - \frac{\partial \tau_{ij}^*}{\partial x_j^*},
\end{equation}
together with the continuity equation,
\begin{equation}\label{eq:ns-continuity}
  \frac{\partial \tilde u_i^*}{\partial x_i^*} = 0,
\end{equation}
so that, in the notation of \eqref{eq:fw-general-pde},
\begin{equation}\label{eq:ns-rhs}
  R_i^*(\tilde u^*) = - \tilde u_j^* \frac{\partial \tilde u_i^*}{\partial x_j^*}
  - \frac{1}{\rho}\frac{\partial \tilde p^*}{\partial x_i^*}
  + \nu^* \frac{\partial^2 \tilde u_i^*}{\partial x_j^* \partial x_j^*}
  - \frac{\partial \tau_{ij}^*}{\partial x_j^*}.
\end{equation}
In addition to the rescaling \eqref{eq:fw-nondim} of $\tilde u_i$ and $x_j$, the
pressure is nondimensionalized by the dynamic pressure,
\begin{equation}\label{eq:ns-pressure-nondim}
  \tilde p = \frac{\tilde p^*}{\rho U^2},
\end{equation}
and the subgrid-scale stress, which carries units of velocity squared, by
$U^2$,
\begin{equation}\label{eq:ns-sgs-nondim}
  \tau_{ij} = \frac{\tau_{ij}^*}{U^2}.
\end{equation}
With the filter width nondimensionalized as $\Delta = \Delta^*/L$, the
definitions carry over unchanged, $\tilde u_i = \mathcal{G}_\Delta(u_i)$ and
$\tau_{ij} = \mathcal{G}_\Delta(u_i u_j) - \tilde u_i \tilde u_j$.
Term by term, the convective, pressure, and subgrid-scale terms scale as
$U^2/L$,
\begin{equation}\label{eq:ns-convection}
  \tilde u_j^* \frac{\partial \tilde u_i^*}{\partial x_j^*} = \frac{U^2}{L} \tilde u_j \frac{\partial \tilde u_i}{\partial x_j},
  \qquad
  \frac{1}{\rho}\frac{\partial \tilde p^*}{\partial x_i^*} = \frac{U^2}{L} \frac{\partial \tilde p}{\partial x_i},
  \qquad
  \frac{\partial \tau_{ij}^*}{\partial x_j^*} = \frac{U^2}{L} \frac{\partial \tau_{ij}}{\partial x_j},
\end{equation}
while the viscous term scales as $\nu^* U / L^2$,
\begin{equation}\label{eq:ns-viscous}
  \nu^* \frac{\partial^2 \tilde u_i^*}{\partial x_j^* \partial x_j^*} = \frac{\nu^* U}{L^2}\frac{\partial^2 \tilde u_i}{\partial x_j \partial x_j}.
\end{equation}
The clock is chosen as the inverse of the advection timescale
\begin{equation}\label{eq:ns-clock}
  \tau = \frac{L}{U}.
\end{equation}
In log form, as required by \eqref{eq:fw-ln-dtstar},
\begin{equation}\label{eq:ns-ln-clock}
  \ln \tau_k = \ln L_k - \ln U_k .
\end{equation}
Applying $\tau / U = L/U^2$ to \eqref{eq:ns-rhs} and introducing the
instantaneous Reynolds number
\begin{equation}\label{eq:ns-reynolds}
  Re = \frac{U L}{\nu^*},
\end{equation}
the rescaling lemma \eqref{eq:fw-lemma} gives the nondimensional momentum
equation
\begin{equation}\label{eq:ns-nondim-momentum-base}
  \frac{\partial \tilde u_i}{\partial t}
  = \sigma_L\, x_j \frac{\partial \tilde u_i}{\partial x_j}
  - \sigma_U\, \tilde u_i
  - \tilde u_j \frac{\partial \tilde u_i}{\partial x_j}
  - \frac{\partial \tilde p}{\partial x_i}
  + \frac{1}{Re} \frac{\partial^2 \tilde u_i}{\partial x_j \partial x_j}
  - \frac{\partial \tau_{ij}}{\partial x_j},
\end{equation}
and the continuity equation is unchanged by the rescaling,
\begin{equation}\label{eq:ns-nondim-continuity}
  \frac{\partial \tilde u_i}{\partial x_i} = 0.
\end{equation}

\subsection{Subgrid stress model}\label{ssec:ns-viscosity}
The subgrid-scale stress $\tau_{ij}^*$ of \eqref{eq:ns-momentum} is closed with
an eddy-viscosity model. Its deviatoric part is taken proportional to the
resolved strain rate,
\begin{equation}\label{eq:ns-eddy-closure}
  \tau_{ij}^* - \frac{1}{3}\tau_{kk}^*\delta_{ij} = -2\nu_{\mathrm{t}}^* \tilde S_{ij}^*,
  \qquad
  \tilde S_{ij}^* = \frac{1}{2}\left(\frac{\partial \tilde u_i^*}{\partial x_j^*} + \frac{\partial \tilde u_j^*}{\partial x_i^*}\right),
\end{equation}
with the isotropic part $\frac{1}{3}\tau_{kk}^*\delta_{ij}$ absorbed into the
pressure. The eddy viscosity $\nu_{\mathrm{t}}^*$ carries the units of $UL$, so
it is nondimensionalized as $\nu_{\mathrm{t}} = \nu_{\mathrm{t}}^*/(UL)$. Noting
that $\tilde S_{ij}^* = (U/L)\, \tilde S_{ij}$, the nondimensional subgrid stress
\eqref{eq:ns-sgs-nondim} is
\begin{equation}\label{eq:ns-eddy-closure-nondim}
  \tau_{ij} - \frac{1}{3}\tau_{kk}\delta_{ij} = -2 \nu_{\mathrm{t}} \tilde S_{ij}.
\end{equation}
The models below differ only in the choice of $\nu_{\mathrm{t}}^*$.

\subsubsection{Smagorinsky model}\label{sssec:ns-sgs-smagorinsky}
The standard Smagorinsky model takes the eddy viscosity
\begin{equation}\label{eq:ns-eddy-viscosity}
  \nu_{\mathrm{t}}^* = (C_s \Delta^*)^2 |\tilde S^*|,
  \qquad
  |\tilde S^*| = \sqrt{2 \tilde S_{ij}^* \tilde S_{ij}^*},
\end{equation}
where $C_s$ is the Smagorinsky constant and $\Delta^*$ is the filter width.
With the nondimensional filter width $\Delta = \Delta^*/L$ taken equal to the
grid spacing, $\Delta = \Delta x$, the eddy viscosity
becomes
\begin{equation}\label{eq:ns-eddy-viscosity-nondim}
  \nu_{\mathrm{t}}^* = U L\, (C_s \Delta)^2 |\tilde S|,
  \qquad
  \nu_{\mathrm{t}} = (C_s \Delta)^2 |\tilde S|.
\end{equation}

\subsubsection{Self-similar constant eddy viscosity}\label{sssec:ns-sgs-constant}
The second model replaces the local eddy viscosity by a constant that is
uniform in space, chosen so that the dissipative scale of the flow tracks the
grid. Throughout this paper the dissipation rate $\epsilon^*$ is half the
trace of the dissipation tensor (\S\ref{sssec:ns-pp-quantities}),
$\epsilon^* = \nu^*\, \overline{(\partial u_i^*/\partial x_k^*)(\partial u_i^*/\partial x_k^*)}$
\eqref{eq:ns-pp-diss-trace}, the sink in the equation for
$\overline{k^*} = \tfrac{1}{2}\overline{u_i^* u_i^*}$, with $\nu_{\mathrm{t}}^*$ in place of
$\nu^*$ in the LES. It is not the thermodynamic dissipation rate
$2\nu^*\, \overline{S_{ij}^* S_{ij}^*}$, the rate at which kinetic energy is
converted to heat: the two differ by a divergence that vanishes in homogeneous
turbulence, where they are equal, but not in the inhomogeneous flow
considered here. With zero molecular viscosity, the Kolmogorov scale built
from the eddy viscosity and this $\epsilon^*$,
\begin{equation}\label{eq:ns-kolmogorov-scale}
  \eta^* = \left( \frac{(\nu_{\mathrm{t}}^*)^3}{\epsilon^*} \right)^{1/4},
\end{equation}
is required to be proportional to the filter width,
\begin{equation}\label{eq:ns-kolmogorov-mesh}
  \eta^* = C_\eta \Delta^*,
\end{equation}
with $C_\eta$ a constant. In homogeneous turbulence
\eqref{eq:ns-kolmogorov-scale} is the conventional Kolmogorov scale. The
dissipation rate is estimated from the large scales of the self-similar frame,
\begin{equation}\label{eq:ns-eps-estimate}
  \epsilon^* \approx C_\epsilon \frac{U^3}{L},
\end{equation}
with $C_\epsilon$ a constant, so that \eqref{eq:ns-kolmogorov-scale} and
\eqref{eq:ns-kolmogorov-mesh} give
\begin{equation}
  C_\eta \Delta^* = \left( \frac{(\nu_{\mathrm{t}}^*)^3 L}{C_\epsilon U^3} \right)^{1/4}.
\end{equation}
Solving for $\nu_{\mathrm{t}}^*$ and writing $\Delta^* = L \Delta$,
\begin{equation}\label{eq:ns-constant-viscosity}
  \nu_{\mathrm{t}}^* = C_\nu (\Delta^*)^{4/3} \frac{U}{L^{1/3}} = C_\nu \Delta^{4/3}\, U L,
  \qquad
  \nu_{\mathrm{t}} = C_\nu \Delta^{4/3},
\end{equation}
where
\begin{equation}\label{eq:ns-constant-viscosity-coeff}
  C_\nu = C_\eta^{4/3} C_\epsilon^{1/3},
\end{equation}
and $\Delta = \Delta^*/L = \Delta x$ as in the Smagorinsky model. The eddy
viscosity is uniform in space, so with \eqref{eq:ns-continuity} the subgrid
stress term reduces to a Laplacian,
\begin{equation}\label{eq:ns-constant-viscosity-laplacian}
  -\frac{\partial \tau_{ij}^*}{\partial x_j^*}
  = \nu_{\mathrm{t}}^* \frac{\partial^2 \tilde u_i^*}{\partial x_j^* \partial x_j^*},
  \qquad
  -\frac{\partial \tau_{ij}}{\partial x_j}
  = C_\nu \Delta^{4/3} \frac{\partial^2 \tilde u_i}{\partial x_j \partial x_j}.
\end{equation}
In the self-similar frame the model is therefore a constant viscosity, of the
same form as the molecular term of \eqref{eq:ns-nondim-momentum-base} with an
effective Reynolds number $Re_{\mathrm{t}} = 1/(C_\nu \Delta^{4/3})$. Unlike $Re$,
which changes with $UL$ as the flow evolves, $Re_{\mathrm{t}}$ is fixed for the
run, while in physical variables, where $\Delta^* = L\Delta$ grows with the
flow, $\nu_{\mathrm{t}}^*$ follows the product $UL$.

\subsubsection*{The momentum equation}
In the limit of infinite Reynolds number, $1/Re = 0$, and with either model the
momentum equation \eqref{eq:ns-nondim-momentum-base} becomes
\begin{equation}\label{eq:ns-nondim-momentum}
  \frac{\partial \tilde u_i}{\partial t}
  = \sigma_L\, x_j \frac{\partial \tilde u_i}{\partial x_j}
  - \sigma_U\, \tilde u_i
  - \tilde u_j \frac{\partial \tilde u_i}{\partial x_j}
  - \frac{\partial \tilde p}{\partial x_i}
  + \frac{\partial}{\partial x_j}\left[ 2 \nu_{\mathrm{t}}\, \tilde S_{ij} \right],
\end{equation}
where $\tilde p$ now denotes the modified pressure that includes
$\frac{1}{3}\tau_{kk}$, and $\nu_{\mathrm{t}}$ is given by
\eqref{eq:ns-eddy-viscosity-nondim} or \eqref{eq:ns-constant-viscosity}. For the
self-similar constant eddy viscosity the last term is
$C_\nu \Delta^{4/3}\, \partial^2 \tilde u_i / \partial x_j \partial x_j$, by
\eqref{eq:ns-constant-viscosity-laplacian}.

\subsection{Sponge forcing}\label{ssec:ns-sponge}
The true boundary conditions we would like to enforce is $\tilde u_i = 0$ at the edge of the domain. However, we are using a pseudo-spectral solver so
this boundary condition cannot be directly enforced. To solve this we introduce a fringe region at the edge of the domain \cite{nordstrom1999fringe},
which introduces a forcing term to the right hand side of the Navier--Stokes equation
\begin{equation}\label{eq:ns-sponge-force}
  f_i = -W(r)\, \tilde u_i,
\end{equation}
where $r = \sqrt{x_ix_i}$ is the distance from the origin and $W(r)$ is a weighting function that increases with increasing $r$.

The weighting function is
\begin{equation}\label{eq:ns-sponge-weight}
  W(r) = W_{\max}\, S\!\left( \frac{r - r_0}{r_1 - r_0} \right),
\end{equation}
where the smooth step $S$ is
\begin{equation}\label{eq:ns-sponge-step}
  S(\xi) =
  \begin{cases}
    0, & \xi \le 0, \\[4pt]
    \dfrac{1}{1 + \exp\!\left( \dfrac{1}{\xi - 1} + \dfrac{1}{\xi} \right)}, & 0 < \xi < 1, \\[10pt]
    1, & \xi \ge 1 .
  \end{cases}
\end{equation}
The function $S$ is $C^\infty$ on the whole real line. Every derivative of it
vanishes at $\xi = 0$ and at $\xi = 1$, it rises monotonically between them,
and it is antisymmetric about its midpoint, $S(1 - \xi) = 1 - S(\xi)$, so that
$S(1/2) = 1/2$. The weight \eqref{eq:ns-sponge-weight} is therefore exactly zero in
the sphere $r \le r_0$, rises smoothly across the shell $r_0 < r < r_1$, and
is exactly $W_{\max}$ for $r \ge r_1$. Adding \eqref{eq:ns-sponge-force} to
\eqref{eq:ns-nondim-momentum}, the momentum equation actually solved is
\begin{equation}\label{eq:ns-nondim-momentum-sponge}
  \frac{\partial \tilde u_i}{\partial t}
  = \sigma_L\, x_j \frac{\partial \tilde u_i}{\partial x_j}
  - \sigma_U\, \tilde u_i
  - \tilde u_j \frac{\partial \tilde u_i}{\partial x_j}
  - \frac{\partial \tilde p}{\partial x_i}
  + \frac{\partial}{\partial x_j}\left[ 2 \nu_{\mathrm{t}}\, \tilde S_{ij} \right]
  - W(r)\, \tilde u_i .
\end{equation}
\subsection{Fourier representation}\label{ssec:ns-fourier}
Using the Fourier convention \eqref{eq:fw-fourier-convention}, the pressure
gradient becomes
\begin{equation}\label{eq:ns-fourier-pressure}
  \widehat{\frac{\partial \tilde p}{\partial x_i}} = \mathrm{i} \kappa_i \hat{\tilde p}.
\end{equation}
The convective term
\begin{equation}\label{eq:ns-fourier-convection}
  \hat{N}_i = \widehat{\tilde u_j \frac{\partial \tilde u_i}{\partial x_j}},
\end{equation}
will be evaluated pseudo-spectrally. How the subgrid stress term
\begin{equation}\label{eq:ns-fourier-sgs-generic}
  \hat{G}_i = \widehat{\frac{\partial}{\partial x_j}\left[ 2 \nu_{\mathrm{t}} \tilde S_{ij} \right]}
\end{equation}
is evaluated depends on the model of \S\ref{ssec:ns-viscosity}. In both cases
the strain rate is available spectrally,
\begin{equation}\label{eq:ns-fourier-strain}
  \hat{\tilde S}_{ij} = \frac{\mathrm{i}}{2}\left( \kappa_j \hat{\tilde u}_i + \kappa_i \hat{\tilde u}_j \right).
\end{equation}

\paragraph{Smagorinsky model.} The eddy viscosity
\eqref{eq:ns-eddy-viscosity-nondim} depends on the field, so $\hat{G}_i$ is
evaluated pseudo-spectrally: $\tilde S_{ij}$ is transformed to physical space, the
products $|\tilde S|$, $\nu_{\mathrm{t}} = (C_s \Delta)^2 |\tilde S|$ and
$T_{ij} = 2 \nu_{\mathrm{t}} \tilde S_{ij}$ are formed there, and the divergence is
taken spectrally,
\begin{equation}\label{eq:ns-fourier-sgs}
  \hat{G}_i = \mathrm{i} \kappa_j \hat{T}_{ij}.
\end{equation}
Since $|\tilde S|$ is not a polynomial in $\tilde u$, $\hat{G}_i$ carries aliasing that the
$2/3$ rule does not remove entirely.

\paragraph{Self-similar constant eddy viscosity.} The eddy viscosity
\eqref{eq:ns-constant-viscosity} is uniform, and by
\eqref{eq:ns-constant-viscosity-laplacian}
\begin{equation}\label{eq:ns-fourier-sgs-constant}
  \hat{G}_i = -C_\nu \Delta^{4/3} \kappa^2 \hat{\tilde u}_i .
\end{equation}
This is exact and diagonal in $\boldsymbol{\kappa}$: no product is formed in physical
space, so $\hat{G}_i$ needs no transforms, carries no aliasing, and stays in
the retained band. By \eqref{eq:ns-fourier-continuity} it is also solenoidal,
$\kappa_i \hat{G}_i = 0$, so it contributes nothing to the pressure below and
passes through the projection \eqref{eq:ns-projection-tensor} unchanged.

\paragraph{Sponge forcing.} The sponge forcing
\eqref{eq:ns-sponge-force} is formed as a product in physical space
and transformed,
\begin{equation}\label{eq:ns-fourier-sponge}
  \hat{F}_i = \widehat{W \tilde u_i},
\end{equation}
with $W$ tabulated once on the grid \eqref{eq:alg-grid}. Because $W$ is smooth
but not band-limited, $\hat{F}_i$ also carries content beyond the retained
band, which the $2/3$ mask removes along with that of $\hat{N}_i$ and,
under the Smagorinsky model, $\hat{G}_i$. The scaling terms transform as in
\eqref{eq:fw-scaling-pde-fourier} with $d = 3$, so the transform of
\eqref{eq:ns-nondim-momentum-sponge} is
\begin{equation}\label{eq:ns-nondim-momentum-fourier}
  \frac{\partial \hat{\tilde u}_i}{\partial t}
  = -\sigma_L \left( \kappa_j \frac{\partial \hat{\tilde u}_i}{\partial \kappa_j} + 3\, \hat{\tilde u}_i \right)
  - \sigma_U\, \hat{\tilde u}_i
  - \hat{N}_i
  - \mathrm{i} \kappa_i \hat{\tilde p}
  + \hat{G}_i
  - \hat{F}_i,
\end{equation}
and the transform of the continuity equation
\eqref{eq:ns-nondim-continuity} is
\begin{equation}\label{eq:ns-fourier-continuity}
  \kappa_i \hat{\tilde u}_i = 0.
\end{equation}

\subsection{Pressure and the solenoidal projection}\label{ssec:ns-projection}
Contracting \eqref{eq:ns-nondim-momentum-fourier} with $\kappa_i$ and using the
continuity constraint gives
\begin{equation}\label{eq:ns-pressure-contracted}
  0 = -\sigma_L \left( \kappa_i \kappa_j \frac{\partial \hat{\tilde u}_i}{\partial \kappa_j} + d\, \kappa_i \hat{\tilde u}_i \right)
  - \sigma_U\, \kappa_i \hat{\tilde u}_i
  - \kappa_i \hat{N}_i
  - \mathrm{i} \kappa^2 \hat{\tilde p}
  + \kappa_i \hat{G}_i
  - \kappa_i \hat{F}_i.
\end{equation}
The remaining scaling contribution also vanishes: differentiating
\eqref{eq:ns-fourier-continuity} with respect to $\kappa_j$ gives
\begin{equation}\label{eq:ns-continuity-dk}
  \frac{\partial}{\partial \kappa_j}\left( \kappa_i \hat{\tilde u}_i \right)
  = \hat{\tilde u}_j + \kappa_i \frac{\partial \hat{\tilde u}_i}{\partial \kappa_j} = 0
  \quad\Longrightarrow\quad
  \kappa_i \frac{\partial \hat{\tilde u}_i}{\partial \kappa_j} = -\hat{\tilde u}_j,
\end{equation}
so that
\begin{equation}
  \kappa_i \kappa_j \partial \hat{\tilde u}_i / \partial \kappa_j = -\kappa_j \hat{\tilde u}_j = 0.
\end{equation}
Therefore, the equation for pressure is
\begin{equation}\label{eq:ns-fourier-poisson}
  \mathrm{i} \kappa^2 \hat{\tilde p} = -\kappa_i \left( \hat{N}_i - \hat{G}_i + \hat{F}_i \right)
  \quad\Longrightarrow\quad
  \hat{\tilde p} = \mathrm{i}\,\frac{\kappa_l \left( \hat{N}_l - \hat{G}_l + \hat{F}_l \right)}{\kappa^2},
  \qquad \kappa \neq 0,
\end{equation}
with $\hat{\tilde p}$ at $\kappa = 0$ fixed arbitrarily, since only its gradient enters
the momentum equation. Substituting \eqref{eq:ns-fourier-poisson} into
\eqref{eq:ns-nondim-momentum-fourier}, the pressure term becomes
\begin{equation}\label{eq:ns-fourier-pressure-substituted}
  -\mathrm{i} \kappa_i \hat{\tilde p} = \frac{\kappa_i \kappa_l}{\kappa^2} \left( \hat{N}_l - \hat{G}_l + \hat{F}_l \right),
\end{equation}
which combines with $-\hat{N}_i + \hat{G}_i - \hat{F}_i$ to give the solenoidal
projection of the convective, subgrid stress, and sponge terms together. Defining the projection tensor
\begin{equation}\label{eq:ns-projection-tensor}
  P_{il}(\boldsymbol{\kappa}) = \delta_{il} - \frac{\kappa_i \kappa_l}{\kappa^2},
\end{equation}
the evolution equation in fourier space is
\eqref{eq:fw-lemma},
\begin{equation}\label{eq:ns-fourier-final}
  \frac{\partial \hat{\tilde u}_i}{\partial t}
  = \underbrace{-\sigma_L \left( \kappa_j \frac{\partial \hat{\tilde u}_i}{\partial \kappa_j} + 3\, \hat{\tilde u}_i \right)
  - \sigma_U\, \hat{\tilde u}_i}_{\text{scaling terms}}
  \underbrace{- P_{il}(\boldsymbol{\kappa}) \left( \hat{N}_l - \hat{G}_l + \hat{F}_l \right)}_{\text{physical terms}}.
\end{equation}

\subsection{The physics step}\label{ssec:ns-physics-step}

The physics operator \eqref{eq:fw-physics-flow} is the flow of the bracketed
physical terms of \eqref{eq:ns-fourier-final},
\begin{equation}\label{eq:ns-physics-flow}
  \Phi_{\Delta t}^{(\mathrm{P})}\bigl(\hat{\tilde u}_i(t)\bigr)
  = \hat{\tilde u}_i(t) - \int_{t}^{t+\Delta t}
    P_{il}(\boldsymbol{\kappa}) \left( \hat{N}_l - \hat{G}_l + \hat{F}_l \right) \mathrm{d}t .
\end{equation}
We take $\Phi^{(\mathrm{P})}$ to be a single step of the classical fourth-order
Runge--Kutta method (RK4) applied to the full physical right-hand side, the
subgrid stress term $\hat{G}_i$ and the sponge forcing $\hat{F}_i$ marched
together with the convection. The convective term has eigenvalues on the
imaginary axis with $|\lambda| \le \max|\tilde{\mathbf{u}}|\, \kappa_{\max}$,
where $\kappa_{\max} = 2\pi (N/3)/H$ is the top mode retained by the $2/3$
rule, and RK4 is stable on the imaginary axis for $|\lambda \Delta t| \le 2\sqrt{2}$.
This gives the advective limit
\begin{equation}\label{eq:ns-cfl-dt}
  \Delta t \le \Delta t_{\mathrm{CFL}}(\tilde u) = \frac{2\sqrt{2}}{\max|\tilde{\mathbf{u}}|\, \kappa_{\max}},
\end{equation}
which depends on the state and is monitored through the realized CFL number
\eqref{eq:alg-realized-cfl} of \S\ref{ssec:alg-timestep}. The sponge forcing
and, under the self-similar constant eddy viscosity, the subgrid term are
linear decays with eigenvalues on the negative real axis, where RK4 is stable
for $|\lambda \Delta t| \lesssim 2.785$. This gives the sponge limit
$\Delta t\, W_{\max} \le 2.785$ and the diffusive limit
$\Delta t\, C_\nu \Delta^{4/3} \kappa_{\max}^2 \le 2.785$, both fixed by the
run parameters. None of the three is enforced: the constant step
$\Delta t$ must be chosen to satisfy them.

\subsection{Constraint functionals}\label{ssec:ns-constraints}
The constraint functionals are built from the resolved kinetic energy, which
we define first.

\paragraph{Kinetic energy.}
The resolved kinetic energy is
\begin{equation}\label{eq:ns-kinetic-energy-nondim}
  \tilde k^* = U^2\, \tilde k,
  \qquad
  \tilde k = \frac{1}{2} \tilde u_i \tilde u_i .
\end{equation}
Contracting \eqref{eq:ns-sgs-def} gives the filtered kinetic energy of the
full field \cite{pope2000turbulent},
\begin{equation}\label{eq:ns-kinetic-energy-split}
  \mathcal{G}_\Delta(k) = \tilde k + k_{\mathrm{sgs}},
  \qquad
  k_{\mathrm{sgs}} = \frac{1}{2}\tau_{kk}.
\end{equation}
The isotropic part of $\tau_{ij}$ is absorbed into the pressure, so the solver
never computes $k_{\mathrm{sgs}}$ and we must estimate it. It enters neither
the integration nor the constraints below; we use it
only as a diagnostic of how much energy the resolved field misses. For both
subgrid stress models we estimate it with the Smagorinsky eddy viscosity
\eqref{eq:ns-eddy-viscosity-nondim}, so that the estimate is the same function
of the resolved field in either case. Equating the production
$\nu_{\mathrm{t}} |\tilde S|^2 = C_s^2 \Delta^2 |\tilde S|^3$ of subgrid
kinetic energy to its sink $C_d k_{\mathrm{sgs}}^{3/2}/\Delta$, with $C_d$ a
constant of order one, gives
\begin{equation}\label{eq:ns-sgs-energy}
  k_{\mathrm{sgs}} = \left( \frac{C_s^2}{C_d} \right)^{2/3} \Delta^2 |\tilde S|^2,
\end{equation}
with $k_{\mathrm{sgs}}^* = U^2 k_{\mathrm{sgs}}$.

\paragraph{The functionals.}
The length and velocity scales will be set such that the zeroth and normalized
second radial moments of the resolved energy density
$\tilde u_i \tilde u_i = 2 \tilde k$, twice the resolved kinetic energy per
unit mass \eqref{eq:ns-kinetic-energy-nondim}, are constant. Mathematically,
these moments are defined as
\begin{equation}\label{eq:ns-functionals}
  F_1[\tilde u] = \int \tilde u_i \tilde u_i\, \mathrm{d}\mathbf{x} \equiv E_0,
  \qquad
  F_2[\tilde u] = \frac{1}{E_0} \int r^2\, \tilde u_i \tilde u_i\, \mathrm{d}\mathbf{x} \equiv E_2,
\end{equation}
where $r^2 = x_j x_j$ is the squared distance from the centre of the domain
\eqref{eq:alg-grid}, and the $E_0$ in $F_2$ is the target value
\eqref{eq:alg-targets}, a constant. As in the 1D case, $E_0$ measures the
amplitude of the field and $E_2$ its spread.
\par
The functionals are built from the resolved field alone: the subgrid kinetic
energy $k_{\mathrm{sgs}}$ of \eqref{eq:ns-kinetic-energy-split} is neglected,
so $E_0$ is twice the resolved kinetic energy rather than that of the full
field, and $E_2$ is the mean-square radius of the resolved energy. This is
accurate to the extent that the resolved field carries most of the energy. We
monitor the resolved fraction
$K_{\mathrm{res}} / (K_{\mathrm{res}} + K_{\mathrm{sgs}})$, with
$K_{\mathrm{res}} = \int \tilde k\, \mathrm{d}\mathbf{x}$ and
$K_{\mathrm{sgs}} = \int k_{\mathrm{sgs}}\, \mathrm{d}\mathbf{x}$ from the
estimate \eqref{eq:ns-sgs-energy}, a posteriori against the criterion that a
well-resolved LES carries at least $80\%$ of the kinetic energy in the
resolved field \cite{pope2000turbulent}.
\par
Writing $E_0^{\mathrm{in}}, E_2^{\mathrm{in}}$ for the moments of the field
being scaled and substituting $\mathbf{y} = \mathbf{x}\, \lambda_L$ in each
integral, with $d = 3$, the scaled field carries
\begin{equation}\label{eq:ns-e0-homogeneity}
  \int \frac{1}{\lambda_U^2}\, \tilde u_i \tilde u_i\!\left( \mathbf{x}\, \lambda_L \right) \mathrm{d}\mathbf{x}
  = \frac{1}{\lambda_U^2 \lambda_L^{3}}\, E_0^{\mathrm{in}},
\end{equation}
\begin{equation}\label{eq:ns-e2-homogeneity}
  \frac{1}{E_0} \int r^2\, \frac{1}{\lambda_U^2}\, \tilde u_i \tilde u_i\!\left( \mathbf{x}\, \lambda_L \right) \mathrm{d}\mathbf{x}
  = \frac{1}{\lambda_U^2 \lambda_L^{5}}\, E_2^{\mathrm{in}},
\end{equation}
so the homogeneity exponents of \eqref{eq:fw-homogeneity} are
\begin{equation}\label{eq:ns-exponents}
  (p_1, q_1) = (2, 3), \qquad (p_2, q_2) = (2, 5),
\end{equation}
with $p_1 q_2 - p_2 q_1 = 4 \neq 0$, so the constraint system
\eqref{eq:fw-constraint-system} is uniquely solvable. Solving it,
\begin{equation}\label{eq:ns-length-ratio}
  \lambda_L = \sqrt{\frac{E_0\, E_2^{\mathrm{in}}}{E_0^{\mathrm{in}}\, E_2}},
\end{equation}
\begin{equation}\label{eq:ns-velocity-ratio}
  \lambda_U = \sqrt{\frac{E_0^{\mathrm{in}}}{E_0\, \lambda_L^{3}}} ,
\end{equation}
where the length ratio is identical in form to its 1D counterpart
\eqref{eq:1d-length-ratio}.

\subsection{Clock control of the scaling step}\label{ssec:ns-clock-control}
The constrained ratios \eqref{eq:ns-length-ratio}, \eqref{eq:ns-velocity-ratio}
fix $E_0, E_2$ exactly at every step, and so pass every fluctuation of the
moments of the evolved field into $\ln \lambda_U, \ln \lambda_L$, and through
them into the rates $\sigma_X$. For a turbulent field these fluctuations are
large from step to step. In particular the clock advances over step $k$ by
$\Delta_{\ln\tau,k} = \ln \tau_{k+1} - \ln \tau_k$, at the rate
\begin{equation}\label{eq:ns-cc-rate}
  c_k = \frac{\Delta_{\ln\tau,k}}{\Delta t},
\end{equation}
which is $\sigma_L - \sigma_U$ by \eqref{eq:ns-ln-clock}, and on the power law
the constant $c_\tau$ of \eqref{eq:ns-pp-rates-clock}. Without clock control
$c_k$ is whatever the constraints produce; clock control prescribes it, makes
it a slowly varying state of the integrator, and lets the second constraint
absorb the difference. \par
Throughout this subsection $\Delta^{\bullet}_{\ln\tau,k}$ denotes an increment
of $\ln \tau$ over step $k$, from $t_k$ to $t_{k+1}$, and
$c^{\bullet}_k = \Delta^{\bullet}_{\ln\tau,k} / \Delta t$ the corresponding
rate, with the same label $\bullet$. Like $\Delta t_k$, both carry the index
$k$ of the step they belong to.


\paragraph{The clock and the miss.}
By \eqref{eq:ns-ln-clock} ratios $\lambda_U, \lambda_L$ advance the clock by
$\ln \lambda_L - \ln \lambda_U$. For a field $\tilde u$, let
$\ln \lambda_X^{\mathrm{con}}[\tilde u]$ be the constrained ratios
\eqref{eq:ns-length-ratio}, \eqref{eq:ns-velocity-ratio} with
$E_m^{\mathrm{in}} = E_m[\tilde u]$. The \emph{miss}
\begin{equation}\label{eq:ns-cc-miss-def}
  \mu[\tilde u] = \ln \lambda_L^{\mathrm{con}}[\tilde u] - \ln \lambda_U^{\mathrm{con}}[\tilde u]
\end{equation}
is the amount by which $\ln \tau$ would have to advance to put $\tilde u$
back on the targets with no physics in between. 

\paragraph{The clock rate.}
Under clock control the rate obeys
\begin{equation}\label{eq:ns-cc-ode}
  \frac{\mathrm{d} c}{\mathrm{d} t} = \gamma \left( c^{\mathrm{tgt}} - c \right),
\end{equation}
with the gain $\gamma > 0$ and the rate $c_{-1} = c_{\mathrm{init}}$ the first
step relaxes from as run parameters. Within step $k$, let
$\Delta^{\mathrm{con}}_{\ln\tau,k}$ be the increment that would put the evolved
field on the targets. It consists of
the miss $\mu_k$ the step started with and the increment the physics added
over the step,
$\Delta^{\mathrm{phys}}_{\ln\tau,k} = \Delta^{\mathrm{con}}_{\ln\tau,k} - \mu_k$.
The target
\begin{equation}\label{eq:ns-cc-target}
  c^{\mathrm{tgt}}_k = c^{\mathrm{phys}}_k + \frac{\mu_k}{T_\mu}
  = \frac{\Delta^{\mathrm{con}}_{\ln\tau,k} - \mu_k}{\Delta t} + \frac{\mu_k}{T_\mu},
  \qquad T_\mu = \frac{4}{\gamma},
\end{equation}
follows the physics and closes the miss over the time $T_\mu$. Holding it over
the step, \eqref{eq:ns-cc-ode} is integrated exactly,
\begin{equation}\label{eq:ns-cc-update}
  c_k = c^{\mathrm{tgt}}_k
  + \left( c_{k-1} - c^{\mathrm{tgt}}_k \right) e^{-\gamma \Delta t},
\end{equation}
which is stable for any $\gamma \Delta t$, and the step advances $\ln \tau$ by
$\Delta_{\ln\tau,k} = c_k \Delta t$.

\paragraph{The ratios.}
The clock fixes one combination of $\ln \lambda_U, \ln \lambda_L$; the other
holds $E_0$ exactly. With the exponents \eqref{eq:ns-exponents}, the ratios
applied solve
\begin{equation}\label{eq:ns-cc-system}
  \begin{bmatrix} 2 & 3 \\ -1 & 1 \end{bmatrix}
  \begin{bmatrix} \ln \lambda_U \\ \ln \lambda_L \end{bmatrix}
  = \begin{bmatrix} \ln E_0^{\mathrm{in}} - \ln E_0 \\ c_k \Delta t - \Delta^{\mathrm{pred}}_{\ln\tau,k} \end{bmatrix},
\end{equation}
with $\Delta^{\mathrm{pred}}_{\ln\tau,k} = 0$ under \textsc{LieTrotter}. The
determinant is $5$, and in closed form
\begin{equation}\label{eq:ns-cc-ratios}
  \ln \lambda_L = \frac{R_0 + 2 R_\tau}{5}, \qquad
  \ln \lambda_U = \frac{R_0 - 3 R_\tau}{5},
\end{equation}
with $R_0$ and $R_\tau$ the two entries of the right-hand side. Of all ratios
that advance the clock as prescribed, these are the ones closest to the
constrained ratios in the sense that the first constraint is still met
exactly: the whole difference goes into $E_2$, which the rescaled field
misses by
\begin{equation}\label{eq:ns-cc-f2}
  \ln \frac{E_2[\tilde u_{k+1}]}{E_2} = \frac{4}{5}\, \mu_{k+1},
\end{equation}
where $4 = p_1 q_2 - p_2 q_1$ is the determinant of the constraint system
\eqref{eq:fw-constraint-system}. As with the projection inside
\textsc{Scale}, the next \textsc{ConstrainedScale} measures the moments
afresh, so the miss is carried forward through $\mu_{k+1}$ and not lost.

\paragraph{Dynamics of the miss.}
Over a step the miss changes by what the physics added less what the clock
applied,
\begin{equation}\label{eq:ns-cc-miss-step}
  \mu_{k+1} = \mu_k + \Delta^{\mathrm{phys}}_{\ln\tau,k} - \Delta_{\ln\tau,k}
  = \mu_k + \bigl( c^{\mathrm{phys}}_k - c_k \bigr) \Delta t .
\end{equation}
Let $\Delta t \to 0$ with $\gamma$ and $T_\mu$ held fixed, and assume the
physics advances the clock by $O(\Delta t)$ per step, so that
$c^{\mathrm{phys}}_k$ has a limit. The rates then become functions $c(t)$,
$c^{\mathrm{phys}}(t)$ and the miss a function $\mu(t)$, and dividing
\eqref{eq:ns-cc-miss-step} by $\Delta t$,
\begin{equation}\label{eq:ns-cc-limit}
  \dot \mu = c^{\mathrm{phys}} - c .
\end{equation}
Substituting the target \eqref{eq:ns-cc-target},
$c^{\mathrm{tgt}} = c^{\mathrm{phys}} + \mu/T_\mu$, into \eqref{eq:ns-cc-ode}
and then using \eqref{eq:ns-cc-limit},
\begin{equation}\label{eq:ns-cc-cdot}
  \dot c = \gamma \left( c^{\mathrm{phys}} - c \right) + \frac{\gamma}{T_\mu}\, \mu
  = \gamma \dot \mu + \frac{\gamma}{T_\mu}\, \mu .
\end{equation}
Differentiating \eqref{eq:ns-cc-limit} and inserting
\eqref{eq:ns-cc-cdot},
$\ddot \mu = \dot c^{\mathrm{phys}} - \dot c
= \dot c^{\mathrm{phys}} - \gamma \dot \mu - (\gamma/T_\mu)\, \mu$, or
\begin{equation}\label{eq:ns-cc-miss}
  \ddot \mu + \gamma \dot \mu + \frac{\gamma}{T_\mu}\, \mu = \dot c^{\mathrm{phys}},
\end{equation}
a damped oscillator. Matching it to the standard form
$\ddot \mu + 2 \zeta \omega_0 \dot \mu + \omega_0^2 \mu$ gives the natural
frequency $\omega_0 = \sqrt{\gamma/T_\mu}$ and the damping ratio
$\zeta = \gamma / (2 \omega_0) = \tfrac12 \sqrt{\gamma T_\mu}$.
Critical damping, $\zeta = 1$, requires $\gamma T_\mu = 4$, so the choice
$T_\mu = 4/\gamma$ makes it critically damped, independently of
$\Delta t$. Closing the whole miss in one step instead, $T_\mu = \Delta t$,
would give $\zeta = \tfrac12 \sqrt{\gamma \Delta t} \ll 1$, an oscillation at
the frequency $\sqrt{\gamma/\Delta t}$ that grows as the step is refined. \par
The natural frequency is then $\omega_0 = \gamma/2$, so the clock rate follows
$c^{\mathrm{phys}}$ through a low-pass filter with corner frequency
$\gamma/2$: a smaller $\gamma$ gives a smoother clock at the cost of a larger
miss in $E_2$.


\section{HIT Post Processing thought experiment}
In HIT the kinetic energy evolution equation reduces to
\begin{equation}
  \frac{\partial k}{\partial t} = -\epsilon,
\end{equation}
which holds for both the physical flow and the simulated flow, each with its own
$k$ and $\epsilon$. For the simulated flow, with kinetic energy $k_s$ and eddy
viscosity $\nu_{\mathrm{t}}$, the dissipation equation is
\begin{equation}
  \frac{\partial \epsilon}{\partial t} = \frac{1}{\nu_{\mathrm{t}}}\frac{\partial \nu_{\mathrm{t}}}{\partial t} \epsilon + P_\epsilon - Y
  \approx \frac{1}{\nu_{\mathrm{t}}}\frac{\partial \nu_{\mathrm{t}}}{\partial t} \epsilon - C_{\epsilon 2} \frac{\epsilon^2}{k_s}.
\end{equation}
The time dependent viscosity $\nu_{\mathrm{t}}$ moves the tail of the kinetic
energy spectrum $E(\kappa)$ away from where the constant physical viscosity
$\nu_0$ would put it. Let $\eta_0$ be the Kolmogorov scale of the physical
flow, built from $\nu_0$, and $\eta$ the Kolmogorov scale of the simulated
flow, built from $\nu_{\mathrm{t}}$. Consequently the kinetic energy can be broken up into two contributions
\begin{equation}
  k = k_s -  k_a,
\end{equation}
where $k_s$ is the simulated kinetic energy
\begin{equation}
  k_s \approx \int_0^{2\pi/\eta} E(\kappa)\, \mathrm{d}\kappa,
\end{equation}
and $k_a$ is the kinetic energy added by the viscosity adjustment
\begin{equation}
  k_a \approx \int_{2\pi/\eta_0}^{2\pi/\eta} E(\kappa)\, \mathrm{d}\kappa.
\end{equation}
The decomposition holds for either ordering of the two scales. When
$\eta < \eta_0$ the simulation resolves more of the tail than the physical
flow has and $k_a > 0$; when $\eta > \eta_0$ the simulation truncates the
tail and $k_a < 0$ is the energy it is missing.
Take the length scale $L$ to be constant. Assume that in the inertial range the spectrum is self-similar,
$E(\kappa, t) = U^2 L\, \hat E(\kappa L)$, with
\begin{equation}
  \hat E(q) = \hat A (q + \hat\alpha)^{\beta_E},
\end{equation}
so that in dimensional form
\begin{equation}
  E(\kappa) = A_E(\kappa + \alpha_E)^{\beta_E},
  \qquad
  A_E = \hat A\, U^2 L^{1 + \beta_E},
  \qquad
  \alpha_E = \frac{\hat\alpha}{L}.
\end{equation}
For a Kolmogorov spectrum, $\beta_E = -5/3$, $\hat\alpha = 0$ and, with
\eqref{eq:ns-eps-estimate}, $\hat A = C_K C_\epsilon^{2/3}$, where $C_K$ is the Kolmogorov constant.
Using
\begin{equation}
  \eta = C_\eta \Delta L,
\end{equation}
and
\begin{equation}
  \eta_0 = \left(\frac{\nu_0^3 L}{C_\epsilon U^3}\right)^{1/4}.
\end{equation}
Plugging in
\begin{equation}
  k_a = \int_{2\pi / \eta_0}^{2\pi/\eta} A_E(\kappa + \alpha_E)^{\beta_E}\, \mathrm{d}\kappa = \frac{A_E}{\beta_E+1} \left[\left(\frac{2\pi}{\eta} + \alpha_E\right)^{\beta_E + 1} - \left(\frac{2\pi}{\eta_0} + \alpha_E\right)^{\beta_E + 1}\right],
\end{equation}
for $\beta_E \neq -1$.
Plugging in equations for Kolmogorov scale
\begin{equation}\label{eq:hit-ka}
  k_a = \frac{A_E}{\beta_E + 1}\left[\left(\frac{2\pi}{C_\eta \Delta L} + \alpha_E\right)^{\beta_E+1} - \left(\frac{2\pi C_\epsilon^{1/4} U^{3/4}}{\nu_0^{3/4}L^{1/4}} + \alpha_E\right)^{\beta_E + 1}\right].
\end{equation}

\subsection{The dissipation correction}\label{ssec:hit-eps}
The dissipation splits in the same way,
\begin{equation}\label{eq:hit-eps-split}
  \epsilon = \epsilon_s - \epsilon_a,
\end{equation}
where $\epsilon_s$ is the dissipation of the simulated flow and $\epsilon_a$
the dissipation added by the viscosity adjustment. There are two ways to
estimate $\epsilon_a$: from the spectrum, as for $k_a$, and from the energy
budget.

\paragraph{From the spectrum.}
With $\epsilon$ half the trace of the dissipation tensor,
$\epsilon = \nu\, \overline{(\partial u_i/\partial x_k)(\partial u_i/\partial x_k)}
= 2\nu \int_0^\infty \kappa^2 E(\kappa)\, \mathrm{d}\kappa$, so
\begin{equation}
  \epsilon_s \approx 2\nu_{\mathrm{t}} \int_0^{2\pi/\eta} \kappa^2 E(\kappa)\, \mathrm{d}\kappa,
  \qquad
  \epsilon \approx 2\nu_0 \int_0^{2\pi/\eta_0} \kappa^2 E(\kappa)\, \mathrm{d}\kappa,
\end{equation}
and
\begin{equation}\label{eq:hit-eps-a-parts}
  \epsilon_a = 2(\nu_{\mathrm{t}} - \nu_0) \int_0^{2\pi/\eta} \kappa^2 E(\kappa)\, \mathrm{d}\kappa
  + 2\nu_0 \int_{2\pi/\eta_0}^{2\pi/\eta} \kappa^2 E(\kappa)\, \mathrm{d}\kappa.
\end{equation}
Unlike $k_a$, which comes only from the band between the two cutoffs,
$\epsilon_a$ has a second part from the change of viscosity over the whole
resolved range. Take $\alpha_E = 0$ and $\beta_E > -3$, so that the integrals
converge at $\kappa = 0$. Then
\begin{equation}
  \epsilon_s = \frac{2 A_E}{\beta_E + 3}\, \nu_{\mathrm{t}} \left(\frac{2\pi}{\eta}\right)^{\beta_E + 3},
\end{equation}
and likewise for $\epsilon$ with $\nu_0$ and $\eta_0$. With the self-similar
constant eddy viscosity \eqref{eq:ns-constant-viscosity}, $\eta = C_\eta \Delta L$
is the Kolmogorov scale of $\nu_{\mathrm{t}}$ and $C_\epsilon U^3/L$, just as
$\eta_0$ is the Kolmogorov scale of $\nu_0$ and $C_\epsilon U^3/L$.
Substituting $2\pi/\eta = 2\pi (C_\epsilon U^3/(\nu_{\mathrm{t}}^3 L))^{1/4}$
and its counterpart for $\eta_0$,
\begin{equation}\label{eq:hit-eps-a-spectral}
  \epsilon_a = \frac{2 A_E}{\beta_E + 3} (2\pi)^{\beta_E + 3}
  \left(\frac{C_\epsilon U^3}{L}\right)^{(\beta_E + 3)/4}
  \left[\nu_{\mathrm{t}}^{-(3\beta_E + 5)/4} - \nu_0^{-(3\beta_E + 5)/4}\right].
\end{equation}
For the Kolmogorov spectrum, $\beta_E = -5/3$, the bracket vanishes and
$\epsilon_a = 0$ for any $\nu_{\mathrm{t}}$: the dissipation is fixed by the
flux through the inertial range, not by the viscosity, and moving the cutoff
changes $k$ but not $\epsilon$. A nonzero spectral $\epsilon_a$ comes only from
departures from $-5/3$ ($\beta_E \neq -5/3$ or $\alpha_E \neq 0$). The sharp
cutoff is crude here: with $\beta_E = -5/3$ the spectral $\epsilon_s$ equals
$\tfrac32 (2\pi)^{4/3} C_K\, C_\epsilon U^3/L$, consistent with
\eqref{eq:ns-eps-estimate} only if $\tfrac32 (2\pi)^{4/3} C_K = 1$. The
cutoff constant cancels from \eqref{eq:hit-eps-a-spectral} but not from
\eqref{eq:hit-ka}.

\paragraph{From the energy budget.}
Both flows obey $\partial k/\partial t = -\epsilon$, so subtracting the two
with $k_a = k_s - k$ gives
\begin{equation}\label{eq:hit-eps-a-budget}
  \epsilon_a = -\frac{\partial k_a}{\partial t}.
\end{equation}
This is exact, given the two energy equations, and it does not vanish for the
Kolmogorov spectrum. The two estimates disagree because the spectral one
assumes that both flows share the same inertial range, $U^2 L \hat E$, which
cannot hold exactly when $k_s \neq k$; the error is of relative size
$k_a/k_s$, the same size as $\epsilon_a/\epsilon_s$ from
\eqref{eq:hit-eps-a-budget}. The budget therefore gives the leading
nonzero $\epsilon_a$. \par
For the Kolmogorov spectrum, $\hat A = C_K C_\epsilon^{2/3}$, and
\eqref{eq:hit-ka} with $\alpha_E = 0$ is
\begin{equation}\label{eq:hit-ka-kolmogorov}
  k_a = -\frac{3}{2} (2\pi)^{-2/3} C_K C_\epsilon^{2/3}\, U^2
  \left[\left(\frac{\eta}{L}\right)^{2/3} - \left(\frac{\eta_0}{L}\right)^{2/3}\right],
\end{equation}
with $\eta/L = C_\eta \Delta$ fixed by the grid and
$\eta_0/L = C_\epsilon^{-1/4} \mathrm{Re}_0^{-3/4}$, $\mathrm{Re}_0 = UL/\nu_0$,
so that $(\eta_0/L)^{2/3} \propto U^{-1/2}$. With $L$ and $\Delta$ constant,
\begin{equation}
  \frac{\partial k_a}{\partial t} = -\frac{3}{2} (2\pi)^{-2/3} C_K C_\epsilon^{2/3}\,
  U \frac{\partial U}{\partial t}
  \left[2\left(\frac{\eta}{L}\right)^{2/3} - \frac{3}{2}\left(\frac{\eta_0}{L}\right)^{2/3}\right].
\end{equation}
The rate $\partial U/\partial t$ follows from the simulated flow: with $L$
constant, self-similarity gives $k_s \propto U^2$, so
$U\, \partial U/\partial t = \tfrac12 (U^2/k_s)\, \partial k_s/\partial t
= -\tfrac12 (U^2/k_s)\, \epsilon_s$. Then \eqref{eq:hit-eps-a-budget} gives
\begin{equation}\label{eq:hit-eps-a}
  \epsilon_a = -\frac{3}{4} (2\pi)^{-2/3} C_K C_\epsilon^{2/3}\,
  \frac{U^2}{k_s}\, \epsilon_s
  \left[2\left(\frac{\eta}{L}\right)^{2/3} - \frac{3}{2}\left(\frac{\eta_0}{L}\right)^{2/3}\right].
\end{equation}
Unlike $k_a$, $\epsilon_a$ does not vanish when $\eta = \eta_0$: the grid holds
$\eta/L$ fixed while $\eta_0/L$ grows as $U$ decays, so the two cutoffs drift
apart and $k_a$ changes even at the instant they coincide. When
$(\eta_0/L)^{2/3} > \tfrac43 (\eta/L)^{2/3}$, which includes the
well-resolved case $\eta \ll \eta_0$, the bracket is negative and
$\epsilon_a > 0$: the simulated flow dissipates faster than the physical one
because the excess energy $k_a$ in its tail decays faster than $k$.

\subsection{Mesh dependance}

\section{Navier--Stokes results}\label{sec:ns-results}

\subsection{Post-processing}\label{ssec:ns-postproc}
This subsection sets out how statistics of the self-similar fields are
related to the physical flow and measured, defines the dissipation rate, and
gives the scaling rates in the stationary state. The budgets of
\S\ref{ssec:ns-pp-rs-budget}--\S\ref{ssec:ns-pp-eps-budget} build on all
three.

\subsubsection{Scaling of point statistics}\label{sssec:ns-pp-scaling}
Let $Q^*$ be a pointwise quantity that nondimensionalizes as
\begin{equation}\label{eq:ns-pp-generic}
  Q^* = U^{m_U} L^{m_L}\, Q,
\end{equation}
where the nondimensional $Q$ is evaluated from the self-similar fields.
Substituting the power-law ansatz \eqref{eq:fw-powerlaw} for
$U$ and $L$,
\begin{equation}\label{eq:ns-pp-powerlaw}
  Q^*(L\mathbf{x}, t^*) = A_U^{m_U} A_L^{m_L}\, (t^* + \alpha_U)^{m_U \beta_U} (t^* + \alpha_L)^{m_L \beta_L}
  \,Q(\mathbf{x}, t) .
\end{equation}
The origin, $\mathbf{x}^* = \mathbf{x} = 0$, is the only point fixed in both
frames ($n = N/2$ for even $N$). Elsewhere \eqref{eq:ns-pp-powerlaw} relates
the physical field to the self-similar one along the ray
$\mathbf{x}^* \propto L$, not at a fixed physical position.

\paragraph{Scaling of the ensemble average.}
Let $\overline{\,\cdot\,}$ denote the average over an ensemble of initial
conditions. We make two assumptions:
\begin{enumerate}
  \item The parameters $A_U, \alpha_U, \beta_U, A_L, \alpha_L, \beta_L$ are the same for
  every member of the ensemble, i.e.\ $U(t^*)$ and $L(t^*)$ do not depend on
  the initial condition.
  \item In the self-similar frame the flow is statistically stationary, so
  that $\overline{Q(\mathbf{x},t)}$ is independent of $t$.
\end{enumerate}
Then
\begin{equation}\label{eq:ns-pp-ensemble}
  \overline{Q^*}(L\mathbf{x}, t^*) = A_U^{m_U} A_L^{m_L}\, (t^* + \alpha_U)^{m_U \beta_U} (t^* + \alpha_L)^{m_L \beta_L}
  \,\overline{Q}(\mathbf{x}) .
\end{equation}
\par
 We assume that the ensemble statistics are isotropic and therfore only depend on the distance from the origin 
$r = \sqrt{x_i x_i}$ so \eqref{eq:ns-pp-ensemble} becomes
\begin{equation}\label{eq:ns-pp-ensemble-radial}
  \overline{Q^*}(r^*, t^*) = A_U^{m_U} A_L^{m_L}\, (t^* + \alpha_U)^{m_U \beta_U} (t^* + \alpha_L)^{m_L \beta_L}
  \,\overline{Q}(r^*/L) .
\end{equation}
Every point on the sphere $|\mathbf{x}| = r$ is then a sample of the same
statistic, so we can average over that sphere within a single snapshot. This
is the spherical mean,
\begin{equation}\label{eq:ns-pp-sphere-mean}
  \langle Q \rangle_r(t) = \frac{1}{4\pi} \oint_{|\hat{\mathbf{n}}| = 1}
  Q(r \hat{\mathbf{n}}, t)\, \mathrm{d}\Omega(\hat{\mathbf{n}}).
\end{equation}
The sphere does not pass through grid points, so
\eqref{eq:ns-pp-sphere-mean} is evaluated on the spectral interpolant of $Q$,
where it reduces to a weighted sum over the Fourier modes
(Appendix~\ref{app:radial}). At $r = 0$ it is the point value,
$\langle Q \rangle_0 = Q(0, t)$.
\par
We also assume that in the self similar frame the system is ergodic so the ensemble average equals a time average,
\begin{equation}\label{eq:ns-pp-time-average}
  \overline{Q}(r) = \frac{1}{T}\int_0^T \langle Q \rangle_r(t)\, \mathrm{d}t,
\end{equation}
and a single long simulation suffices. In practice $T$ is finite and the
initial transient is discarded, so \eqref{eq:ns-pp-time-average} only
estimates $\overline{Q}(r)$; Appendix~\ref{app:stat-error} describes how its
statistical error is computed. The spherical mean reduces this error compared
with a single point. The reduction is larger at larger $r$, where the sphere
spans more correlation lengths, and absent at $r = 0$.

\subsubsection{The dissipation rate}\label{sssec:ns-pp-quantities}
\paragraph{The dissipation tensor and its trace.}
As in \S\ref{ssec:ns-pp-k-budget}, we assume that the ensemble-mean velocity
vanishes, so $u_i^*$ is the fluctuating velocity. The evolution equation for the Reynolds
stress tensor $\overline{u_i^* u_j^*}$, with $\overline{\,\cdot\,}$ the ensemble
average of \S\ref{sssec:ns-pp-scaling}, contains the sink $-\varepsilon_{ij}^*$
(\S\ref{ssec:ns-pp-rs-budget}),
where
\begin{equation}\label{eq:ns-pp-diss-tensor}
  \varepsilon_{ij}^* = 2\nu^*\, \overline{\frac{\partial u_i^*}{\partial x_k^*} \frac{\partial u_j^*}{\partial x_k^*}}
\end{equation}
is the dissipation tensor \cite{pope2000turbulent}. We report half its trace,
\begin{equation}\label{eq:ns-pp-diss-trace}
  \epsilon^* = \tfrac{1}{2}\varepsilon_{ii}^* = \nu^*\, \overline{\frac{\partial u_i^*}{\partial x_k^*} \frac{\partial u_i^*}{\partial x_k^*}},
\end{equation}
which is the sink in the equation for $\overline{k^*} = \tfrac{1}{2}\overline{u_i^* u_i^*}$.
\paragraph{Estimate from the filtered field.}
The LES does not resolve the gradients in \eqref{eq:ns-pp-diss-trace}. As for
the kinetic energy below, we therefore estimate the filtered integrand,
$\mathcal{G}_{\Delta^*}$ of it. We split it into a resolved and a subgrid part,
\begin{equation}\label{eq:ns-pp-diss-split}
  \mathcal{G}_{\Delta^*}\!\left( 2\nu^*\, \frac{\partial u_i^*}{\partial x_k^*} \frac{\partial u_j^*}{\partial x_k^*} \right)
  = 2\nu^*\, \frac{\partial \tilde u_i^*}{\partial x_k^*} \frac{\partial \tilde u_j^*}{\partial x_k^*}
  + 2\nu^* \left[ \mathcal{G}_{\Delta^*}\!\left( \frac{\partial u_i^*}{\partial x_k^*} \frac{\partial u_j^*}{\partial x_k^*} \right)
  - \frac{\partial \tilde u_i^*}{\partial x_k^*} \frac{\partial \tilde u_j^*}{\partial x_k^*} \right],
\end{equation}
where the bracket is the unclosed subgrid gradient correlation. The
eddy-viscosity closure \eqref{eq:ns-eddy-closure} treats the subgrid scales as
an additional viscosity $\nu_{\mathrm{t}}^*$ acting on the resolved field. For
consistency we close the bracket the same way, by gradient transport,
\begin{equation}\label{eq:ns-pp-diss-closure}
  \mathcal{G}_{\Delta^*}\!\left( \frac{\partial u_i^*}{\partial x_k^*} \frac{\partial u_j^*}{\partial x_k^*} \right)
  \approx \left( 1 + \frac{\nu_{\mathrm{t}}^*}{\nu^*} \right)
  \frac{\partial \tilde u_i^*}{\partial x_k^*} \frac{\partial \tilde u_j^*}{\partial x_k^*}.
\end{equation}
The bracket in \eqref{eq:ns-pp-diss-split} then becomes
$(\nu_{\mathrm{t}}^*/\nu^*)\, (\partial \tilde u_i^*/\partial x_k^*)(\partial \tilde u_j^*/\partial x_k^*)$,
and
\begin{equation}\label{eq:ns-pp-diss-tensor-les}
  \mathcal{G}_{\Delta^*}\!\left( 2\nu^*\, \frac{\partial u_i^*}{\partial x_k^*} \frac{\partial u_j^*}{\partial x_k^*} \right)
  \approx 2\nu^*\, \frac{\partial \tilde u_i^*}{\partial x_k^*} \frac{\partial \tilde u_j^*}{\partial x_k^*}
  + 2\nu_{\mathrm{t}}^*\, \frac{\partial \tilde u_i^*}{\partial x_k^*} \frac{\partial \tilde u_j^*}{\partial x_k^*}
  = 2\left( \nu^* + \nu_{\mathrm{t}}^* \right) \frac{\partial \tilde u_i^*}{\partial x_k^*} \frac{\partial \tilde u_j^*}{\partial x_k^*}.
\end{equation}
Taking half the trace gives the pointwise estimate of the integrand of
\eqref{eq:ns-pp-diss-trace},
\begin{equation}\label{eq:ns-pp-diss-trace-les}
  \epsilon^* = \left( \nu^* + \nu_{\mathrm{t}}^* \right) \frac{\partial \tilde u_i^*}{\partial x_k^*} \frac{\partial \tilde u_i^*}{\partial x_k^*}.
\end{equation}
Since the filter commutes with the ensemble average, the average of
\eqref{eq:ns-pp-diss-trace-les} estimates \eqref{eq:ns-pp-diss-trace}. With
$\nu_{\mathrm{t}}^* = UL\, \nu_{\mathrm{t}}$, $\nu^* = UL/Re$, and
$\partial \tilde u_i^*/\partial x_k^* = (U/L)\, \partial \tilde u_i/\partial x_k$,
\begin{equation}\label{eq:ns-pp-diss-nondim}
  \epsilon^* = \frac{U^3}{L}\, \epsilon,
  \qquad
  \epsilon = \left( \frac{1}{Re} + \nu_{\mathrm{t}} \right) \frac{\partial \tilde u_i}{\partial x_k} \frac{\partial \tilde u_i}{\partial x_k},
\end{equation}
which is of the form \eqref{eq:ns-pp-generic} with $m_U = 3$, $m_L = -1$. In
the limit $1/Re = 0$ the resolved part vanishes and only the subgrid part
remains. By \eqref{eq:ns-eddy-viscosity-nondim} and
\eqref{eq:ns-constant-viscosity},
\begin{equation}\label{eq:ns-pp-diss-models}
  \epsilon = \nu_{\mathrm{t}}\, \frac{\partial \tilde u_i}{\partial x_k} \frac{\partial \tilde u_i}{\partial x_k} =
  \begin{cases}
    (C_s \Delta)^2 |\tilde S|\, \dfrac{\partial \tilde u_i}{\partial x_k} \dfrac{\partial \tilde u_i}{\partial x_k}, & \text{Smagorinsky model}, \\[10pt]
    C_\nu \Delta^{4/3}\, \dfrac{\partial \tilde u_i}{\partial x_k} \dfrac{\partial \tilde u_i}{\partial x_k}, & \text{self-similar constant eddy viscosity}.
  \end{cases}
\end{equation}

\subsubsection{Scaling rates in the stationary state}\label{sssec:ns-pp-rates}
The budgets below contain the rates $\sigma_U$ and $\sigma_L$
\eqref{eq:fw-rates}, which the power-law ansatz fixes. Take
$\beta_L - \beta_U = 1$ and a shared offset $\alpha_U = \alpha_L = \alpha$
in \eqref{eq:fw-powerlaw}. The clock \eqref{eq:ns-clock} is then
\begin{equation}\label{eq:ns-pp-rates-clock}
  \tau = \frac{L}{U} = c_\tau\, (t^* + \alpha),
  \qquad
  c_\tau = \frac{A_L}{A_U},
\end{equation}
and by \eqref{eq:pp-rate-chain} the rates are constant,
\begin{equation}\label{eq:ns-pp-rates}
  \sigma_X = \beta_X \frac{\tau}{t^* + \alpha} = c_\tau \beta_X ,
\end{equation}
so that
\begin{equation}\label{eq:ns-pp-rates-diff}
  \sigma_U - \sigma_L = c_\tau \left( \beta_U - \beta_L \right) = -c_\tau .
\end{equation}
Equivalently, by \eqref{eq:ns-ln-clock}, $c_\tau$ is the rate at which
$\ln \tau = \ln L - \ln U$ grows in $t$. In the budgets we use
\eqref{eq:ns-pp-rates} in place of the instantaneous rates; in a run whose
scales are prescribed by the power law it holds exactly. Under the clock
control of \S\ref{ssec:ns-clock-control}, \eqref{eq:ns-pp-rates-diff} holds at every
step with the instantaneous clock rate $c_k$ in place of $c_\tau$.

\subsection{Reynolds stress budget}\label{ssec:ns-pp-rs-budget}
This subsection derives the evolution equation of the resolved Reynolds
stress $\overline{\tilde u_i \tilde u_j}$. As for the $\epsilon$ budget of
\S\ref{ssec:ns-pp-eps-budget}, we proceed in three steps:
\begin{enumerate}
  \item write the dimensional budget, \eqref{eq:ns-pp-rs-rate};
  \item nondimensionalize it into the self-similar frame,
  \eqref{eq:ns-pp-rs-rate-nondim};
  \item assume stationarity and isotropy of the statistics, which reduces it
  to two balances between radial profiles,
  \eqref{eq:ns-pp-rs-radial-ctau}.
\end{enumerate}

\paragraph{Assumptions.}
The assumptions are those of the $\epsilon$ budget
(\S\ref{ssec:ns-pp-eps-budget}): $1/Re = 0$, the subgrid model is the
self-similar constant eddy viscosity \eqref{eq:ns-constant-viscosity}, the
ensemble-mean velocity vanishes, $\overline{\tilde u_i} = 0$, and we work
inside the sphere $r \le r_0$, where $W = 0$.

\paragraph{Dimensional budget.}
Under these assumptions the subgrid stress term is the Laplacian
\eqref{eq:ns-constant-viscosity-laplacian}, and the filtered momentum equation
\eqref{eq:ns-momentum} becomes
\begin{equation}\label{eq:ns-pp-rs-momentum}
  \frac{\partial \tilde u_i^*}{\partial t^*}
  + \tilde u_k^* \frac{\partial \tilde u_i^*}{\partial x_k^*}
  = -\frac{1}{\rho}\frac{\partial \tilde p^*}{\partial x_i^*}
  + \nu_{\mathrm{t}}^* \frac{\partial^2 \tilde u_i^*}{\partial x_k^* \partial x_k^*} ,
\end{equation}
with $\nu_{\mathrm{t}}^*$ uniform in space \eqref{eq:ns-pp-eps-nut}. Multiply
the equation for $\tilde u_i^*$ by $\tilde u_j^*$, add the same with $i$ and
$j$ exchanged, and average. By continuity \eqref{eq:ns-continuity}, the
advection term is the divergence
$\partial_k^* (\tilde u_i^* \tilde u_j^* \tilde u_k^*)$, and since
$\nu_{\mathrm{t}}^*$ is uniform the viscous term splits as
\begin{equation*}
  \nu_{\mathrm{t}}^* \left( \tilde u_j^* \frac{\partial^2 \tilde u_i^*}{\partial x_k^* \partial x_k^*}
  + \tilde u_i^* \frac{\partial^2 \tilde u_j^*}{\partial x_k^* \partial x_k^*} \right)
  = \nu_{\mathrm{t}}^* \frac{\partial^2 (\tilde u_i^* \tilde u_j^*)}{\partial x_k^* \partial x_k^*}
  - 2 \nu_{\mathrm{t}}^* \frac{\partial \tilde u_i^*}{\partial x_k^*} \frac{\partial \tilde u_j^*}{\partial x_k^*} .
\end{equation*}
Unlike $\epsilon^*$, the Reynolds stress does not contain $\nu_{\mathrm{t}}^*$,
so the change of the eddy viscosity in time does not enter. The dimensional
budget is \cite{pope2000turbulent}
\begin{equation}\label{eq:ns-pp-rs-rate}
  \frac{\partial \overline{\tilde u_i^* \tilde u_j^*}}{\partial t^*}
  = \tilde T_{ij}^* + \tilde \Pi_{ij}^* + \tilde D_{ij}^* - \varepsilon_{ij}^* ,
\end{equation}
with
\begin{equation}\label{eq:ns-pp-rs-terms}
\begin{aligned}
  \tilde T_{ij}^* &= -\frac{\partial}{\partial x_k^*}\, \overline{\tilde u_i^* \tilde u_j^* \tilde u_k^*}
  && \text{turbulent transport,} \\
  \tilde \Pi_{ij}^* &= -\frac{1}{\rho} \left( \overline{\tilde u_i^* \frac{\partial \tilde p^*}{\partial x_j^*}}
  + \overline{\tilde u_j^* \frac{\partial \tilde p^*}{\partial x_i^*}} \right)
  && \text{velocity--pressure gradient,} \\
  \tilde D_{ij}^* &= \nu_{\mathrm{t}}^* \frac{\partial^2 \overline{\tilde u_i^* \tilde u_j^*}}{\partial x_k^* \partial x_k^*}
  && \text{viscous diffusion,} \\
  \varepsilon_{ij}^* &= 2 \nu_{\mathrm{t}}^*\, \overline{\frac{\partial \tilde u_i^*}{\partial x_k^*} \frac{\partial \tilde u_j^*}{\partial x_k^*}}
  && \text{dissipation.}
\end{aligned}
\end{equation}
The dissipation is the tensor \eqref{eq:ns-pp-diss-tensor-les} with
$1/Re = 0$, so that $\tfrac12 \varepsilon_{ii}^* = \overline{\epsilon^*}$
\eqref{eq:ns-pp-eps-def}. With zero mean velocity the production vanishes,
and so does the mean advection. The velocity--pressure gradient term contains
both the pressure--strain redistribution and the pressure transport
\cite{pope2000turbulent}; we keep it whole, since only its sum enters.

\paragraph{Nondimensionalization.}
The substitutions are those of the $\epsilon$ budget, $\tilde u_i^* = U \tilde u_i$,
$x_j^* = L x_j$, $\tilde p^* = \rho U^2 \tilde p$, and
$\nu_{\mathrm{t}}^* = UL\, \nu_{\mathrm{t}}$ \eqref{eq:ns-pp-eps-nut}.
\begin{itemize}
  \item \emph{Right-hand side.} Each of the four terms carries the factor
  $U^3/L$. For example,
  $\tilde T_{ij}^* = (U^3/L)\, \tilde T_{ij}$ and
  $\varepsilon_{ij}^* = (UL)(U/L)^2\, \varepsilon_{ij} = (U^3/L)\, \varepsilon_{ij}$.
  \item \emph{Left-hand side.} The chain rule \eqref{eq:fw-lhs-chain}, with
  the amplitude $U^2$ (rate $2\sigma_U$) in place of $U$ and the clock
  $\tau = L/U$ \eqref{eq:ns-clock}, gives
  \begin{equation*}
    \frac{\partial \overline{\tilde u_i^* \tilde u_j^*}}{\partial t^*}
    = \frac{U^3}{L}\left( \frac{\partial \overline{\tilde u_i \tilde u_j}}{\partial t}
    + 2\sigma_U\, \overline{\tilde u_i \tilde u_j}
    - \sigma_L\, x_k \frac{\partial \overline{\tilde u_i \tilde u_j}}{\partial x_k} \right).
  \end{equation*}
\end{itemize}
Dividing by $U^3/L$ gives
\begin{equation}\label{eq:ns-pp-rs-rate-nondim}
  \frac{\partial \overline{\tilde u_i \tilde u_j}}{\partial t}
  = \sigma_L\, x_k \frac{\partial \overline{\tilde u_i \tilde u_j}}{\partial x_k}
  - 2\sigma_U\, \overline{\tilde u_i \tilde u_j}
  + \tilde T_{ij} + \tilde \Pi_{ij} + \tilde D_{ij} - \varepsilon_{ij} ,
\end{equation}
with the nondimensional terms
\begin{equation}\label{eq:ns-pp-rs-terms-nondim}
\begin{aligned}
  \tilde T_{ij} &= -\frac{\partial}{\partial x_k}\, \overline{\tilde u_i \tilde u_j \tilde u_k},
  &
  \tilde \Pi_{ij} &= -\overline{\tilde u_i \frac{\partial \tilde p}{\partial x_j}}
  - \overline{\tilde u_j \frac{\partial \tilde p}{\partial x_i}},
  \\
  \tilde D_{ij} &= \nu_{\mathrm{t}} \frac{\partial^2 \overline{\tilde u_i \tilde u_j}}{\partial x_k \partial x_k},
  &
  \varepsilon_{ij} &= 2 \nu_{\mathrm{t}}\, \overline{\frac{\partial \tilde u_i}{\partial x_k} \frac{\partial \tilde u_j}{\partial x_k}} ,
\end{aligned}
\end{equation}
of the same form as \eqref{eq:ns-pp-rs-terms} without stars. The two scaling
terms follow directly from those of the momentum equation
\eqref{eq:ns-nondim-momentum}: the first advects the stress radially as the
frame dilates, and the second, $-2\sigma_U$, is twice the velocity rescaling
$-\sigma_U \tilde u_i$, since the stress is quadratic in the velocity.

\paragraph{Stationary state.}
In the stationary state $\partial \overline{\tilde u_i \tilde u_j}/\partial t = 0$
and the rates are constant, $\sigma_X = c_\tau \beta_X$
(\S\ref{sssec:ns-pp-rates}), so \eqref{eq:ns-pp-rs-rate-nondim} becomes
\begin{equation}\label{eq:ns-pp-rs-stationary}
  0 = c_\tau \beta_L\, x_k \frac{\partial \overline{\tilde u_i \tilde u_j}}{\partial x_k}
  - 2 c_\tau \beta_U\, \overline{\tilde u_i \tilde u_j}
  + \tilde T_{ij} + \tilde \Pi_{ij} + \tilde D_{ij} - \varepsilon_{ij} .
\end{equation}
Unlike the $\epsilon$ budget \eqref{eq:ns-pp-eps-radial-ctau}, whose
rescaling term depends only on $\beta_U - \beta_L = -1$, the rescaling term
here depends on $\beta_U$ itself.

\paragraph{Isotropic form.}
Under the isotropy assumption of \S\ref{sssec:ns-pp-scaling}, the statistics
are invariant under rotations about the origin. A symmetric second-rank
tensor field with this invariance has the form
\begin{equation}\label{eq:ns-pp-rs-isotropic}
  \overline{\tilde u_i \tilde u_j}
  = R_\perp(r)\, \delta_{ij} + \left( R_r(r) - R_\perp(r) \right) n_i n_j ,
  \qquad
  n_i = \frac{x_i}{r} ,
\end{equation}
where $R_r = \overline{\tilde u_r \tilde u_r}$, with $\tilde u_r = n_i \tilde u_i$,
is the radial normal stress and $R_\perp$ the normal stress in either
direction tangent to the sphere; the off-diagonal components in spherical
coordinates vanish. Each term of \eqref{eq:ns-pp-rs-stationary} has the same
form, so for a term $X_{ij}$ we define its radial and tangential components
\begin{equation}\label{eq:ns-pp-rs-projections}
  X_r = n_i n_j X_{ij} ,
  \qquad
  X_\perp = \tfrac{1}{2} \left( \delta_{ij} - n_i n_j \right) X_{ij} ,
\end{equation}
which are evaluated as spherical means \eqref{eq:ns-pp-sphere-mean} of the
projected pointwise integrands. The trace is $X_{ii} = X_r + 2 X_\perp$.
\par
Two of the terms can be written in terms of $R_r$ and $R_\perp$. Since
$n_i$ is constant along rays, $x_k\, \partial_k n_i = 0$, and the advection
term preserves the form \eqref{eq:ns-pp-rs-isotropic},
\begin{equation*}
  x_k \frac{\partial \overline{\tilde u_i \tilde u_j}}{\partial x_k}
  = r \frac{\mathrm{d} R_\perp}{\mathrm{d} r}\, \delta_{ij}
  + r \frac{\mathrm{d} (R_r - R_\perp)}{\mathrm{d} r}\, n_i n_j .
\end{equation*}
For the diffusion, $\nabla^2 (n_i n_j) = (2\delta_{ij} - 6 n_i n_j)/r^2$ and
the cross term $2\, \partial_k (R_r - R_\perp)\, \partial_k (n_i n_j)$
vanishes, again because $n_i$ is constant along rays, so the projections
\eqref{eq:ns-pp-rs-projections} of $\tilde D_{ij}$ are
\begin{equation}\label{eq:ns-pp-rs-diffusion}
  \tilde D_r = \nu_{\mathrm{t}} \left( \nabla_r^2 R_r - \frac{4}{r^2} \left( R_r - R_\perp \right) \right),
  \qquad
  \tilde D_\perp = \nu_{\mathrm{t}} \left( \nabla_r^2 R_\perp + \frac{2}{r^2} \left( R_r - R_\perp \right) \right),
\end{equation}
with the radial Laplacian
$\nabla_r^2 f = r^{-2}\, \mathrm{d}(r^2\, \mathrm{d}f/\mathrm{d}r)/\mathrm{d}r$.
The extra terms in $R_r - R_\perp$ couple the two components through the
curvature of the spheres and cancel in the trace. Projecting
\eqref{eq:ns-pp-rs-stationary} with \eqref{eq:ns-pp-rs-projections} gives the
budget we evaluate,
\begin{equation}\label{eq:ns-pp-rs-radial-ctau}
\begin{aligned}
  0 &= c_\tau \beta_L\, r \frac{\mathrm{d} R_r}{\mathrm{d} r}
  - 2 c_\tau \beta_U\, R_r
  + \tilde T_r + \tilde \Pi_r + \tilde D_r - \varepsilon_r , \\
  0 &= c_\tau \beta_L\, r \frac{\mathrm{d} R_\perp}{\mathrm{d} r}
  - 2 c_\tau \beta_U\, R_\perp
  + \tilde T_\perp + \tilde \Pi_\perp + \tilde D_\perp - \varepsilon_\perp ,
\end{aligned}
\end{equation}
each term a function of $r$ alone. As for the $\epsilon$ budget, every term is
estimated independently, so the residuals of \eqref{eq:ns-pp-rs-radial-ctau}
measure the statistical and discretization error of the profiles. Regularity
at the origin requires $R_r(0) = R_\perp(0)$ and
$\mathrm{d}R_r/\mathrm{d}r = \mathrm{d}R_\perp/\mathrm{d}r = 0$ at $r = 0$,
where the stress is isotropic.

\subsection{Kinetic energy budget}\label{ssec:ns-pp-k-budget}
This subsection derives the evolution equation of the mean resolved kinetic
energy $\overline{\tilde k} = \tfrac12 \overline{\tilde u_i \tilde u_i}$,
following the same three steps. The assumptions are those of
\S\ref{ssec:ns-pp-rs-budget}. Since $\tilde k$ is half the trace of the
Reynolds stress, each step is half the trace of its counterpart there; we
give the result of each step and the form the terms take.

\paragraph{Dimensional budget.}
Half the trace of \eqref{eq:ns-pp-rs-rate} is
\begin{equation}\label{eq:ns-pp-k-rate}
  \frac{\partial \overline{\tilde k^*}}{\partial t^*}
  = \tilde T_k^* + \tilde \Pi_k^* + \tilde D_k^* - \overline{\epsilon^*} ,
\end{equation}
with
\begin{equation}\label{eq:ns-pp-k-terms}
\begin{aligned}
  \tilde T_k^* &= -\frac{\partial}{\partial x_j^*}\, \overline{\tilde u_j^* \tilde k^*}
  && \text{turbulent transport,} \\
  \tilde \Pi_k^* &= -\frac{1}{\rho} \frac{\partial}{\partial x_j^*}\, \overline{\tilde u_j^* \tilde p^*}
  && \text{pressure transport,} \\
  \tilde D_k^* &= \nu_{\mathrm{t}}^* \frac{\partial^2 \overline{\tilde k^*}}{\partial x_j^* \partial x_j^*}
  && \text{viscous diffusion,} \\
  \overline{\epsilon^*} &= \nu_{\mathrm{t}}^*\, \overline{\frac{\partial \tilde u_i^*}{\partial x_m^*} \frac{\partial \tilde u_i^*}{\partial x_m^*}}
  && \text{dissipation,}
\end{aligned}
\end{equation}
where $\tilde k^* = \tfrac12 \tilde u_i^* \tilde u_i^*$. The pressure term is
half the trace of $\tilde \Pi_{ij}^*$, which by continuity
\eqref{eq:ns-continuity} is a divergence: the pressure--strain part is
traceless and only redistributes energy among the components. The sink is the
$\epsilon^*$ \eqref{eq:ns-pp-eps-def} whose budget is formed in
\S\ref{ssec:ns-pp-eps-budget}. It appears here with the Laplacian
$\tilde D_k^*$ of $\overline{\tilde k^*}$, rather than with the thermodynamic
dissipation and the divergence of the viscous stress work, which is why
$\epsilon^*$ is defined as half the trace of the dissipation tensor
(\S\ref{sssec:ns-sgs-constant}).

\paragraph{Nondimensionalization.}
With the substitutions of \S\ref{ssec:ns-pp-rs-budget}, the kinetic energy
and the dissipation scale as
\begin{equation}\label{eq:ns-pp-k-scales}
  \overline{\tilde k^*} = U^2\, \overline{\tilde k} ,
  \qquad
  \overline{\epsilon^*} = \frac{U^3}{L}\, \overline{\epsilon} ,
\end{equation}
by \eqref{eq:ns-kinetic-energy-nondim} and \eqref{eq:ns-pp-eps-def}. Each term
on the right of \eqref{eq:ns-pp-k-rate} carries the factor $U^3/L$, and the
left-hand side transforms as in \S\ref{ssec:ns-pp-rs-budget}, with the
amplitude $U^2$. Dividing by $U^3/L$,
\begin{equation}\label{eq:ns-pp-k-rate-nondim}
  \frac{\partial \overline{\tilde k}}{\partial t}
  = \sigma_L\, x_j \frac{\partial \overline{\tilde k}}{\partial x_j}
  - 2\sigma_U\, \overline{\tilde k}
  + \tilde T_k + \tilde \Pi_k + \tilde D_k - \overline{\epsilon} ,
\end{equation}
with the nondimensional terms
\begin{equation}\label{eq:ns-pp-k-terms-nondim}
  \tilde T_k = -\frac{\partial}{\partial x_j}\, \overline{\tilde u_j \tilde k},
  \qquad
  \tilde \Pi_k = -\frac{\partial}{\partial x_j}\, \overline{\tilde u_j \tilde p},
  \qquad
  \tilde D_k = \nu_{\mathrm{t}} \frac{\partial^2 \overline{\tilde k}}{\partial x_j \partial x_j},
  \qquad
  \overline{\epsilon} = \nu_{\mathrm{t}}\, \overline{\frac{\partial \tilde u_i}{\partial x_m} \frac{\partial \tilde u_i}{\partial x_m}} .
\end{equation}

\paragraph{Stationary state.}
In the stationary state, with $\sigma_X = c_\tau \beta_X$,
\begin{equation}\label{eq:ns-pp-k-stationary}
  0 = c_\tau \beta_L\, x_j \frac{\partial \overline{\tilde k}}{\partial x_j}
  - 2 c_\tau \beta_U\, \overline{\tilde k}
  + \tilde T_k + \tilde \Pi_k + \tilde D_k - \overline{\epsilon} .
\end{equation}

\paragraph{Isotropic form.}
Under the isotropy assumption, $\overline{\tilde k}$ and $\overline{\epsilon}$
depend only on $r$, and the mean fluxes $\overline{\tilde u_j \tilde k}$ and
$\overline{\tilde u_j \tilde p}$ are radial, $F_j = n_j F(r)$. As in
\S\ref{sssec:ns-pp-eps-model}, the advection term becomes
$x_j\, \partial_j \overline{\tilde k} = r\, \mathrm{d}\overline{\tilde k}/\mathrm{d}r$
and the divergence of a radial flux is $r^{-2}\, \mathrm{d}(r^2 F)/\mathrm{d}r$,
so the transport terms are
\begin{equation}\label{eq:ns-pp-k-radial-terms}
  \tilde T_k = -\frac{1}{r^2} \frac{\mathrm{d}}{\mathrm{d} r}\left( r^2\, \overline{\tilde u_r \tilde k} \right),
  \qquad
  \tilde \Pi_k = -\frac{1}{r^2} \frac{\mathrm{d}}{\mathrm{d} r}\left( r^2\, \overline{\tilde u_r \tilde p} \right),
  \qquad
  \tilde D_k = \frac{\nu_{\mathrm{t}}}{r^2} \frac{\mathrm{d}}{\mathrm{d} r}\left( r^2 \frac{\mathrm{d} \overline{\tilde k}}{\mathrm{d} r} \right),
\end{equation}
with $\tilde u_r = n_j \tilde u_j$. The budget we evaluate is
\begin{equation}\label{eq:ns-pp-k-radial-ctau}
  0 = c_\tau \beta_L\, r \frac{\mathrm{d} \overline{\tilde k}}{\mathrm{d} r}
  - 2 c_\tau \beta_U\, \overline{\tilde k}
  + \tilde T_k + \tilde \Pi_k + \tilde D_k - \overline{\epsilon} ,
\end{equation}
each term a function of $r$ alone. It is the sum of the radial equation and
twice the tangential equation of \eqref{eq:ns-pp-rs-radial-ctau}, halved,
with $\overline{\tilde k} = \tfrac12 (R_r + 2 R_\perp)$ and
$\overline{\epsilon} = \tfrac12 (\varepsilon_r + 2\varepsilon_\perp)$; it is
not an independent check of the profiles, but it is the budget that the
$k$--$\epsilon$ model of \S\ref{sssec:ns-pp-eps-model} closes. Regularity at
the origin requires $\mathrm{d}\overline{\tilde k}/\mathrm{d}r = 0$ at $r = 0$.

\subsubsection{The dimensional budget at the origin}\label{sssec:ns-pp-k-origin}
We now apply the isotropy assumption directly to the dimensional budget
\eqref{eq:ns-pp-k-rate}, without passing to the self-similar frame, and
evaluate it at the origin. Let $r^* = |\mathbf{x}^*|$ and
$n_j = x_j^* / r^*$. The statistics $\overline{\tilde k^*}$ and
$\overline{\epsilon^*}$ depend only on $r^*$ and $t^*$, and the mean fluxes are
radial,
\begin{equation}\label{eq:ns-pp-k-origin-flux}
  \overline{\tilde u_j^* \tilde k^*} = n_j\, F_k^*(r^*, t^*) ,
  \qquad
  \overline{\tilde u_j^* \tilde p^*} = n_j\, F_p^*(r^*, t^*) ,
\end{equation}
with $F_k^* = \overline{\tilde u_r^* \tilde k^*}$,
$F_p^* = \overline{\tilde u_r^* \tilde p^*}$ and
$\tilde u_r^* = n_j \tilde u_j^*$. As in \eqref{eq:ns-pp-k-radial-terms}, the
divergence of a radial flux and the Laplacian of a radial function give
\begin{equation}\label{eq:ns-pp-k-origin-radial}
  \frac{\partial \overline{\tilde k^*}}{\partial t^*}
  = -\frac{1}{r^{*2}} \frac{\partial}{\partial r^*}\left( r^{*2} F_k^* \right)
  - \frac{1}{\rho}\, \frac{1}{r^{*2}} \frac{\partial}{\partial r^*}\left( r^{*2} F_p^* \right)
  + \frac{\nu_{\mathrm{t}}^*}{r^{*2}} \frac{\partial}{\partial r^*}\left( r^{*2} \frac{\partial \overline{\tilde k^*}}{\partial r^*} \right)
  - \overline{\epsilon^*} .
\end{equation}

\paragraph{Behaviour near the origin.}
Each term of \eqref{eq:ns-pp-k-origin-radial} has the form $0/0$ at
$r^* = 0$, and its value there follows from the regularity of the mean fields.
A smooth function of $r^*$ alone is even in $r^*$, so
\begin{equation}\label{eq:ns-pp-k-origin-k-expansion}
  \overline{\tilde k^*} = k_0^*(t^*) + k_2^*(t^*)\, r^{*2} + O(r^{*4}) .
\end{equation}
A smooth radial vector field $n_j F^*(r^*) = x_j^*\, F^*/r^*$ requires $F^*/r^*$
to be smooth, so $F^*$ is odd in $r^*$ and vanishes at the origin,
\begin{equation}\label{eq:ns-pp-k-origin-flux-expansion}
  F^* = a^*(t^*)\, r^* + O(r^{*3}) ,
  \qquad
  a^* = \left.\frac{\partial F^*}{\partial r^*}\right|_{r^*=0} .
\end{equation}
The flux itself vanishes at the origin, since there is no preferred direction
there, but its divergence does not:
\begin{equation}\label{eq:ns-pp-k-origin-limits}
  \lim_{r^* \to 0} \frac{1}{r^{*2}} \frac{\partial}{\partial r^*}\left( r^{*2} F^* \right)
  = 3\, a^* ,
  \qquad
  \lim_{r^* \to 0} \frac{1}{r^{*2}} \frac{\partial}{\partial r^*}\left( r^{*2} \frac{\partial \overline{\tilde k^*}}{\partial r^*} \right)
  = 6\, k_2^*
  = 3 \left.\frac{\partial^2 \overline{\tilde k^*}}{\partial r^{*2}}\right|_{r^*=0} .
\end{equation}
Equivalently, $n_j F^* = a^* x_j^* + O(r^{*3})$ has divergence $3 a^*$ at the
origin, and $\nabla^{*2} (k_2^* r^{*2}) = 6 k_2^*$.

\paragraph{Budget at the origin.}
Substituting \eqref{eq:ns-pp-k-origin-limits} into
\eqref{eq:ns-pp-k-origin-radial},
\begin{equation}\label{eq:ns-pp-k-origin}
  \frac{\mathrm{d} \overline{\tilde k^*}(0, t^*)}{\mathrm{d} t^*}
  = \underbrace{-3 \left.\frac{\partial F_k^*}{\partial r^*}\right|_{0}}_{\tilde T_k^*(0)}
  \;\underbrace{-\,\frac{3}{\rho} \left.\frac{\partial F_p^*}{\partial r^*}\right|_{0}}_{\tilde \Pi_k^*(0)}
  \;+\; \underbrace{3\, \nu_{\mathrm{t}}^* \left.\frac{\partial^2 \overline{\tilde k^*}}{\partial r^{*2}}\right|_{0}}_{\tilde D_k^*(0)}
  \;-\; \overline{\epsilon^*}(0, t^*) .
\end{equation}
The transport terms do not vanish at the origin in general. They are
divergences, so they integrate to zero over a volume on whose boundary the
fluxes vanish, but locally they exchange energy between the centre and its
surroundings. If the turbulent flux points outward near the origin,
$\partial F_k^* / \partial r^* > 0$ there and $\tilde T_k^*(0) < 0$: turbulent
transport drains energy from the centre. The same limits apply
to the self-similar budget \eqref{eq:ns-pp-k-radial-ctau}, whose additional
advection term $c_\tau \beta_L\, r\, \mathrm{d}\overline{\tilde k}/\mathrm{d}r$
vanishes at $r = 0$, so that there
\begin{equation}\label{eq:ns-pp-k-origin-ctau}
  0 = - 2 c_\tau \beta_U\, \overline{\tilde k}(0)
  + \tilde T_k(0) + \tilde \Pi_k(0) + \tilde D_k(0) - \overline{\epsilon}(0) ,
\end{equation}
with $\tilde T_k(0)$, $\tilde \Pi_k(0)$ and $\tilde D_k(0)$ given by the
nondimensional counterparts of \eqref{eq:ns-pp-k-origin}. When the profiles
are evaluated numerically, these limits replace the radial forms
\eqref{eq:ns-pp-k-radial-terms} at $r = 0$.

\subsection{\texorpdfstring{$\epsilon$}{epsilon} budget}\label{ssec:ns-pp-eps-budget}
This subsection derives the evolution equation of the mean dissipation
$\overline{\epsilon}$. In the stationary state it reduces to a balance
\eqref{eq:ns-pp-eps-radial} between radial profiles that are each measured
independently, so its residual tests the profiles; its dimensional form
\eqref{eq:ns-pp-eps-rate} is the starting point of \S\ref{sssec:ns-pp-eps-model}.
We proceed in three steps:
\begin{enumerate}
  \item write the dimensional budget by adapting the dissipation equation of
  \cite{mansour1987reynolds}, \eqref{eq:ns-pp-eps-rate};
  \item nondimensionalize it into the self-similar frame,
  \eqref{eq:ns-pp-eps-rate-constant};
  \item evaluate each term as a radial profile and form the stationary
  balance, \eqref{eq:ns-pp-eps-radial}.
\end{enumerate}

\subsubsection{Exact budget}\label{sssec:ns-pp-eps-exact}
\paragraph{Assumptions.}
Throughout, $1/Re = 0$, and the subgrid model is the self-similar
constant eddy viscosity \eqref{eq:ns-constant-viscosity}; the budget is not
formed under the Smagorinsky model. As in \S\ref{ssec:ns-pp-k-budget}, the ensemble-mean
velocity vanishes, $\overline{\tilde u_i} = 0$, so $\tilde u_i$ is the
fluctuating velocity. We work inside the sphere $r \le r_0$, where $W = 0$, so the
sponge does not enter the budget.
The quantity whose budget we form is half the trace of the dissipation tensor of
\S\ref{sssec:ns-pp-quantities}, which by \eqref{eq:ns-pp-diss-trace-les} and
\eqref{eq:ns-pp-diss-nondim} with $1/Re = 0$ is
\begin{equation}\label{eq:ns-pp-eps-def}
  \epsilon^* = \nu_{\mathrm{t}}^*\, \frac{\partial \tilde u_i^*}{\partial x_m^*} \frac{\partial \tilde u_i^*}{\partial x_m^*} ,
  \qquad
  \epsilon = \nu_{\mathrm{t}}\, \frac{\partial \tilde u_i}{\partial x_m} \frac{\partial \tilde u_i}{\partial x_m} ,
  \qquad
  \epsilon^* = \frac{U^3}{L}\, \epsilon .
\end{equation}

\paragraph{Dimensional budget.}
Under the self-similar constant eddy viscosity,
$\nu_{\mathrm{t}} = C_\nu \Delta^{4/3}$ with $\Delta = \Delta x$ fixed in the
self-similar frame. Since $\Delta^* = L\, \Delta$, the dimensional eddy
viscosity is
\begin{equation}\label{eq:ns-pp-eps-nut}
  \nu_{\mathrm{t}}^* = C_\nu \Delta^{4/3}\, U L = UL\, \nu_{\mathrm{t}} .
\end{equation}
It is uniform in space but not in time. To account for this, differentiate
\eqref{eq:ns-pp-eps-def} by the product rule,
\begin{equation*}
  \frac{\partial \overline{\epsilon^*}}{\partial t^*}
  = \frac{1}{\nu_{\mathrm{t}}^*}\frac{\mathrm{d}\nu_{\mathrm{t}}^*}{\mathrm{d}t^*}\, \overline{\epsilon^*}
  + \nu_{\mathrm{t}}^*\, \frac{\partial}{\partial t^*}\, \overline{\frac{\partial \tilde u_i^*}{\partial x_m^*} \frac{\partial \tilde u_i^*}{\partial x_m^*}} .
\end{equation*}
The first term comes from the change of the eddy viscosity alone. In the
second, the time derivative of the mean squared velocity gradient involves
the viscosity only as the
instantaneous coefficient of the viscous term, so the derivation of
\cite{mansour1987reynolds} applies unchanged: it gives their dissipation
equation for zero mean velocity, with $\nu^*$ replaced by $\nu_{\mathrm{t}}^*$. The
dimensional budget is therefore
\begin{equation}\label{eq:ns-pp-eps-rate}
  \frac{\partial \overline{\epsilon^*}}{\partial t^*}
  = \frac{1}{\nu_{\mathrm{t}}^*}\frac{\mathrm{d}\nu_{\mathrm{t}}^*}{\mathrm{d}t^*}\, \overline{\epsilon^*}
  + \tilde P_\epsilon^* + \tilde T_\epsilon^* + \tilde \Pi_\epsilon^* + \tilde D_\epsilon^* - \tilde \Upsilon^*,
\end{equation}
with
\begin{equation}\label{eq:ns-pp-eps-terms}
\begin{aligned}
  \tilde P_\epsilon^* &= -2\nu_{\mathrm{t}}^*\, \overline{\frac{\partial \tilde u_i^*}{\partial x_k^*} \frac{\partial \tilde u_i^*}{\partial x_m^*} \frac{\partial \tilde u_k^*}{\partial x_m^*}}
  && \text{turbulent production,} \\
  \tilde T_\epsilon^* &= -\frac{\partial}{\partial x_k^*}\, \overline{\tilde u_k^* \epsilon^*}
  && \text{turbulent transport,} \\
  \tilde \Pi_\epsilon^* &= -\frac{2\nu_{\mathrm{t}}^*}{\rho}\, \frac{\partial}{\partial x_k^*}\, \overline{\frac{\partial \tilde p^*}{\partial x_m^*} \frac{\partial \tilde u_k^*}{\partial x_m^*}}
  && \text{pressure transport,} \\
  \tilde D_\epsilon^* &= \nu_{\mathrm{t}}^* \frac{\partial^2 \overline{\epsilon^*}}{\partial x_k^* \partial x_k^*}
  && \text{viscous diffusion,} \\
  \tilde \Upsilon^* &= 2 \left(\nu_{\mathrm{t}}^*\right)^2 \overline{\frac{\partial^2 \tilde u_i^*}{\partial x_k^* \partial x_m^*} \frac{\partial^2 \tilde u_i^*}{\partial x_k^* \partial x_m^*}}
  && \text{dissipation of } \epsilon^* .
\end{aligned}
\end{equation}
The tilde marks terms computed from the filtered velocity and pressure with
the eddy viscosity. $\tilde P_\epsilon^*$ is the turbulent production
$P_\epsilon^4$ of \cite{mansour1987reynolds}; their mean-advection term and
the productions $P_\epsilon^1$, $P_\epsilon^2$, and $P_\epsilon^3$ involve the
mean velocity and vanish. Compared with the equation of
\cite{mansour1987reynolds}, the eddy-viscosity change is the only term the LES
adds.

\paragraph{Nondimensionalization.}
We now move \eqref{eq:ns-pp-eps-rate} into the self-similar frame, where the
profiles are measured. The substitutions are
$\tilde u_i^* = U \tilde u_i$, $x_j^* = L x_j$, $\tilde p^* = \rho U^2 \tilde p$,
$\nu_{\mathrm{t}}^* = UL\, \nu_{\mathrm{t}}$ \eqref{eq:ns-pp-eps-nut}, and
$\epsilon^* = (U^3/L)\, \epsilon$ \eqref{eq:ns-pp-eps-def}.
\begin{itemize}
  \item \emph{Right-hand side.} Each of the five terms carries exactly the
  factor $U^4/L^2$: each velocity gradient contributes $U/L$, each further
  derivative $1/L$, $\nu_{\mathrm{t}}^*$ contributes $UL$, and the pressure
  $\rho U^2$. For example,
  $\tilde P_\epsilon^* = (UL)(U/L)^3\, \tilde P_\epsilon = (U^4/L^2)\, \tilde P_\epsilon$.
  \item \emph{Eddy-viscosity change.} With $\mathrm{d}t^*/\mathrm{d}t = \tau$
  and the rates \eqref{eq:fw-rates},
  \begin{equation}\label{eq:ns-pp-eps-nut-rate}
    \frac{1}{\nu_{\mathrm{t}}^*}\frac{\mathrm{d}\nu_{\mathrm{t}}^*}{\mathrm{d}t^*}
    = \frac{\sigma_U + \sigma_L}{\tau} ,
  \end{equation}
  and with the clock $\tau = L/U$ \eqref{eq:ns-clock} the term becomes
  $(U^4/L^2)(\sigma_U + \sigma_L)\, \overline{\epsilon}$.
  \item \emph{Left-hand side.} The chain rule \eqref{eq:fw-lhs-chain}, with
  the amplitude $U^3/L$ (rate $3\sigma_U - \sigma_L$) in place of $U$, gives
  \begin{equation*}
    \frac{\partial \overline{\epsilon^*}}{\partial t^*}
    = \frac{U^4}{L^2}\left( \frac{\partial \overline{\epsilon}}{\partial t}
    + \left( 3\sigma_U - \sigma_L \right) \overline{\epsilon}
    - \sigma_L\, x_j \frac{\partial \overline{\epsilon}}{\partial x_j} \right).
  \end{equation*}
\end{itemize}
Dividing by $U^4/L^2$ and collecting the terms in $\overline{\epsilon}$,
$(\sigma_U + \sigma_L) - (3\sigma_U - \sigma_L) = -2(\sigma_U - \sigma_L)$, gives
\begin{equation}\label{eq:ns-pp-eps-rate-constant}
  \frac{\partial \overline{\epsilon}}{\partial t}
  = \sigma_L\, x_j \frac{\partial \overline{\epsilon}}{\partial x_j}
  - 2\left( \sigma_U - \sigma_L \right) \overline{\epsilon}
  + \tilde P_\epsilon + \tilde T_\epsilon + \tilde \Pi_\epsilon + \tilde D_\epsilon - \tilde \Upsilon ,
\end{equation}
with the nondimensional terms
\begin{equation}\label{eq:ns-pp-eps-terms-nondim}
\begin{aligned}
  \tilde P_\epsilon &= -2\nu_{\mathrm{t}}\, \overline{\frac{\partial \tilde u_i}{\partial x_k} \frac{\partial \tilde u_i}{\partial x_m} \frac{\partial \tilde u_k}{\partial x_m}},
  &
  \tilde T_\epsilon &= -\frac{\partial}{\partial x_k}\, \overline{\tilde u_k \epsilon},
  &
  \tilde \Pi_\epsilon &= -2\nu_{\mathrm{t}}\, \frac{\partial}{\partial x_k}\, \overline{\frac{\partial \tilde p}{\partial x_m} \frac{\partial \tilde u_k}{\partial x_m}},
  \\
  \tilde D_\epsilon &= \nu_{\mathrm{t}} \frac{\partial^2 \overline{\epsilon}}{\partial x_k \partial x_k},
  &
  \tilde \Upsilon &= 2 \nu_{\mathrm{t}}^2\, \overline{\frac{\partial^2 \tilde u_i}{\partial x_k \partial x_m} \frac{\partial^2 \tilde u_i}{\partial x_k \partial x_m}} ,
\end{aligned}
\end{equation}
of the same form as \eqref{eq:ns-pp-eps-terms} without stars. The two new
terms on the right of \eqref{eq:ns-pp-eps-rate-constant} are the scaling terms
of the budget. The first, $\sigma_L\, x_j\, \partial_j \overline{\epsilon}$,
advects $\overline{\epsilon}$ radially as the frame dilates. The second
combines the rescaling of $\epsilon$ with the eddy-viscosity change: since
$\nu_{\mathrm{t}}$ is constant in the self-similar frame, $\epsilon$ scales as
the square of the gradient scale $U/L$, which changes at the rate
$\sigma_U - \sigma_L$.
\par
In the stationary state $\partial \overline{\epsilon}/\partial t = 0$ and the
rates are constant, $\sigma_X = c_\tau \beta_X$ (\S\ref{sssec:ns-pp-rates}),
so by \eqref{eq:ns-pp-eps-rate-constant} the profiles satisfy
\begin{equation}\label{eq:ns-pp-eps-radial}
  0 = \sigma_L\, r \frac{\mathrm{d} \overline{\epsilon}}{\mathrm{d} r}
  - 2\left( \sigma_U - \sigma_L \right) \overline{\epsilon}
  + \tilde P_\epsilon + \tilde T_\epsilon + \tilde \Pi_\epsilon + \tilde D_\epsilon - \tilde \Upsilon ,
\end{equation}
each term a function of $r$.
\par
The second scaling term depends on the rates only through their difference
$\sigma_U - \sigma_L$, which by \eqref{eq:ns-pp-rates-diff} is $-c_\tau$.
Substituting it, together with $\sigma_L = c_\tau \beta_L$, into
\eqref{eq:ns-pp-eps-radial} gives the budget we evaluate,
\begin{equation}\label{eq:ns-pp-eps-radial-ctau}
  0 = c_\tau \beta_L\, r \frac{\mathrm{d} \overline{\epsilon}}{\mathrm{d} r}
  + 2 c_\tau\, \overline{\epsilon}
  + \tilde P_\epsilon + \tilde T_\epsilon + \tilde \Pi_\epsilon + \tilde D_\epsilon - \tilde \Upsilon .
\end{equation}
The coefficients $c_\tau \beta_L$ and $2 c_\tau$ are constants, and the
profiles $\overline{\epsilon}$, $\tilde P_\epsilon$, $\tilde T_\epsilon$,
$\tilde \Pi_\epsilon$, $\tilde D_\epsilon$, and $\tilde \Upsilon$ are
statistics of the stationary state, functions of $r$ alone; every term of
\eqref{eq:ns-pp-eps-radial-ctau} is therefore independent of $t$.
Since every term is estimated independently, the residual of
\eqref{eq:ns-pp-eps-radial-ctau} measures the statistical and discretization
error of the profiles.
\par
\subsubsection{Modeling terms}\label{sssec:ns-pp-eps-model}
The budget \eqref{eq:ns-pp-eps-rate} is exact, but its production,
dissipation, and transport terms are unclosed. We approximate them in terms
of $\overline{k^*}$ and $\overline{\epsilon^*}$, following the standard
$k$--$\epsilon$ model \cite{pope2000turbulent}. The eddy-viscosity change and
the viscous diffusion $\tilde D_\epsilon^*$ are already closed and are kept
exactly.
\paragraph{Dimensional model.}
With zero mean velocity, the production $\tilde P_\epsilon^*$ and the
dissipation $\tilde \Upsilon^*$ are the two large terms of the budget and
nearly balance. We model their difference as
\begin{equation}\label{eq:ns-pp-eps-model-PY}
  \tilde P_\epsilon^* - \tilde \Upsilon^*
  \approx -C_{\epsilon 2}\, \frac{\left(\overline{\epsilon^*}\right)^2}{\overline{k^*}} ,
\end{equation}
and the turbulent and pressure transport as gradient diffusion,
\begin{equation}\label{eq:ns-pp-eps-model-T}
  \tilde T_\epsilon^* + \tilde \Pi_\epsilon^*
  \approx \frac{\partial}{\partial x_j^*}\left( \frac{\nu_\epsilon^*}{\sigma_\epsilon}\,
  \frac{\partial \overline{\epsilon^*}}{\partial x_j^*} \right),
  \qquad
  \nu_\epsilon^* = C_\mu \frac{\left(\overline{k^*}\right)^2}{\overline{\epsilon^*}} ,
\end{equation}
where $C_{\epsilon 2}$, $C_\mu$, and $\sigma_\epsilon$ are constants and
$\nu_\epsilon^*$ is the turbulent viscosity of the $k$--$\epsilon$ model,
distinct from the LES eddy viscosity $\nu_{\mathrm{t}}^*$. Since our
$\epsilon^*$ is the sink of $\overline{k^*}$, as in \cite{pope2000turbulent},
the constants are directly comparable with the standard ones.
Since $\nu_{\mathrm{t}}^*$ is uniform in space,
$\tilde D_\epsilon^* = \partial_j^* (\nu_{\mathrm{t}}^*\, \partial_j^* \overline{\epsilon^*})$
combines with \eqref{eq:ns-pp-eps-model-T} into a single diffusion term.
Substituting \eqref{eq:ns-pp-eps-model-PY} and \eqref{eq:ns-pp-eps-model-T}
into \eqref{eq:ns-pp-eps-rate} gives the modeled dimensional budget
\begin{equation}\label{eq:ns-pp-eps-model-dim}
  \frac{\partial \overline{\epsilon^*}}{\partial t^*}
  = \frac{1}{\nu_{\mathrm{t}}^*}\frac{\mathrm{d}\nu_{\mathrm{t}}^*}{\mathrm{d}t^*}\, \overline{\epsilon^*}
  + \frac{\partial}{\partial x_j^*}\left[ \left( \nu_{\mathrm{t}}^* + \frac{\nu_\epsilon^*}{\sigma_\epsilon} \right)
  \frac{\partial \overline{\epsilon^*}}{\partial x_j^*} \right]
  - C_{\epsilon 2}\, \frac{\left(\overline{\epsilon^*}\right)^2}{\overline{k^*}} .
\end{equation}
\paragraph{Nondimensionalization.}
The modeled terms scale like the exact terms they replace. With
$\overline{k^*} = U^2\, \overline{k}$ and
$\overline{\epsilon^*} = (U^3/L)\, \overline{\epsilon}$ \eqref{eq:ns-pp-k-scales},
\begin{equation}\label{eq:ns-pp-eps-model-scales}
  \frac{\left(\overline{\epsilon^*}\right)^2}{\overline{k^*}}
  = \frac{U^4}{L^2}\, \frac{\overline{\epsilon}^{\,2}}{\overline{k}} ,
  \qquad
  \nu_\epsilon^* = UL\, \nu_\epsilon ,
  \quad
  \nu_\epsilon = C_\mu \frac{\overline{k}^{\,2}}{\overline{\epsilon}} ,
\end{equation}
so $\nu_\epsilon^*$ scales like $\nu_{\mathrm{t}}^*$ \eqref{eq:ns-pp-eps-nut},
and the diffusion term, with one factor $U^3/L$ from $\overline{\epsilon^*}$
and two derivatives $1/L$, carries $(UL)(U^3/L)/L^2 = U^4/L^2$. Every term of
\eqref{eq:ns-pp-eps-model-dim} therefore carries the factor $U^4/L^2$, as in
\eqref{eq:ns-pp-eps-rate-constant}, and the left-hand side and the
eddy-viscosity change transform as there. The modeled nondimensional budget
is
\begin{equation}\label{eq:ns-pp-eps-model-nondim}
  \frac{\partial \overline{\epsilon}}{\partial t}
  = \sigma_L\, x_j \frac{\partial \overline{\epsilon}}{\partial x_j}
  - 2\left( \sigma_U - \sigma_L \right) \overline{\epsilon}
  + \frac{\partial}{\partial x_j}\left[ \left( \nu_{\mathrm{t}} + \frac{\nu_\epsilon}{\sigma_\epsilon} \right)
  \frac{\partial \overline{\epsilon}}{\partial x_j} \right]
  - C_{\epsilon 2}\, \frac{\overline{\epsilon}^{\,2}}{\overline{k}} ,
\end{equation}
with $\nu_{\mathrm{t}} = C_\nu \Delta^{4/3}$ \eqref{eq:ns-constant-viscosity}.
\paragraph{Isotropic form.}
Under the isotropy assumption of \S\ref{sssec:ns-pp-scaling},
$\overline{k}$ and $\overline{\epsilon}$, and hence $\nu_\epsilon$, depend on
space only through $r = \sqrt{x_i x_i}$. With $\partial r/\partial x_j = x_j/r$,
the gradient of a radial function $f$ is
$\partial f/\partial x_j = (x_j/r)\, \mathrm{d}f/\mathrm{d}r$, so the
advection term becomes
$x_j\, \partial \overline{\epsilon}/\partial x_j = r\, \partial \overline{\epsilon}/\partial r$,
and the divergence of the radial flux
$F_j = (x_j/r)\, F(r)$ is $r^{-2}\, \mathrm{d}(r^2 F)/\mathrm{d}r$. The
modeled budget \eqref{eq:ns-pp-eps-model-nondim} then becomes
\begin{equation}\label{eq:ns-pp-eps-model-radial}
  \frac{\partial \overline{\epsilon}}{\partial t}
  = \sigma_L\, r \frac{\partial \overline{\epsilon}}{\partial r}
  - 2\left( \sigma_U - \sigma_L \right) \overline{\epsilon}
  + \frac{1}{r^2} \frac{\partial}{\partial r}\left[ r^2 \left( \nu_{\mathrm{t}} + \frac{\nu_\epsilon}{\sigma_\epsilon} \right)
  \frac{\partial \overline{\epsilon}}{\partial r} \right]
  - C_{\epsilon 2}\, \frac{\overline{\epsilon}^{\,2}}{\overline{k}} ,
  \qquad
  \nu_\epsilon = C_\mu \frac{\overline{k}^{\,2}}{\overline{\epsilon}} ,
\end{equation}
a one-dimensional equation for $\overline{\epsilon}(r, t)$. In the stationary
state $\partial \overline{\epsilon}/\partial t = 0$ and the rates are constant,
$\sigma_X = c_\tau \beta_X$ (\S\ref{sssec:ns-pp-rates}), so
\eqref{eq:ns-pp-eps-model-radial} reduces to the ordinary differential
equation
\begin{equation}\label{eq:ns-pp-eps-model-stationary}
  0 = \sigma_L\, r \frac{\mathrm{d} \overline{\epsilon}}{\mathrm{d} r}
  - 2\left( \sigma_U - \sigma_L \right) \overline{\epsilon}
  + \frac{1}{r^2} \frac{\mathrm{d}}{\mathrm{d} r}\left[ r^2 \left( \nu_{\mathrm{t}} + \frac{\nu_\epsilon}{\sigma_\epsilon} \right)
  \frac{\mathrm{d} \overline{\epsilon}}{\mathrm{d} r} \right]
  - C_{\epsilon 2}\, \frac{\overline{\epsilon}^{\,2}}{\overline{k}} .
\end{equation}
Regularity at the origin requires
$\mathrm{d}\overline{\epsilon}/\mathrm{d}r = 0$ at $r = 0$.

\appendix


\section{Navier--Stokes numerical details}\label{app:ns-numerics}

This appendix collects two implementation details of the Navier--Stokes
solver of \S\ref{sec:ns}: the construction of the initial condition and the
evaluation of the scaling operator by spectral interpolation in three
dimensions.

\subsection{Initial condition}\label{ssec:ns-ic}
Because the curl of a vector feild is divergence-free we will define $\tilde u_i$ as the curl of some feild $A_i$
\begin{equation}\label{eq:ns-ic-curl}
  \tilde u_i = \varepsilon_{ijk} \frac{\partial A_k}{\partial x_j},
  \qquad\text{equivalently}\qquad
  \hat{\tilde u}_i = \mathrm{i}\, \varepsilon_{ijk}\, \kappa_j \hat{A}_k.
\end{equation}
$\tilde u_i$ is built on the grid \eqref{eq:alg-grid} in four steps, followed
by a rescaling.
\begin{enumerate}
  \item \emph{Noise.} Each component $A_k(\mathbf{x_n})$ is an independent
  standard normal, drawn from a seeded Mersenne Twister.  
  \item \emph{Spectral filter.} The noise is given a spectrum,
  \begin{equation}\label{eq:ns-ic-filter}
    \hat{A}_k \gets \hat{A}_k\, e^{-\kappa^2/\kappa_p^2}.
  \end{equation}
  \item \emph{Window.} The potential is localized in physical space,
  \begin{equation}\label{eq:ns-ic-window}
    A_k \gets A_k\, e^{-r^2/2\sigma^2},
  \end{equation}
  with $r$ the same distance from the centre of the domain that $E_2$ uses in
  \eqref{eq:ns-functionals}.  
  \item \emph{Rescaling.} Finally $\tilde u_i \gets \tilde u_i \sqrt{E_0 / F_1[\tilde u]}$,
  with $F_1$ evaluated as the Riemann sum \eqref{eq:alg-riemann}, so that the
  field delivered to the integrator carries a prescribed $E_0$.
\end{enumerate}
\subsection{Spectral interpolation in three dimensions}\label{ssec:ns-interpolation}

The scaling operator \eqref{eq:fw-scale-step} needs the same band-limited
evaluation \eqref{eq:fw-spectral-interp} as in one dimension, but at $d = 3$ a
direct evaluation of it is no longer affordable, and its realization on the
real transform differs enough from \eqref{eq:1d-interp} to be worth writing
out.

\paragraph{The sum on the real transform.}
Write $\hat{\tilde u}_i^{(k)}(\mathbf{j}) = N^{-3}\,\mathrm{RFFT}(\tilde u_k)_{\mathbf{j}}$
for the normalized coefficients, with the first axis carrying only the
non-negative frequencies $j_1 = 0, 1, \dots, N/2$ and the remaining two
carrying the full signed set $j_d = 0, 1, \dots, N/2, -(N/2-1), \dots, -1$.
Two corrections turn \eqref{eq:fw-spectral-interp} into a sum over that
layout. First, the omitted Hermitian conjugates are restored by the
multiplicity
\begin{equation}\label{eq:ns-interp-multiplicity}
  m_{j_1} =
  \begin{cases}
    1, & j_1 = 0, \ \text{or } j_1 = N/2 \text{ with } N \text{ even}, \\
    2, & \text{otherwise},
  \end{cases}
\end{equation}
exactly as the factor of $2$ does in \eqref{eq:1d-interp}. Second, the
transform phases its coefficients to the sample index $n$, which on the
centred grid \eqref{eq:alg-grid} places its origin at $x = -H/2$ rather than
at $x = 0$. Since the scaling operator dilates about the centre of the domain,
each mode must be evaluated at $\lambda_L x + H/2$, not at
$\lambda_L (x + H/2)$ --- the two agree only at $\lambda_L = 1$. As $H/2$ is
half a period, the correction is the constant phase
$\exp(\mathrm{i} \kappa_j H/2) = (-1)^{j}$ per axis. Together,
\begin{equation}\label{eq:ns-interp}
  \Phi_{\Delta t}^{(\mathrm{S})}(\tilde u_k)_i(\mathbf{x})
  = \frac{1}{\lambda_{U,k+1}} \operatorname{Re}
    \sum_{j_1 = 0}^{N/2} \sum_{j_2} \sum_{j_3}
      m_{j_1}\, \hat{\tilde u}_i^{(k)}(\mathbf{j})
      \prod_{d=1}^{3} (-1)^{j_d}
      \exp\!\left( \frac{2\pi \mathrm{i}\, j_d\, \lambda_{L,k+1}\, x_d}{H} \right).
\end{equation}

\paragraph{Separable evaluation.}
Evaluated as written, \eqref{eq:ns-interp} costs $O(N^{2d})$: every retained
mode is summed at every grid point, which at $N = 128$ is $4 \times 10^{12}$
complex multiply--adds per scaling substep. The kernel, however, is a product
of one factor per axis, so the triple sum separates and can be applied one
axis at a time. Define the two interpolation matrices
\begin{equation}\label{eq:ns-interp-matrices}
  W^{(1)}_{n j_1} = m_{j_1} (-1)^{j_1}
    \exp\!\left( \frac{2\pi \mathrm{i}\, j_1 \lambda_{L,k+1} x_n}{H} \right),
  \qquad
  W^{(\mathrm{f})}_{n j} = (-1)^{j}
    \exp\!\left( \frac{2\pi \mathrm{i}\, j\, \lambda_{L,k+1} x_n}{H} \right),
\end{equation}
the first of size $N \times (N/2 + 1)$ for the half-spectrum axis and the
second of size $N \times N$ for the two full axes, both rebuilt whenever
$\lambda_{L,k+1}$ changes. Then \eqref{eq:ns-interp} is the sequence of
contractions
\begin{equation}\label{eq:ns-interp-contractions}
  \hat{\tilde u}_i \;\xrightarrow{\ W^{(\mathrm{f})} \text{ on axis } 3\ }\;
  \;\xrightarrow{\ W^{(\mathrm{f})} \text{ on axis } 2\ }\;
  \;\xrightarrow{\ W^{(1)} \text{ on axis } 1\ }\;
  \operatorname{Re}(\cdot),
\end{equation}
each of which is a matrix--matrix product once the contracted axis is brought
to the front. The cost falls to $O(d\, N^{d+1})$, i.e.\ $8 \times 10^{8}$
multiply--adds at $N = 128$, a factor of $5 \times 10^{3}$ below the direct
sum, and the work is expressed as three dense $\mathrm{GEMM}$ calls. The
factorization is exact --- it is a regrouping of the same finite sum, not an
approximation --- so the guarantee of \S\ref{ssec:fw-interpolation} that the
scaling step contributes no spatial discretization error is retained. For
$d = 1$ the sequence \eqref{eq:ns-interp-contractions} collapses to the single
contraction with $W^{(1)}$, which is \eqref{eq:1d-interp}.

\paragraph{The Nyquist bin.}
For even $N$ the bins $j_d = \pm N/2$ of a real field's transform coincide,
and the two axis layouts resolve that differently. On the half-spectrum axis
the mode is stored once and is given multiplicity $m_{N/2} = 1$ by
\eqref{eq:ns-interp-multiplicity}. On a full axis it is stored in the single
bin $j_d = N/2$ but represents the conjugate pair, so it must be split evenly
between $+N/2$ and $-N/2$; the two halves combine to
\begin{equation}\label{eq:ns-interp-nyquist}
  W^{(\mathrm{f})}_{n,\, N/2} = (-1)^{N/2}
    \cos\!\left( \frac{\pi N \lambda_{L,k+1} x_n}{H} \right),
\end{equation}
the same cosine that closes \eqref{eq:1d-interp}. With
\eqref{eq:ns-interp-nyquist} in place the contraction
\eqref{eq:ns-interp-contractions} reproduces the direct sum
\eqref{eq:ns-interp} to roundoff for either parity of $N$.

\paragraph{The out-of-box cutoff.}
The box indicator $\chi$ of \eqref{eq:fw-dilate-cutoff} is a product of one
factor per axis, like the kernel, so it separates the same way. It is imposed by
zeroing row $n$ of each interpolation matrix \eqref{eq:ns-interp-matrices}
whenever the query of that axis leaves the box,
\begin{equation}\label{eq:ns-interp-cutoff}
  W^{(1)}_{n j_1} \gets 0, \quad W^{(\mathrm{f})}_{n j} \gets 0
  \qquad \text{whenever } |x_n \lambda_{L,k+1}| > H/2 .
\end{equation}
After the three contractions \eqref{eq:ns-interp-contractions} the point
$\mathbf{x_n}$ then carries the product of the three row factors, which is
$\chi(\mathbf{x_n}\lambda_{L,k+1})$ exactly. The cutoff is thus applied at no
cost beyond building the matrices, and the direct sum it is tested against
applies $\chi$ pointwise. For $\lambda_{L,k+1}$ within a few percent of unity
the zeroed rows are the outermost one or two grid points per side. These lie
deep in the sponge layer (\S\ref{ssec:ns-sponge}), where the field has already
been damped. Zeroing a region in physical space is not a
solenoidal operation, so it reintroduces a divergence
proportional to the amplitude it removes.

\paragraph{Loss of solenoidality and the Leray projection.}
The continuous scaling operator preserves continuity: if
$\partial \tilde u_j/\partial x_j = 0$, then
$\partial/\partial x_j \bigl[\lambda_U^{-1} \tilde u_j(\mathbf{x} \lambda_L)\bigr]
= \lambda_U^{-1} \lambda_L\, (\partial \tilde u_j/\partial x_j)(\mathbf{x} \lambda_L) = 0$,
the counterpart in physical space of \eqref{eq:ns-continuity-dk}. Its discrete
evaluation does not preserve it exactly. The dilated field carries the
wavenumbers $\lambda_L \boldsymbol{\kappa}$, which lie off the lattice of the grid, so
sampling \eqref{eq:ns-interp} at the grid points and transforming back
redistributes each mode over neighbouring lattice wavenumbers without regard
to $\kappa_i \hat{\tilde u}_i = 0$; the out-of-box cutoff \eqref{eq:ns-interp-cutoff}
further removes amplitude in physical space. Both leave a divergence on the
grid. When the option $\mathit{project}$ of \S\ref{ssec:alg-substeps} is
enabled, \textsc{Scale} therefore ends with the Leray projection of the
scaled field back onto the solenoidal subspace, using the projection tensor
\eqref{eq:ns-projection-tensor},
\begin{equation}\label{eq:ns-leray-scale}
  \hat{\tilde u}_i(\boldsymbol{\kappa}) \gets P_{il}(\boldsymbol{\kappa})\, \hat{\tilde u}_l(\boldsymbol{\kappa})
  = \hat{\tilde u}_i - \frac{\kappa_i \kappa_l \hat{\tilde u}_l}{\kappa^2},
  \qquad \boldsymbol{\kappa} \neq \mathbf{0},
\end{equation}
followed by the transform back to physical space. The mean mode
$\boldsymbol{\kappa} = \mathbf{0}$ is left unchanged. The projection is orthogonal, so
it can only remove energy, and the cost is one forward and one inverse
transform per scaling step. Under \textsc{Strang} it is applied after each of
the two scaling half steps, so the field entering \textsc{Evolve} is
solenoidal as well as the field the step returns.

\paragraph{Interaction with dealiasing.}
The dilation maps each represented wavenumber to $\lambda_{L,k+1} \boldsymbol{\kappa}$,
so a field that fills the spectrum up to the Nyquist wavenumber would wrap
under any $\lambda_{L,k+1} > 1$. The $2/3$ rule already required by the
pseudo-spectral product \eqref{eq:ns-fourier-convection} supplies the margin:
the top retained mode sits at $N/3$ and moves to $\lambda_{L,k+1} N/3$, which
stays below $N/2$ for every $\lambda_{L,k+1} < 3/2$. Runs in the self-similar
frame hold $\lambda_{L,k+1}$ within a few percent of unity, so the margin is
not approached, but it is the condition under which the interpolation stays
free of aliasing. Finally, \eqref{eq:ns-interp-contractions} is applied to
each velocity component independently and mixes neither components nor the
wavenumbers within an axis; at the exceptional value $\lambda_{L,k+1} = 1$ it
reduces to the plain inverse transform, which is taken instead at
$O(N^d \log N)$.


\section{Radial profiles from the spectral fields}\label{app:radial}

This appendix sets out how the radial profiles of
\S\ref{ssec:ns-postproc} are computed from a snapshot of the self-similar
fields: the spherical mean \eqref{eq:ns-pp-sphere-mean} at a prescribed set of
radii, the pressure that enters $\tilde \Pi_\epsilon$, the
refined grid on which the integrands are formed, and the sampling in
time.

\subsection{The spherical mean of a band-limited field}\label{ssec:radial-mean}
Let $Q(\mathbf{x_n})$ be a real field on the grid \eqref{eq:alg-grid}, with
discrete Fourier coefficients
\begin{equation}\label{eq:radial-dft}
  \hat Q_{\mathbf{m}} = \sum_{\mathbf{n}} Q(\mathbf{x_n})\,
  e^{-\mathrm{i} \boldsymbol{\kappa}_{\mathbf{m}} \cdot \mathbf{n}\, \Delta x},
  \qquad
  \boldsymbol{\kappa}_{\mathbf{m}} = \frac{2\pi}{H}\, \mathbf{m},
\end{equation}
over the signed indices $m_j \in \{-N/2, \dots, N/2 - 1\}$. The transform is
phased to the sample index, whose origin sits at $\mathbf{x} = -(H/2)(1,1,1)$
rather than at $\mathbf{x} = 0$, so its trigonometric interpolant is
\begin{equation}\label{eq:radial-interpolant}
  Q(\mathbf{x}) = \frac{1}{N^3} \sum_{\mathbf{m}} \hat s_{\mathbf{m}}\,
  e^{\mathrm{i} \boldsymbol{\kappa}_{\mathbf{m}} \cdot \mathbf{x}},
  \qquad
  \hat s_{\mathbf{m}} = (-1)^{m_1 + m_2 + m_3}\, \hat Q_{\mathbf{m}},
\end{equation}
with the phase-origin correction of \eqref{eq:1d-interp} in each direction,
since $\kappa_{m_j} H/2 = \pi m_j$. The spherical mean of a single mode is
\begin{equation}\label{eq:radial-mode-mean}
  \frac{1}{4\pi} \oint e^{\mathrm{i} \boldsymbol{\kappa} \cdot r \hat{\mathbf{n}}}\,
  \mathrm{d}\Omega
  = j_0(|\boldsymbol{\kappa}|\, r),
  \qquad
  j_0(z) = \frac{\sin z}{z}, \quad j_0(0) = 1,
\end{equation}
so the spherical mean \eqref{eq:ns-pp-sphere-mean} of the interpolant is
\begin{equation}\label{eq:radial-mean-modes}
  \langle Q \rangle_r = \frac{1}{N^3} \sum_{\mathbf{m}}
  \mathrm{Re}\bigl( \hat s_{\mathbf{m}} \bigr)\, j_0(|\boldsymbol{\kappa}_{\mathbf{m}}|\, r),
\end{equation}
the real part because $Q$ is real and $j_0$ is even. Evaluated with the real
transform, which stores the half spectrum $m_1 \ge 0$, each stored mode with
$0 < m_1 < N/2$ stands also for its omitted conjugate and carries weight $2$,
as in \eqref{eq:1d-interp}; the planes $m_1 = 0$ and $m_1 = N/2$ carry weight
$1$. At $r = 0$, \eqref{eq:radial-mean-modes} is the interpolant at the
origin, which is the grid value there.
\par
Equation \eqref{eq:radial-mean-modes} is exact for the interpolant; no
interpolation is done in physical space. For the resolved velocity, which is
band-limited to the retained $2/3$ band, the interpolant is the field itself.
For a product such as $\epsilon$, formed pointwise on the grid, it is the
interpolant of the product as sampled, the same spectral interpolation as
\S\ref{ssec:fw-interpolation}.
\par
The kernel depends on $\boldsymbol{\kappa}_{\mathbf{m}}$ only through the
integer $s = m_1^2 + m_2^2 + m_3^2 \le 3 (N/2)^2$. Summing the weighted
coefficients over each such value first,
\begin{equation}\label{eq:radial-group}
  g_s = \sum_{\mathbf{m} : |\mathbf{m}|^2 = s} c_{\mathbf{m}}\,
  \mathrm{Re}\bigl( \hat s_{\mathbf{m}} \bigr),
  \qquad
  \langle Q \rangle_r = \frac{1}{N^3} \sum_{s} g_s\,
  j_0\!\left( \frac{2\pi}{H} \sqrt{s}\, r \right),
\end{equation}
with $c_{\mathbf{m}} \in \{1, 2\}$ the weights above, turns the sum into a
regrouping of the spectrum, done once per field, followed by a product with
the matrix $J_{is} = j_0\bigl( (2\pi/H) \sqrt{s}\, r_i \bigr)$, tabulated once
per run for the radii $r_i$ below. The cost per field is one forward transform
and a matrix--vector product of size $n_r \times 3(N/2)^2$.
\par
The radial derivative of the profile follows from the same coefficients.
Since $j_0^\prime = -j_1$, with $j_1(z) = \sin z / z^2 - \cos z / z$,
\begin{equation}\label{eq:radial-derivative}
  r \frac{\partial \langle Q \rangle_r}{\partial r}
  = -\frac{1}{N^3} \sum_{s} g_s\, \kappa_s r\, j_1(\kappa_s r),
  \qquad \kappa_s = \frac{2\pi}{H} \sqrt{s},
\end{equation}
a second tabulated kernel applied to the same $g_s$. Because the spherical
mean commutes with $x_j\, \partial/\partial x_j = r\, \partial/\partial r$,
this is $\langle x_j\, \partial Q / \partial x_j \rangle_r$, the advection term
of \eqref{eq:ns-pp-eps-rate-constant}. It is evaluated this way rather than from
the pointwise product $x_j\, \partial Q/\partial x_j$, which is not periodic
and whose interpolant rings at the faces of the box. It vanishes at $r = 0$.

\subsection{Radial projections of tensor fields}\label{ssec:radial-tensor}
The Reynolds stress budget needs the spherical means of the projections
\eqref{eq:ns-pp-rs-projections} of symmetric tensor fields $X_{ij}$. The
weight $n_i n_j = x_i x_j / r^2$ is not smooth at the origin, so
$n_i n_j X_{ij}$ is not band-limited and the mean of its interpolant would
not be exact. Instead, the projection is moved onto the kernel. Differentiating
$\langle e^{\mathrm{i} \boldsymbol{\kappa} \cdot \mathbf{x}} \rangle_r = j_0(|\boldsymbol{\kappa}| r)$
twice with respect to $\boldsymbol{\kappa}$ gives
\begin{equation}\label{eq:radial-tensor}
  \frac{1}{4\pi} \oint n_i n_j\, e^{\mathrm{i} \boldsymbol{\kappa} \cdot r \hat{\mathbf{n}}}\,
  \mathrm{d}\Omega
  = \frac{j_1(z)}{z}\, \delta_{ij} - j_2(z)\, \hat\kappa_i \hat\kappa_j ,
  \qquad z = |\boldsymbol{\kappa}| r ,
\end{equation}
with $\hat{\boldsymbol{\kappa}} = \boldsymbol{\kappa}/|\boldsymbol{\kappa}|$ and
$j_2(z) = (3/z^2 - 1) \sin z / z - 3 \cos z / z^2$. Hence
$\langle n_i n_j X_{ij} \rangle_r$ is \eqref{eq:radial-group} applied to two
fields of the components, the trace $\hat X_{ii}$ with the kernel $j_1(z)/z$
and the contraction $\hat\kappa_i \hat\kappa_j \hat X_{ij}$ with the kernel
$-j_2(z)$. Both are formed in Fourier space from the band-limited components
and grouped by $s$ like any other field, so the projection is exact. The
tangential projection follows as
$X_\perp = \tfrac12 (\langle X_{ii} \rangle_r - X_r)$. The radial derivatives
use $z\, \mathrm{d}(j_1/z)/\mathrm{d}z = -j_2$ and
$z\, j_2^\prime = z j_1 - 3 j_2$, as in \eqref{eq:radial-derivative}. At
$r = 0$ the kernels are $1/3$ and $0$, so $X_r = X_\perp = X_{ii}/3$.

\subsection{The radii}\label{ssec:radial-radii}
The profiles are evaluated at $r = 0$ and at $n_r$ logarithmically spaced
radii,
\begin{equation}\label{eq:radial-radii}
  r_i = r_{\min} \left( \frac{r_{\max}}{r_{\min}} \right)^{(i-1)/(n_r - 1)},
  \qquad i = 1, \dots, n_r,
\end{equation}
together with the few radii at which the statistical error is reported,
which are added to the set exactly rather than taken from the nearest $r_i$.
By default $n_r = 32$, $r_{\min} = \Delta x$, and $r_{\max} = r_0$, the radius
inside which the sponge forcing vanishes and the $\epsilon$ budget of
\S\ref{ssec:ns-pp-eps-budget} applies.

\subsection{The pressure}\label{ssec:radial-pressure}
The pressure transport $\tilde \Pi_\epsilon$ needs the nondimensional pressure
itself, which the projection \eqref{eq:ns-projection-tensor} eliminates and
the solver never forms. It is recovered from the Poisson equation
\eqref{eq:ns-fourier-poisson} with the convective term in divergence form,
$\hat N_l = \mathrm{i} \kappa_j \widehat{\tilde u_l \tilde u_j}$, which equals
\eqref{eq:ns-fourier-convection} for a solenoidal field,
\begin{equation}\label{eq:radial-pressure}
  \hat{\tilde p} = -\frac{\kappa_l \kappa_j\, \widehat{\tilde u_l \tilde u_j}}{\kappa^2},
  \qquad \kappa \ne 0,
  \qquad \hat{\tilde p}(\mathbf{0}) = 0,
\end{equation}
without the sponge term $\mathrm{i} \kappa_l \hat F_l / \kappa^2$ of
\eqref{eq:ns-fourier-poisson}, since the sponge is not part of the budget
(\S\ref{ssec:ns-pp-eps-budget}). The products $\tilde u_l \tilde u_j$ are formed on the refined grid of
\S\ref{ssec:radial-integrands}, where they are exact. The solver itself forms the convective
term in rotational form, $-\tilde{\mathbf{u}} \times \boldsymbol{\omega}$,
which differs from $\tilde u_j\, \partial \tilde u_i / \partial x_j$ by the
gradient $\partial_i (\tfrac12 \tilde u_k \tilde u_k)$; the pressure it implies
is therefore the total head $\tilde p + \tfrac12 \tilde u_k \tilde u_k$, not
the $\tilde p$ of the convective form in which the budget
\eqref{eq:ns-pp-eps-rate-constant} is written. The subgrid term
$\mathrm{i} \kappa_l \hat G_l / \kappa^2$ is zero under the self-similar
constant eddy viscosity, since $\kappa_l \hat G_l = 0$. The scaling terms contribute nothing to the pressure
(\S\ref{ssec:ns-projection}).

\subsection{The integrands}\label{ssec:radial-integrands}
The integrands
\eqref{eq:ns-pp-eps-integrands} are cubic in $\tilde u_i$, and on the
solver's grid they alias badly. The cubic mask leaves the field filled up to
$|m_d| = N/3$ in every direction, so a product of $n$ such fields reaches
$|m_d| = nN/3$, and a cubic one reaches the box wavenumber $N$, twice the Nyquist
wavenumber $N/2$. The aliased content folds onto the represented modes and
corrupts both the interpolant between grid points, which the spherical mean
\eqref{eq:radial-group} uses, and every spectral derivative of a product. The
divergence terms $t_\epsilon$ and $\pi_\epsilon$ suffer most.

We therefore form every product on a refined grid of
$N_p = 2 \lceil a N / 2 \rceil$ points per direction, with padding factor
$a \ge 1$, onto which the field is carried by zero-padding its spectrum. The
normalized coefficients $\hat Q_{\mathbf{m}}/N^3$ of \eqref{eq:radial-dft}
carry over unchanged, so the field on the refined grid is the spectral
interpolant of the stored one; the Nyquist modes of the solver's grid, which
the mask keeps empty, are dropped. A product of $n$ fields is then exact while
$nN/3 < N_p/2$, that is for $a > 2n/3$: $a = 2$ makes every cubic integrand
exact, and is the default. The eddy viscosity keeps the filter width of the
run, $\Delta = H/N$, whatever $N_p$ is, and $a = 1$ recovers evaluation on the
solver's grid.

On the refined grid every derivative is taken spectrally,
$\widehat{\partial Q/\partial x_j} = \mathrm{i} \kappa_j \hat Q$, and every product
is formed pointwise in physical space. The divergence terms
$t_\epsilon$, $\pi_\epsilon$, and $d_\epsilon$ are formed as products first and
differentiated spectrally afterwards, so that each is the divergence of the
grid flux and not a sum of separately interpolated terms. Each integrand is
then reduced to its profile by \eqref{eq:radial-group} on the refined grid,
which is exact there as on the solver's grid.

\subsection{Sampling in time}\label{ssec:radial-sampling}
The profiles cost of order a hundred transforms per snapshot on the refined
grid, against the handful of a physics step on the solver's grid, so they are
not evaluated every step. They are evaluated on
the stored field $u_k$ every $n_s$ steps, $n_s = \mathrm{round}(\Delta t_s /
\Delta t)$, so that the samples are equally spaced in $t$ at the interval
$\Delta t_s$. The interval need only be short against the integral time scale
of the profiles, which the batch means of Appendix~\ref{app:stat-error}
measure; by default $\Delta t_s = 0.02$, an order of magnitude below the
integral time scale of $\epsilon(0)$ observed in the runs.

\section{Statistical error of time averages}\label{app:stat-error}

This appendix sets out how the statistical error of a time-averaged quantity
such as \eqref{eq:ns-pp-time-average} is estimated, following
\cite{shirian2023convergence}. Let $q(t)$ be a statistically stationary signal,
here the spherical mean \eqref{eq:ns-pp-sphere-mean} at some radius $r$ of
the nondimensional dissipation rate,
$q(t) = \langle \epsilon \rangle_r(t)$ of \eqref{eq:ns-pp-diss-models},
or of the resolved kinetic energy, $q(t) = \langle \tilde k \rangle_r(t)$ of
\eqref{eq:ns-kinetic-energy-nondim},
both in the self-similar frame, with mean $\langle q \rangle$, fluctuation
$q' = q - \langle q \rangle$, and root-mean-square fluctuation $q'_{\mathrm{rms}}$.
Because the clock \eqref{eq:ns-clock} is the advective time scale $L/U$, the
self-similar time $t$ is already measured in eddy turnover times.

\subsection{Batch means}\label{ssec:stat-batch-means}

Successive values of $q$ are correlated, so the samples of a run cannot be
treated as independent. We instead use the method of batch means:
\begin{enumerate}
  \item Discard the initial transient, during which the flow approaches the
  self-similar state, and let $T$ be the length of the remaining record.
  \item Split the record into $M = T/T_w$ non-overlapping windows of length
  $T_w$, and compute the window means
  \begin{equation}\label{eq:stat-window-mean}
    \langle q \rangle_i = \frac{1}{T_w} \int_{t_i}^{t_i + T_w} q(t)\, \mathrm{d}t ,
    \qquad i = 1, \dots, M .
  \end{equation}
  The samples $q_k = q(t_k)$ are recorded at equally spaced $t_k$, once per
  step (\S\ref{ssec:alg-timestep}) or once per fixed sampling interval. Each
  integral \eqref{eq:stat-window-mean} is evaluated as the exact integral of the piecewise-linear interpolant of the samples: the trapezoid
  rule on $(t_k, q_k)$, with $q$ interpolated linearly to window edges that
  fall between samples. The windows then tile $[t_0, t_0 + M T_w]$ exactly,
  and the remainder of the record shorter than $T_w$ is dropped.
  \item Estimate the mean and its standard error (SEM) as
  \begin{equation}\label{eq:stat-sem}
    q_{\mathrm{est}} = \frac{1}{M} \sum_{i=1}^{M} \langle q \rangle_i ,
    \qquad
    \mathrm{SEM} \simeq \frac{\mathrm{STD}(\langle q \rangle_i)}{\sqrt{M}} .
  \end{equation}
  \item Report the $95\%$ confidence interval
  $q_{\mathrm{est}} \pm 1.96\,\mathrm{SEM}$.
\end{enumerate}
The confidence interval assumes a Gaussian error, which the central limit
theorem supplies when the windows are uncorrelated. Equation \eqref{eq:stat-sem}
holds only if
\begin{enumerate}
  \item the windows are uncorrelated, $T_w \gg \tau_{\mathrm{corr}}$, where
  $\tau_{\mathrm{corr}}$ is the integral time scale of $q$; and
  \item there are enough windows, $M \gg 1$.
\end{enumerate}
For a fixed $T$ the two conditions pull in opposite directions, so $T_w$ has to
be chosen with care.

\subsection{Choosing the window length}\label{ssec:stat-window}

The window length is chosen by plotting $\mathrm{SEM}\sqrt{T}$ against $T_w$,
both in units of the eddy turnover time. The curve has three regimes:
\begin{itemize}
  \item \emph{Small $T_w$.} Adjacent windows are correlated, so $M$ overcounts
  the independent samples and
  $\mathrm{SEM} \approx q'_{\mathrm{rms}} \sqrt{T_w / T}$, a slope of $1/2$
  on log axes. This underestimates the true error.
  \item \emph{Plateau.} Both conditions hold and
  $\mathrm{SEM} \approx q'_{\mathrm{rms}} \sqrt{\tau_{\mathrm{corr}} / T}$,
  independent of $T_w$. This is the true SEM.
  \item \emph{Large $T_w$.} $M$ is too small and the estimate is noisy and
  unreliable.
\end{itemize}
The small-$T_w$ asymptote and the plateau intersect at
$T_w = \tau_{\mathrm{corr}}$, where
\begin{equation}\label{eq:stat-tau-corr}
  \tau_{\mathrm{corr}} = \frac{1}{q'^2_{\mathrm{rms}}}
  \int_{-\infty}^{\infty} C(\tau)\, \mathrm{d}\tau ,
  \qquad
  C(\tau) = \langle q'(t)\, q'(t + \tau) \rangle ,
\end{equation}
so the plot also yields the integral time scale. A clear plateau requires a
long record: \cite{shirian2023convergence} found one only for
$T \gtrsim 10^4$ eddy turnover times, approximately $1.6 \times 10^4$ for
homogeneous isotropic turbulence and $1.4 \times 10^4$ for channel flow at
$Re_\tau = 180$. A dissipation-type quantity, being a second-order statistic of
the velocity gradients, has $\tau_{\mathrm{corr}}$ of roughly five eddy
turnover times in homogeneous isotropic turbulence.

\subsection{Pooling an ensemble of runs}\label{ssec:stat-ensemble}

A record long enough for a clear plateau can exceed the wall-clock limit of a
single computing job. The record is then gathered as an ensemble of $n$
independent runs instead: identical in every parameter except the seed of the
random phases of the initial condition (\S\ref{ssec:ns-ic}), and in their
length $T_{\mathrm{end},j}$. Each run is a draw from the same self-similar
state, and after its own transient is discarded, run $j$ contributes a record
of length $T_j = T_{\mathrm{end},j} - t_{\mathrm{discard}}$.

The batch means of \S\ref{ssec:stat-batch-means} extend to the ensemble
without change. Every record is cut into windows of the common length $T_w$,
starting at $t_{\mathrm{discard}}$, giving $M_j = \lfloor T_j / T_w \rfloor$
window means $\langle q \rangle_{j,i}$ by \eqref{eq:stat-window-mean}, and the
estimate \eqref{eq:stat-sem} is formed over the union of all
$M = \sum_j M_j$ windows,
\begin{equation}\label{eq:stat-ensemble-mean}
  q_{\mathrm{est}} = \frac{1}{M} \sum_{j=1}^{n} \sum_{i=1}^{M_j}
  \langle q \rangle_{j,i} = \sum_{j=1}^{n} \frac{M_j}{M}\, \bar q_j ,
  \qquad
  \bar q_j = \frac{1}{M_j} \sum_{i=1}^{M_j} \langle q \rangle_{j,i} ,
\end{equation}
with $\mathrm{SEM} \simeq \mathrm{STD}(\langle q \rangle_{j,i}) / \sqrt{M}$
over the same union. Each run is therefore weighted by its length. When every
window has the same variance, which holds for $T_w$ on the plateau and runs
sampling a single stationary state, this is the minimum-variance weighting of
the run means $\bar q_j$, whose variances are inversely proportional to $M_j$.
The two conditions of \S\ref{ssec:stat-batch-means} are, if anything, better
met than within a single run: windows of different runs are uncorrelated
exactly, so only the windows adjacent within each run need $T_w \gg
\tau_{\mathrm{corr}}$, and $M$ grows with every run added. The window length
is chosen as in \S\ref{ssec:stat-window}, from the curve of
$\mathrm{SEM}\sqrt{T}$ against $T_w$ for the pooled windows, with
$T = \sum_j T_j$; the longest window scanned is the shortest $T_j$, so that
every run still contributes. The cost of pooling is one transient
$t_{\mathrm{discard}}$ discarded per run, rather than one in all.

Two runs that differ only in $T_{\mathrm{end}}$ follow the same trajectory,
up to the roundoff that eventually separates them, and would be counted twice.
For each seed, only the longest completed run is therefore used.

Pooling assumes that the runs sample the same state, and the ensemble can test
it. Within run $j$ the window means scatter about $\bar q_j$ with the
within-run variance
\begin{equation}\label{eq:stat-ensemble-within}
  s_w^2 = \frac{1}{M - n} \sum_{j=1}^{n} \sum_{i=1}^{M_j}
  \left( \langle q \rangle_{j,i} - \bar q_j \right)^2 ,
\end{equation}
and if the run means differ only by that scatter, the ratio of the
between-run to the within-run variance,
\begin{equation}\label{eq:stat-ensemble-F}
  F = \frac{1}{(n - 1)\, s_w^2} \sum_{j=1}^{n} M_j \left( \bar q_j - q_{\mathrm{est}} \right)^2 ,
\end{equation}
follows the $F$ distribution with $(n - 1, M - n)$ degrees of freedom: close
to one, with standard deviation $\approx \sqrt{2/(n - 1)}$ for $M \gg n$
(the one-way analysis of variance). A value well above one means the runs
disagree by more than their own fluctuations allow: the transient is not over
by $t_{\mathrm{discard}}$, the initial condition leaves a lasting imprint, or
$T_w$ is short of the plateau, which shrinks $s_w$. The unweighted mean of the
$\bar q_j$, with its standard error $\mathrm{STD}(\bar q_j)/\sqrt{n}$ from
the spread of the runs alone, is reported alongside \eqref{eq:stat-ensemble-mean}
as a second check that does not rely on the windows being uncorrelated.

The same pooling applies radius by radius to the spherical means of
Appendix~\ref{app:radial}, whose radii are shared by every run, and to the
coefficients of the modeled $\epsilon$ budget (\S\ref{sssec:ns-pp-eps-model}),
which are fitted to the pooled time averages and, for their confidence
intervals, to the pooled windows.

\bibliographystyle{osajnl2}
\bibliography{references}
\end{document}
